% Généralités sur la transmission du mouvement de rotation — PCT 3ème — Cours signé OnBuch+
% Programme officiel MINESEC. Compiler avec : tectonic N08-transmission-mouvement-rotation.tex
\newcommand{\DOCMATIERE}{PCT}
\newcommand{\DOCNIVEAU}{3ème}
\newcommand{\DOCMODULE}{Module 4 : Technologie et mécanique}
\newcommand{\DOCLECON}{N08}
\newcommand{\DOCTITRE}{Généralités sur la transmission du mouvement de rotation}
% OnBuch+ — préambule « cours signés » ENRICHI (compilé avec tectonic / XeLaTeX)
% Système visuel « Pop » d'OnBuch+ : fond crème (jamais blanc), bordures encre
% épaisses, ombres « dures » décalées, titres Archivo Black en capitales.
% Version MATHÉMATIQUES : graphes (pgfplots/tikz), boîtes pédagogiques
% (définition, propriété, méthode, attention, le savais-tu, exercice résolu,
% exercice, TP, bilan), page de garde et signature OnBuch+ ancrée en bas.
%
% Macros attendues dans main.tex AVANT \input{preamble} :
%   \DOCMATIERE  \DOCNIVEAU  \DOCMODULE  \DOCLECON (numéro)  \DOCTITRE
\documentclass[11pt]{article}

\usepackage{fontspec}
\usepackage[french]{babel}
\usepackage[a4paper,margin=2cm,top=3.4cm,bottom=2.8cm,headheight=1.4cm,footskip=1.3cm]{geometry}
\usepackage{amsmath,amssymb,amsfonts}
\usepackage{mathtools} % \xmapsto, \coloneqq... souvent utilisés par les modèles
\usepackage{cancel} % simplifications barrées (bilans de forces)
% \abs{z} (module, valeur absolue) et \norm{u} (norme), utilisés par les modèles
\providecommand{\abs}{}\let\abs\relax\DeclarePairedDelimiter{\abs}{\lvert}{\rvert}
\providecommand{\norm}{}\let\norm\relax\DeclarePairedDelimiter{\norm}{\lVert}{\rVert}
\usepackage[table]{xcolor} % [table] : fournit aussi \rowcolors (alternance de lignes)
\usepackage{pagecolor}
\usepackage{tikz}
\usetikzlibrary{arrows.meta,positioning,calc,shapes.geometric,shapes.misc,shapes.arrows,decorations.pathmorphing,decorations.pathreplacing,decorations.markings,patterns,shadows,mindmap,trees,fit,backgrounds,matrix,intersections,angles,quotes}
\usepackage{pgfplots}
\usepackage[version=4]{mhchem} % équations nucléaires \ce{^{235}_{92}U}
\usepackage[european,siunitx]{circuitikz} % schémas électriques
\pgfplotsset{compat=1.18}
\usepgfplotslibrary{fillbetween,groupplots} % aires entre courbes (calcul intégral)
\usepackage{enumitem}
\usepackage{fancyhdr}
% [most] charge toutes les bibliothèques tcolorbox (listings, minted, raster,
% documentation...) et gonfle inutilement la mémoire TeX sur un document long
% (~300 boîtes) : on ne charge que ce qui est réellement utilisé.
\usepackage[skins,breakable]{tcolorbox}
\usepackage{graphicx}
\usepackage{colortbl}
\usepackage{booktabs}
\usepackage{tabularx}
\usepackage{multirow}
\usepackage{diagbox} % case d'en-tête diagonale des tableaux à double entrée
% Colonnes extensibles centrée / gauche / droite (C, L, R), souvent utilisées par les modèles
\newcolumntype{C}{>{\centering\arraybackslash}X}
\newcolumntype{L}{>{\raggedright\arraybackslash}X}
\newcolumntype{R}{>{\raggedleft\arraybackslash}X}
\usepackage{array}
\usepackage{multicol}
\usepackage{siunitx}
% parse-numbers=false : accepte aussi \SI{2,5 \times 10^{-3}}{...}
\sisetup{locale=FR,output-decimal-marker={,},group-minimum-digits=4,parse-numbers=false}
\DeclareSIUnit\ohm{\ensuremath{\symup{\Omega}}} % U+2126 absent de la police
\DeclareSIUnit\division{div} % réglages d'oscilloscope (ms/div, V/div)
\DeclareSIUnit\tr{tr} % tours (vitesse de rotation, ex. \SI{1500}{\tr\per\minute})
\DeclareSIUnit\molar{M} % concentration molaire (ex. \SI{3}{\molar}, \SI{1}{\micro\molar})
\DeclareSIUnit\an{an} % année (le modèle confond parfois avec \year, un compteur TeX) — ex. \SI{5}{\kilo\meter\per\an}
\DeclareSIUnit\FCFA{FCFA} % franc CFA (ex. \SI{500}{\FCFA\per\kilo\gram})
\DeclareSIUnit\degreeBrix{\textdegree Brix} % teneur en sucre d'un sirop/jus (réfractométrie)
\DeclareSIUnit\month{mois} % le modèle utilise parfois \month, un compteur TeX, comme unité de durée
\DeclareSIUnit\microliter{\micro\liter} % le modèle écrit parfois \microliter en un seul token
\DeclareSIUnit\million{millions} % dénombrement (ex. \SI{15}{\million\per\mL} spermatozoïdes)
\usepackage{qrcode}
% \degree : commande fréquemment utilisée par les modèles pour un angle ou une
% température en dehors de \SI{}{} (non fournie par les paquets chargés).
\providecommand{\degree}{^{\circ}}
% \qed (fin de démonstration), utilisé par les modèles sans amsthm
\providecommand{\qed}{\hfill$\square$}
% \barre : double barre des valeurs interdites dans un tableau de variations (tkz-tab)
\providecommand{\barre}{\ensuremath{\|}}
% environnement proof (amsthm non chargé) utilisé par les modèles
\ifdefined\proof\else\newenvironment{proof}[1][Démonstration]{\par\noindent\textit{#1.}\ }{\qed\par}\fi
\usepackage{lastpage}
\usepackage{hyperref}
\hypersetup{hidelinks}

% ============================================================
% Palette « Pop » OnBuch+ — mêmes hex que la classe P (plus_ui.dart)
% ============================================================
\definecolor{poporange}{HTML}{F59321}
\definecolor{popdark}{HTML}{E07A0C}
\definecolor{popcream}{HTML}{FFF6E8}
\definecolor{popink}{HTML}{1C1714}
\definecolor{poppurple}{HTML}{7A5AE0}
\definecolor{popgreen}{HTML}{2E9E6A}
\definecolor{poppink}{HTML}{E0457B}
\definecolor{popblue}{HTML}{2F7DE1}
\definecolor{popgold}{HTML}{C98A12}
\definecolor{popmuted}{HTML}{8A7F76}
\definecolor{popline}{HTML}{EADFD2}
\definecolor{popgray}{HTML}{8A7F76}
\colorlet{popgrey}{popgray}
% Alias tolérants (le modèle nomme parfois les couleurs différemment)
\colorlet{popwhite}{white}
\colorlet{popblack}{popink}
\colorlet{popred}{poppink}
\colorlet{popyellow}{popgold}
% Teintes claires (fonds de tableaux/encadrés) — le modèle nomme parfois
% les couleurs avec un suffixe L (ex. popblueL) sans les avoir définies.
\colorlet{poporangeL}{poporange!20}
\colorlet{popdarkL}{popdark!20}
\colorlet{poppurpleL}{poppurple!20}
\colorlet{popgreenL}{popgreen!20}
\colorlet{poppinkL}{poppink!20}
\colorlet{popblueL}{popblue!20}
\colorlet{popgoldL}{popgold!20}
\colorlet{popredL}{popred!20}
\colorlet{popgrayL}{popgray!20}
% Teintes claires pour fonds de tableaux / zones
\colorlet{popblueL}{popblue!10!white}
\colorlet{poporangeL}{poporange!14!white}
\colorlet{popgreenL}{popgreen!12!white}
\colorlet{poppurpleL}{poppurple!12!white}
\colorlet{poppinkL}{poppink!12!white}
\colorlet{popgoldL}{popgold!14!white}

\pagecolor{popcream}

% ============================================================
% Typographie
% ============================================================
\defaultfontfeatures{Ligatures=TeX}
\setmainfont{PlusJakartaSans-Medium.ttf}[Path=../../../fonts/,BoldFont=PlusJakartaSans-Bold.ttf]
% Maths en sans-serif (Fira Math) avec chiffres et lettres droites de Plus
% Jakarta, pour que les formules se fondent dans le texte.
\usepackage[math-style=ISO,bold-style=ISO]{unicode-math}
\setmathfont{FiraMath-Regular.otf}[Path=../../../fonts/]
\setmathfont{PlusJakartaSans-Medium.ttf}[Path=../../../fonts/,range={up/num,up/{Latin,latin},\mathup}]
\usepackage{icomma} % virgule décimale sans espace en mode math
% Caractères mathématiques Unicode absents des polices de texte (Plus Jakarta,
% Archivo Black) : rendus par leur équivalent mathématique, en texte comme en
% formule — sinon ils s'affichent en carré vide (ex. « Arithmétique dans ℤ »).
\usepackage{newunicodechar}
\newunicodechar{ℝ}{\ensuremath{\mathbb{R}}}
\newunicodechar{ℤ}{\ensuremath{\mathbb{Z}}}
\newunicodechar{ℂ}{\ensuremath{\mathbb{C}}}
\newunicodechar{ℕ}{\ensuremath{\mathbb{N}}}
\newunicodechar{ℚ}{\ensuremath{\mathbb{Q}}}
\newunicodechar{𝔻}{\ensuremath{\mathbb{D}}}
\newunicodechar{α}{\ensuremath{\alpha}}
\newunicodechar{β}{\ensuremath{\beta}}
\newunicodechar{γ}{\ensuremath{\gamma}}
\newunicodechar{δ}{\ensuremath{\delta}}
\newunicodechar{ε}{\ensuremath{\varepsilon}}
\newunicodechar{θ}{\ensuremath{\theta}}
\newunicodechar{λ}{\ensuremath{\lambda}}
\newunicodechar{μ}{\ensuremath{\mu}}
\newunicodechar{σ}{\ensuremath{\sigma}}
\newunicodechar{τ}{\ensuremath{\tau}}
\newunicodechar{φ}{\ensuremath{\varphi}}
\newunicodechar{ψ}{\ensuremath{\psi}}
\newunicodechar{ω}{\ensuremath{\omega}}
\newunicodechar{Σ}{\ensuremath{\Sigma}}
% majuscules produites par \MakeUppercase dans les titres (α-aminés -> Α)
\newunicodechar{Α}{\ensuremath{\alpha}}
\newunicodechar{Β}{\ensuremath{\beta}}
\newunicodechar{Γ}{\ensuremath{\Gamma}}
\newunicodechar{Δ}{\ensuremath{\Delta}}
\newunicodechar{Ε}{\ensuremath{\varepsilon}}
\newunicodechar{Θ}{\ensuremath{\Theta}}
\newunicodechar{Λ}{\ensuremath{\Lambda}}
\newunicodechar{Μ}{\ensuremath{\mu}}
\newunicodechar{Τ}{\ensuremath{\tau}}
\newunicodechar{Φ}{\ensuremath{\Phi}}
\newunicodechar{Ψ}{\ensuremath{\Psi}}
\newunicodechar{Ω}{\ensuremath{\Omega}}
\newunicodechar{↦}{\ensuremath{\mapsto}}
\newunicodechar{∈}{\ensuremath{\in}}
\newunicodechar{∉}{\ensuremath{\notin}}
\newunicodechar{∧}{\ensuremath{\wedge}}
\newunicodechar{⇔}{\ensuremath{\Leftrightarrow}}
\newunicodechar{⟺}{\ensuremath{\Longleftrightarrow}}
\newunicodechar{⇒}{\ensuremath{\Rightarrow}}
\newunicodechar{⟹}{\ensuremath{\Longrightarrow}}
\newunicodechar{∘}{\ensuremath{\circ}}
\newunicodechar{∪}{\ensuremath{\cup}}
\newunicodechar{∩}{\ensuremath{\cap}}
\newunicodechar{≡}{\ensuremath{\equiv}}
\newunicodechar{⊂}{\ensuremath{\subset}}
\newunicodechar{⊥}{\ensuremath{\perp}}
\newunicodechar{∃}{\ensuremath{\exists}}
\newunicodechar{∀}{\ensuremath{\forall}}
\newunicodechar{∛}{\ensuremath{\sqrt[3]{\,}}}
\newunicodechar{∜}{\ensuremath{\sqrt[4]{\,}}}
\newunicodechar{⃗}{\ensuremath{{}^{\to}}}
\newunicodechar{ⁿ}{\ifmmode^{n}\else\textsuperscript{\ensuremath{n}}\fi}
\newunicodechar{ˣ}{\ifmmode^{x}\else\textsuperscript{\ensuremath{x}}\fi}
\newunicodechar{ᵇ}{\ifmmode^{b}\else\textsuperscript{\ensuremath{b}}\fi}
\newunicodechar{ʳ}{\ifmmode^{r}\else\textsuperscript{\ensuremath{r}}\fi}
\newunicodechar{ᵘ}{\ifmmode^{u}\else\textsuperscript{\ensuremath{u}}\fi}
\newunicodechar{ʸ}{\ifmmode^{y}\else\textsuperscript{\ensuremath{y}}\fi}
\newunicodechar{ᵃ}{\ifmmode^{a}\else\textsuperscript{\ensuremath{a}}\fi}
\newunicodechar{ᵉ}{\ifmmode^{e}\else\textsuperscript{\ensuremath{e}}\fi}
\newunicodechar{ᵏ}{\ifmmode^{k}\else\textsuperscript{\ensuremath{k}}\fi}
\newunicodechar{ᵖ}{\ifmmode^{p}\else\textsuperscript{\ensuremath{p}}\fi}
\newunicodechar{ᴺ}{\ifmmode^{N}\else\textsuperscript{\ensuremath{N}}\fi}
\newunicodechar{ᶜ}{\ifmmode^{c}\else\textsuperscript{\ensuremath{c}}\fi}
\newunicodechar{ʰ}{\ifmmode^{h}\else\textsuperscript{\ensuremath{h}}\fi}
\newunicodechar{ᵠ}{\ifmmode^{q}\else\textsuperscript{\ensuremath{q}}\fi}
\newunicodechar{ₙ}{\ifmmode_{n}\else\textsubscript{\ensuremath{n}}\fi}
\newunicodechar{ₐ}{\ifmmode_{a}\else\textsubscript{\ensuremath{a}}\fi}
\newunicodechar{ᵢ}{\ifmmode_{i}\else\textsubscript{\ensuremath{i}}\fi}
\newunicodechar{ᵣ}{\ifmmode_{r}\else\textsubscript{\ensuremath{r}}\fi}
\newunicodechar{ₖ}{\ifmmode_{k}\else\textsubscript{\ensuremath{k}}\fi}
\newunicodechar{ᵦ}{\ifmmode_{\beta}\else\textsubscript{\ensuremath{\beta}}\fi}
\newunicodechar{ₓ}{\ifmmode_{x}\else\textsubscript{\ensuremath{x}}\fi}
\newfontfamily\popdisplay{ArchivoBlack-Regular.ttf}[Path=../../../fonts/]
\newfontfamily\poplogo{SpaceGrotesk-Bold.ttf}[Path=../../../fonts/]
\newfontfamily\popsemi{PlusJakartaSans-SemiBold.ttf}[Path=../../../fonts/]
\newfontfamily\popxbold{PlusJakartaSans-ExtraBold.ttf}[Path=../../../fonts/]
\newfontfamily\popxboldls{PlusJakartaSans-ExtraBold.ttf}[Path=../../../fonts/,LetterSpace=3.0]

\setlist[enumerate]{leftmargin=1.7em,itemsep=5pt,topsep=4pt}
\setlist[itemize]{leftmargin=1.7em,itemsep=4pt,topsep=4pt,label={\color{poporange}\textbullet}}
\setlength{\parindent}{0pt}
\setlength{\parskip}{5pt}
\renewcommand{\arraystretch}{1.25}

% pgfplots : style Pop par défaut
\pgfplotsset{
  every axis/.append style={axis line style={popink,line width=1pt},tick style={popink},
    grid style={popline,line width=0.5pt},label style={font=\small},tick label style={font=\footnotesize}},
  popcurve/.style={poporange,line width=1.8pt,no markers,smooth},
}
% Même style disponible dans un \draw TikZ simple (hors pgfplots)
\tikzset{popcurve/.style={poporange,line width=1.8pt}}
\tikzset{popgrid/.style={popline,very thin}}
\colorlet{popcurve}{poporange} % le nom du style est parfois utilisé comme couleur

% ============================================================
% Pill / squiggle
% ============================================================
\newcommand{\poppill}[3]{%
  \tikz[baseline=-0.55ex]\node[draw=popink,line width=1.2pt,rounded corners=2.6pt,%
    fill=#2,inner xsep=6.5pt,inner ysep=3pt]{%
    \popxboldls\fontsize{7}{7}\selectfont\color{#3}\MakeUppercase{#1}};%
}
\newcommand{\popsquiggle}[1]{%
  \tikz[baseline=-0.4ex]{
    \draw[#1,line width=2.6pt,line cap=round]
      plot [smooth] coordinates {(0,0.09) (0.34,0.01) (0.68,0.21) (1.02,0.01) (1.36,0.10)};
  }%
}
% Mot-clé surligné (marqueur orange sous le texte)
\newcommand{\cle}[1]{{\popxbold\color{popdark}#1}}

% ============================================================
% En-tête / pied
% ============================================================
\pagestyle{fancy}
\fancyhf{}
\renewcommand{\headrulewidth}{0pt}
\renewcommand{\footrulewidth}{0pt}
\lhead{\raisebox{-0.15ex}{\poplogo\fontsize{13}{13}\selectfont\color{popink}OnBuch\color{poporange}+}}
\rhead{\poppill{\DOCMATIERE\ \textbullet\ \DOCNIVEAU\ \textbullet\ Leçon \DOCLECON}{white}{popink}}
\lfoot{\popxbold\fontsize{7.6}{7.6}\selectfont\color{popmuted}\MakeUppercase{Cours signé OnBuch+}\ \textbullet\ {\popsemi\DOCTITRE}}
\rfoot{\tikz[baseline=-0.6ex]\node[rounded corners=8.5pt,fill=popink,minimum height=17pt,minimum width=17pt,inner xsep=5pt,inner sep=0]{\popxbold\fontsize{8}{8}\selectfont\color{popcream}\thepage\,/\,\pageref*{LastPage}};}

% ============================================================
% Titres
% ============================================================
\newcommand{\chaptitre}[2]{%
  \begin{tcolorbox}[enhanced,colback=poporange,colframe=popink,boxrule=2.6pt,arc=11pt,
      drop shadow=popink,left=15pt,right=15pt,top=12pt,bottom=11pt,before skip=3pt,after skip=18pt]
    \poppill{#2}{popink}{popcream}\par\vspace{9pt}
    {\raggedright\hyphenpenalty=10000\exhyphenpenalty=10000\popdisplay\fontsize{22}{24}\selectfont\color{popcream}\MakeUppercase{#1}\par}
  \end{tcolorbox}
}
% Section numérotée : pastille orange avec le numéro + titre Archivo
\newcounter{popsec}
\newcommand{\coursec}[1]{%
  \stepcounter{popsec}\setcounter{popsub}{0}%
  \par\vspace{16pt}\noindent
  \begin{minipage}{\linewidth}
  \tikz[baseline=-0.7ex]\node[draw=popink,line width=1.6pt,fill=poporange,rounded corners=4pt,minimum size=22pt,inner sep=0]{\popdisplay\fontsize{12}{12}\selectfont\color{popcream}\thepopsec};%
  \hspace{8pt}{\popdisplay\fontsize{15}{17}\selectfont\color{popink}\MakeUppercase{#1}}\par
  \vspace{4pt}\noindent{\color{poporange}\rule{36pt}{3.6pt}}
  \end{minipage}\par\nopagebreak\vspace{8pt}
}
\newcounter{popsub}
\newcommand{\courssub}[1]{%
  \stepcounter{popsub}%
  \par\vspace{9pt}\noindent{\popxbold\fontsize{11.5}{13}\selectfont\color{popink}{\color{poporange}\thepopsec.\thepopsub}\hspace{6pt}#1}\par\nopagebreak\vspace{4pt}
}

% ============================================================
% Boîtes pédagogiques — même recette : fond blanc, bordure épaisse colorée,
% ombre dure encre, badge plein. Titre optionnel [#1].
% ============================================================
\tcbset{popbase/.style={breakable,enhanced,colback=white,boxrule=2.3pt,arc=9pt,
  shadow={3pt}{-3pt}{0pt}{popink},left=14pt,right=14pt,top=11pt,bottom=11pt,before skip=11pt,after skip=15pt,
  fontupper=\popsemi\fontsize{10.6}{14.5}\selectfont\color{popink},
  fontlower=\popsemi\fontsize{10.6}{14.5}\selectfont\color{popink}}}
% Coche dessinée (le glyphe ✓ n'existe pas dans les polices du document)
\newcommand{\popcheck}[1]{\tikz[baseline=-0.5ex]\draw[#1,line width=1.6pt,line cap=round,line join=round](-0.13,0.0)--(-0.04,-0.09)--(0.13,0.1);}
\providecommand{\coche}{\popcheck{popgreen}}
\providecommand{\resultat}[1]{\par\smallskip\noindent\fcolorbox{popgreen}{popgreenL}{\parbox{\dimexpr\linewidth-2\fboxsep-2\fboxrule\relax}{\textbf{Résultat.} #1}}\par\smallskip}
\newcommand{\popboxhead}[3]{\poppill{#1}{#2}{white}\ifx\relax#3\relax\else\hspace{6pt}{\popxbold\fontsize{10.6}{12}\selectfont\color{popink}#3}\fi\par\vspace{7pt}}

\newtcolorbox{aretenir}[1][]{popbase,colframe=popgreen,before upper={\popboxhead{À retenir}{popgreen}{#1}}}
\newtcolorbox{exemplebox}[1][]{popbase,colframe=poporange,before upper={\popboxhead{Exemple}{poporange}{#1}}}
\newtcolorbox{definition}[1][]{popbase,colframe=popblue,before upper={\popboxhead{Définition}{popblue}{#1}}}
\newtcolorbox{propriete}[1][]{popbase,colframe=poppurple,before upper={\popboxhead{Propriété}{poppurple}{#1}}}
\newtcolorbox{methode}[1][]{popbase,colframe=popgold,before upper={\popboxhead{Méthode}{popgold}{#1}}}
\newtcolorbox{attention}[1][]{popbase,colframe=poppink,before upper={\popboxhead{Attention}{poppink}{#1}}}
\newtcolorbox{experience}[1][]{popbase,colframe=popblue,colback=popblueL,before upper={\popboxhead{Expérience}{popblue}{#1}}}
% « Le savais-tu ? » : carte violette pleine (contexte, culture, Cameroun)
\newtcolorbox{savaistu}[1][]{popbase,colback=poppurple,colframe=popink,
  fontupper=\popsemi\fontsize{10.4}{14.2}\selectfont\color{white},
  before upper={\poppill{Le savais-tu\,?}{white}{poppurple}\ifx\relax#1\relax\else\hspace{6pt}{\popxbold\color{white}#1}\fi\par\vspace{7pt}}}
% Exercice résolu : énoncé en haut, \tcblower puis la solution sur fond crème
\newcounter{popexo}\newcounter{popexr}
\newtcolorbox{exoresolu}[1][]{popbase,colframe=poporange,colbacklower=popcream,
  segmentation style={popink,line width=1.2pt,dashed},
  before upper={\stepcounter{popexr}\popboxhead{Exercice résolu \thepopexr}{poporange}{#1}},
  before lower={\poppill{Solution}{popink}{popcream}\par\vspace{6pt}}}
% Exercice (non résolu en place) — la correction est dans la partie Corrigés
\newtcolorbox{exercice}[1][]{popbase,colframe=popink,
  before upper={\stepcounter{popexo}\popboxhead{Exercice \thepopexo}{popink}{#1}}}
% Corrigé
\newtcolorbox{corrige}[1][]{popbase,colframe=popgreen,colback=popgreenL,
  before upper={\popboxhead{Corrigé}{popgreen}{#1}}}
% Objectifs / prérequis / bilan
\newtcolorbox{objectifs}{popbase,colframe=popink,colback=poporangeL,
  before upper={\popboxhead{Objectifs de la leçon}{popink}{}}}
\newtcolorbox{prerequis}{popbase,colframe=popink,colback=popblueL,
  before upper={\popboxhead{Prérequis}{popblue}{}}}
\newtcolorbox{bilan}{popbase,colframe=popink,colback=white,boxrule=2.8pt,
  before upper={\popboxhead{Fiche bilan}{poporange}{L'essentiel à connaître pour l'examen}}}
% Figure encadrée avec légende
\newtcolorbox{popfigure}[1][]{enhanced,breakable=false,colback=white,colframe=popline,boxrule=1.4pt,arc=8pt,
  left=10pt,right=10pt,top=10pt,bottom=8pt,before skip=10pt,after skip=12pt,halign=center,
  lower separated=false,halign lower=center,
  fontlower=\popsemi\fontsize{9}{11.5}\selectfont\color{popmuted},
  #1}
\newcommand{\legende}[1]{\par\vspace{6pt}{\popsemi\fontsize{9}{11.5}\selectfont\color{popmuted}#1}}

% ============================================================
% Signature OnBuch+
% ============================================================
\newcommand{\poplogoinline}[1]{{\poplogo\fontsize{#1}{#1}\selectfont\color{popink}OnBuch\color{poporange}+}}
\newcommand{\signatureonbuch}{%
  \begin{tcolorbox}[enhanced,colback=popink,colframe=popink,boxrule=2.2pt,arc=10pt,
      drop shadow=poporange,left=14pt,right=14pt,top=10pt,bottom=10pt,before skip=0pt,after skip=0pt]
    \tikz[baseline=-0.7ex]\node[circle,fill=popgreen,draw=popcream,line width=1.3pt,minimum size=22pt,inner sep=0]{\popcheck{white}};%
    \hspace{9pt}\begin{minipage}[c]{0.62\linewidth}
      {\popxbold\fontsize{12}{13}\selectfont\color{popcream}Signé\ \textbullet\ {\poplogo OnBuch\color{poporange}+}}\par\vspace{2pt}
      {\raggedright\popsemi\fontsize{8}{10.5}\selectfont\color{popline}\MakeUppercase{Équipe pédagogique OnBuch+}\\\MakeUppercase{Conforme au programme officiel MINESEC}\par}
    \end{minipage}\hfill
    {\poplogo\fontsize{22}{22}\selectfont\color{popcream}OnBuch\color{poporange}+}
  \end{tcolorbox}%
}

% ============================================================
% Page de garde
% #1 titre · #2 matière · #3 niveau · #4 module · #5 n° leçon · #6 durée ·
% #7 description / objectifs (texte ou itemize)
% ============================================================
\newcommand{\pagegarde}[7]{%
  \thispagestyle{empty}
  \newgeometry{a4paper,margin=1.7cm,top=1.6cm,bottom=1.4cm}%
  \begin{tikzpicture}[remember picture,overlay]
    \fill[poporange!18] ([xshift=-3.2cm,yshift=-3.6cm]current page.north east) circle (4.6cm);
    \fill[poppurple!14] ([xshift=2.4cm,yshift=7.6cm]current page.south west) circle (3.8cm);
  \end{tikzpicture}%
  {\poplogo\fontsize{30}{30}\selectfont\color{popink}OnBuch\color{poporange}+}\hfill
  \raisebox{0.6ex}{\poppill{Cours signé \textbullet\ Premium}{popink}{popcream}}\par
  \vspace{1pt}\popsquiggle{poppurple}\par
  \vspace{8pt}
  \begin{tcolorbox}[enhanced,colback=poporange,colframe=popink,boxrule=2.8pt,arc=13pt,
      drop shadow=popink,left=18pt,right=18pt,top=12pt,bottom=13pt,after skip=10pt]
    \poppill{#2 \textbullet\ #3}{popink}{popcream}\hspace{5pt}\poppill{#4}{white}{popink}\par\vspace{8pt}
    {\popdisplay\fontsize{14}{14}\selectfont\color{popink}LEÇON #5}\par\vspace{4pt}
    {\raggedright\hyphenpenalty=10000\exhyphenpenalty=10000\popdisplay\fontsize{25}{27}\selectfont\color{popcream}\MakeUppercase{#1}\par}
    \vspace{8pt}
    {\popxbold\fontsize{9.5}{9.5}\selectfont\color{popink}\MakeUppercase{Durée conseillée : #6}}
  \end{tcolorbox}
  \begin{tcolorbox}[enhanced,colback=white,colframe=popink,boxrule=2.2pt,arc=10pt,
      drop shadow=popink,left=16pt,right=16pt,top=10pt,bottom=10pt,after skip=8pt]
    \poppill{Dans cette leçon}{poppurple}{white}\par\vspace{5pt}
    \popsemi\fontsize{9.3}{12.6}\selectfont\color{popink}#7
  \end{tcolorbox}
  \noindent\begin{minipage}[t]{0.487\linewidth}
    \begin{tcolorbox}[enhanced,colback=popgreen,colframe=popink,boxrule=2.2pt,arc=10pt,
        drop shadow=popink,left=11pt,right=11pt,top=10pt,bottom=10pt,height=3.1cm,valign=center]
      \tcbox[colback=white,colframe=popink,boxrule=1.3pt,arc=5pt,boxsep=0pt,nobeforeafter,
        left=3pt,right=3pt,top=3pt,bottom=3pt,valign=center]{\qrcode[height=1.9cm]{https://play.google.com/store/apps/details?id=com.luvvix.onbuch&hl=fr}}%
      \hspace{8pt}\begin{minipage}[c]{0.5\linewidth}
        {\popdisplay\fontsize{11}{12.5}\selectfont\color{white}\MakeUppercase{Télécharge OnBuch}}\par\vspace{4pt}
        {\raggedright\popsemi\fontsize{8.4}{11}\selectfont\color{white}Gratuit \textbullet\ Google Play\\Tuteur IA, annales, concours, résultats\par}
      \end{minipage}
    \end{tcolorbox}
  \end{minipage}\hfill
  \begin{minipage}[t]{0.487\linewidth}
    \begin{tcolorbox}[enhanced,colback=poppurple,colframe=popink,boxrule=2.2pt,arc=10pt,
        drop shadow=popink,left=11pt,right=11pt,top=10pt,bottom=10pt,height=3.1cm,valign=center]
      \tikz[baseline=-0.7ex]\node[circle,fill=white,draw=popink,line width=1.3pt,minimum size=1.6cm,inner sep=0]{\popdisplay\fontsize{24}{24}\selectfont\color{poppurple}+};%
      \hspace{8pt}\begin{minipage}[c]{0.6\linewidth}
        {\popdisplay\fontsize{11}{12.5}\selectfont\color{white}\MakeUppercase{Passe à OnBuch+}}\par\vspace{4pt}
        {\raggedright\popsemi\fontsize{8.4}{11}\selectfont\color{white}Tous les cours signés, Tuteur IA illimité, fiches PDF — onglet OnBuch+ de l'app\par}
      \end{minipage}
    \end{tcolorbox}
  \end{minipage}
  \vspace{8pt}
  \popcontenu
  \vfill
  \signatureonbuch
  \restoregeometry
  \newpage
}

% Rangée « Ce cours contient » (page de garde)
\newcommand{\poptile}[3]{% #1 couleur #2 gros texte #3 légende
  \begin{tcolorbox}[enhanced,colback=#1,colframe=popink,boxrule=2pt,arc=9pt,shadow={2.5pt}{-2.5pt}{0pt}{popink},
      left=6pt,right=6pt,top=9pt,bottom=9pt,halign=center,valign=center,height=1.75cm,width=\linewidth,nobeforeafter]
    {\popdisplay\fontsize{12.5}{13}\selectfont\color{white}\MakeUppercase{#2}}\par\vspace{3pt}
    {\centering\popsemi\fontsize{7.8}{9.5}\selectfont\color{white}#3\par}
  \end{tcolorbox}}
\newcommand{\popcontenu}{%
  \par\noindent{\popxboldls\fontsize{7.5}{7.5}\selectfont\color{popmuted}CE COURS SIGNÉ CONTIENT}\par\vspace{7pt}
  \noindent\begin{minipage}[t]{0.235\linewidth}\poptile{popblue}{Cours}{complet, illustré, pas à pas}\end{minipage}\hfill
  \begin{minipage}[t]{0.235\linewidth}\poptile{poporange}{Méthodes}{et exercices résolus}\end{minipage}\hfill
  \begin{minipage}[t]{0.235\linewidth}\poptile{poppink}{Exercices}{type examen + corrigés}\end{minipage}\hfill
  \begin{minipage}[t]{0.235\linewidth}\poptile{popgreen}{Fiche bilan}{l'essentiel à retenir}\end{minipage}\par}

% Sommaire
\newtcolorbox{sommaire}{enhanced,colback=white,colframe=popink,boxrule=2.2pt,arc=10pt,
  drop shadow=popink,left=18pt,right=18pt,top=13pt,bottom=13pt,before skip=6pt,after skip=20pt,
  before upper={\poppill{Sommaire}{poppurple}{white}\par\vspace{9pt}\popsemi\fontsize{10.6}{14.5}\selectfont\color{popink}}}

% Signature de fin de cours — poussée tout en bas de la dernière page
\newcommand{\findecours}{%
  \par\vfill\nopagebreak
  \begin{center}\popsquiggle{poporange}\end{center}\vspace{4pt}
  \signatureonbuch
}
\AtBeginDocument{\def\llbracket{\mathopen{[\mkern-3.5mu[}}\def\rrbracket{\mathclose{]\mkern-3.5mu]}}} % intervalles d'entiers (glyphe ⟦ absent de la police)
% Glyphes absents de Fira Math : redéfinis à partir de glyphes disponibles
\AtBeginDocument{%
  \def\setminus{\mathbin{\backslash}}\let\smallsetminus\setminus
  \def\vdots{\mathord{\vbox{\baselineskip3.2pt\lineskiplimit0pt\kern4pt\hbox{.}\hbox{.}\hbox{.}}}}%
  \def\ddots{\mathinner{\mkern1mu\raise6.5pt\hbox{.}\mkern2mu\raise3.6pt\hbox{.}\mkern2mu\raise0.7pt\hbox{.}\mkern1mu}}%
  \def\triangle{\mathord{\scalebox{0.9}{$\symup{\Delta}$}}}%
  \def\nabla{\mathord{\raisebox{\depth}{\scalebox{1}[-1]{$\symup{\Delta}$}}}}%
  \def\checkmark{\popcheck{popdark}}%
  \def\star{\mathbin{\ast}}\def\bigstar{\mathord{\ast}}%
  \def\diamond{\mathbin{\scalebox{0.6}{$\Diamond$}}}%
  \def\bigcup{\mathop{\mathchoice{\raisebox{-0.3em}{\scalebox{1.7}{$\cup$}}}{\scalebox{1.25}{$\cup$}}{\cup}{\cup}}\displaylimits}%
  \def\bigcap{\mathop{\mathchoice{\raisebox{-0.3em}{\scalebox{1.7}{$\cap$}}}{\scalebox{1.25}{$\cap$}}{\cap}{\cap}}\displaylimits}%
  \def\lesseqqgtr{\mathrel{\lessgtr}}\def\bigoplus{\mathop{\oplus}\displaylimits}%
}
\providecommand{\ssi}{si et seulement si} % abréviation fréquente
\AtBeginDocument{\providecommand{\wideparen}{\overparen}} % arc de cercle

\begin{document}

\pagegarde{Généralités sur la transmission du mouvement de rotation}{PCT}{3ème}{Module 4 : Technologie et mécanique}{N08}{2 séances (1 h chacune)}{Cette leçon établit les bases théoriques et pratiques de la transmission du mouvement de rotation : grandeurs caractéristiques (vitesse angulaire, rapport de transmission), étude des systèmes courants (engrenages, courroies, chaînes, vis sans fin) et règles de sécurité. Elle prépare l'élève à analyser et dimensionner des chaînes cinématiques simples.
\vspace{4pt}
\begin{multicols}{2}\small
\begin{itemize}
\item À la fin de cette leçon, je sais définir la rotation, la vitesse angulaire, la période et la fréquence, et convertir des tr/min en rad/s.
\item À la fin de cette leçon, je sais définir le rapport de transmission et le calculer pour un engrenage (dents) ou une courroie (diamètres).
\item À la fin de cette leçon, je sais identifier le sens de rotation en sortie par rapport à l'entrée selon le type de liaison (engrenage externe, courroie ouverte/croisée, vis sans fin).
\item À la fin de cette leçon, je sais lister les avantages et inconvénients des transmissions par engrenages, courroies, chaînes et vis sans fin.
\item À la fin de cette leçon, je sais expliquer le phénomène de glissement (courroies) et l'irréversibilité de la vis sans fin.
\item À la fin de cette leçon, je sais calculer la vitesse de sortie et le couple de sortie connaissant le rendement d'un système.
\end{itemize}\end{multicols}}

\begin{sommaire}
\begin{multicols}{2}
\begin{enumerate}
\item 1. Rotation et grandeurs caractéristiques du mouvement
\item 2. Rapport de transmission : définition, calcul et signe
\item 3. Les transmissions par contact direct : les engrenages
\item 4. Les transmissions par liaison flexible : courroies et chaînes
\item 5. Systèmes particuliers, rendement et sécurité
\item Activité d'intégration
\item Méthodes et exercices résolus
\item Exercices
\item Corrigés des exercices
\item Fiche bilan
\end{enumerate}
\end{multicols}
\end{sommaire}

\begin{objectifs}
\begin{itemize}
\item À la fin de cette leçon, je sais définir la rotation, la vitesse angulaire, la période et la fréquence, et convertir des tr/min en rad/s.
\item À la fin de cette leçon, je sais définir le rapport de transmission et le calculer pour un engrenage (dents) ou une courroie (diamètres).
\item À la fin de cette leçon, je sais identifier le sens de rotation en sortie par rapport à l'entrée selon le type de liaison (engrenage externe, courroie ouverte/croisée, vis sans fin).
\item À la fin de cette leçon, je sais lister les avantages et inconvénients des transmissions par engrenages, courroies, chaînes et vis sans fin.
\item À la fin de cette leçon, je sais expliquer le phénomène de glissement (courroies) et l'irréversibilité de la vis sans fin.
\item À la fin de cette leçon, je sais calculer la vitesse de sortie et le couple de sortie connaissant le rendement d'un système.
\item À la fin de cette leçon, je sais identifier les zones de pincement et les dispositifs de protection (carters, capots) sur un schéma réel.
\item À la fin de cette leçon, je sais choisir le système de transmission adapté à une situation simple (ex: vélo, moulin, portillon).
\end{itemize}
\end{objectifs}

\begin{prerequis}
\begin{itemize}
\item Notion de cercle : rayon, diamètre, périmètre (classe de 6ème/5ème).
\item Vitesse linéaire, unités (m/s, km/h) et conversions (4ème/3ème début année).
\item Notion de force, de couple (moment de force) et d'énergie (travail, puissance) (3ème Mécanique).
\item Représentation vectorielle simple (direction, sens, norme) (3ème Physique).
\item Lecture de schémas techniques normalisés (vues de face, de profil) (Technologie 4ème/5ème).
\end{itemize}
\end{prerequis}

\newpage

\coursec{Rotation et grandeurs caractéristiques du mouvement}

Au quartier, le \emph{benskin} (moto-taxi) grimpe la côte de Nkolbisson sans caler grâce à sa chaîne de transmission, pendant que le moulin à maïs du marché réduit bruyamment les grains en farine grâce à un système d'engrenages robustes. Sur le terrain de l'école, le vélo du surveillant général change de braquet pour affronter le dénivelé de la montée. Comment ces machines transforment-elles la rotation rapide du moteur ou des jambes en une rotation lente mais puissante à la roue ? Pourquoi le sens de rotation change-t-il parfois ?

Pour répondre à ces questions, il faut d'abord apprendre à \cle{décrire} et à \cle{mesurer} un mouvement de rotation. C'est l'objectif de cette première partie : définir la rotation, introduire ses grandeurs caractéristiques (angle, vitesse angulaire, fréquence, période) et établir la relation fondamentale $v = \omega \times r$ qui relie la vitesse d'un point du solide à sa distance à l'axe.

\courssub{Le mouvement de rotation : axe, sens et angle}

\textbf{Observation de la vie courante.} Prends une porte de classe, une hélice de ventilateur, la roue du vélo du surveillant : tous ces objets tournent. Mais tournent-ils n'importe comment ? Non. La porte pivote autour de ses gonds, le ventilateur autour de son axe central, la roue autour de son moyeu. Dans chaque cas, il existe une \cle{droite fixe} autour de laquelle tout l'objet tourne.

\begin{definition}[Mouvement de rotation]
Un solide est animé d'un \cle{mouvement de rotation} autour d'une droite fixe $(\Delta)$, appelée \cle{axe de rotation}, si tout point $M$ du solide décrit un \cle{cercle} dont :
\begin{itemize}
\item le centre appartient à l'axe $(\Delta)$ ;
\item le plan est perpendiculaire à l'axe $(\Delta)$ ;
\item le rayon est la distance $r$ du point $M$ à l'axe.
\end{itemize}
Tous les points du solide balayent le \emph{même angle} pendant la même durée : on dit que le solide tourne \emph{d'un seul bloc}.
\end{definition}

\begin{attention}[Axe et points de l'axe]
Les points situés \emph{sur} l'axe de rotation ne se déplacent pas : leur vitesse est nulle. C'est pourquoi le centre du moyeu de la roue de vélo reste immobile par rapport au cadre, alors que le pneu file très vite. Ne dis jamais qu'un point de l'axe « tourne sur place » : il est \textbf{fixe}.
\end{attention}

\textbf{Le sens de rotation.} Une rotation possède un sens, que l'on repère face au disque qui tourne :
\begin{itemize}
\item le \cle{sens horaire} : celui des aiguilles d'une montre ;
\item le \cle{sens anti-horaire} (ou sens trigonométrique) : le sens contraire.
\end{itemize}
Ce sens dépend du côté d'où l'on regarde ! Vu de face, le ventilateur tourne dans un sens ; vu de derrière, il semble tourner dans l'autre. C'est pourquoi on précisera toujours le point d'observation, ou mieux, on utilisera le vecteur vitesse angulaire (voir 1.4).

\textbf{L'angle balayé.} Pour mesurer « de combien » le solide a tourné, on utilise l'\cle{angle} $\theta$ (lettre grecque \emph{thêta}) balayé par un rayon. L'unité légale est le \cle{radian} (symbole : rad).

\begin{savaistu}[Qu'est-ce qu'un radian ?]
Un \cle{radian} est l'angle au centre qui intercepte, sur le cercle, un arc dont la longueur est égale au rayon. Comme le périmètre du cercle vaut $2\pi r$, un tour complet correspond à $2\pi$ radians, soit environ $6{,}28$ rad. Un radian vaut environ $57{,}3$ degrés. Retiens les correspondances utiles : $360^{\circ} = 2\pi$ rad ; $180^{\circ} = \pi$ rad ; $90^{\circ} = \dfrac{\pi}{2}$ rad.
\end{savaistu}

\begin{popfigure}
\begin{tikzpicture}[scale=1.1]
% Disque
\fill[popblueL] (0,0) circle (2);
\draw[popdark,thick] (0,0) circle (2);
% Axe de rotation
\draw[popdark,very thick] (0,-2.6) -- (0,2.9);
\fill[popdark] (0,0) circle (0.07);
\node[popdark,right] at (0,2.9) {axe $(\Delta)$};
% Vecteur omega
\draw[-{Stealth[length=3mm]},poporange,very thick] (0,0.4) -- (0,2.4);
\node[poporange,left] at (-0.08,1.6) {$\vec{\omega}$};
% Point M et rayon
\fill[poppurple] (1.2,1.6) circle (0.09);
\node[poppurple,above right] at (1.2,1.6) {$M$};
\draw[poppurple,thick] (0,0) -- (1.2,1.6) node[midway,below left] {$r$};
% Trajectoire circulaire de M
\draw[poppurple,dashed,thick] (0,0) circle (2);
% Vecteur vitesse tangent
\draw[-{Stealth[length=3mm]},popgreen,very thick] (1.2,1.6) -- (0.4,2.4);
\node[popgreen,above left] at (0.55,2.35) {$\vec{v}$};
% Sens de rotation
\draw[-{Stealth[length=2.5mm]},poporange,thick] (2.3,-0.6) arc (-15:45:1.1);
\node[poporange] at (3.15,0.35) {sens de};
\node[poporange] at (3.15,0.0) {rotation};
\end{tikzpicture}
\legende{Disque en rotation autour de l'axe $(\Delta)$. Le point $M$, situé à la distance $r$ de l'axe, décrit un cercle (pointillés violets). Sa vitesse $\vec{v}$ (verte) est \textbf{tangente} à ce cercle. Le vecteur $\vec{\omega}$ (orange) est porté par l'axe.}
\end{popfigure}

\courssub{La vitesse angulaire $\omega$}

\textbf{Intuition.} Deux ventilateurs tournent : l'un fait 5 tours par seconde, l'autre 20 tours par seconde. Le second tourne « plus vite ». Mais comment le dire avec rigueur ? En mesurant l'angle balayé \emph{par unité de temps}. C'est exactement le rôle de la vitesse angulaire.

\begin{experience}[Mesurer la rapidité d'une rotation]
\textbf{Matériel :} un ventilateur de table dont une pale est marquée à la craie, un chronomètre (celui d'un téléphone).
\begin{enumerate}
\item On met le ventilateur en marche et on attend que la rotation soit régulière.
\item On compte le nombre $n$ de tours effectués par la pale marquée pendant une durée $\Delta t = \SI{10}{s}$.
\item On calcule le nombre de tours par seconde : $f = \dfrac{n}{\Delta t}$.
\end{enumerate}
\textbf{Observations :} pour un ventilateur domestique en vitesse moyenne, on compte environ $n = 150$ tours en \SI{10}{s}, soit $f = 15$ tours par seconde.
\textbf{Interprétation :} chaque tour correspond à un angle de $2\pi$ rad. L'angle balayé par seconde est donc $\omega = 2\pi \times f = 2\pi \times 15 \approx \SI{94}{rad/s}$. Cette grandeur $\omega$ caractérise la rapidité de la rotation, indépendamment du point du solide que l'on observe.
\end{experience}

\begin{definition}[Vitesse angulaire]
La \cle{vitesse angulaire} $\omega$ (lettre grecque \emph{oméga}) d'un solide en rotation mesure la rapidité de variation de l'angle de rotation :
\[ \omega = \dfrac{\theta}{t} \]
où $\theta$ est l'angle balayé (en rad) pendant la durée $t$ (en s). L'unité de $\omega$ est le \cle{radian par seconde} (rad/s).
\end{definition}

\begin{propriete}[Lien entre $\omega$, la fréquence $f$ et la période $T$]
Pour une rotation régulière :
\begin{itemize}
\item la \cle{fréquence} $f$ est le nombre de tours par seconde ; son unité est le \cle{hertz} (Hz) : $f = \dfrac{\text{nombre de tours}}{t}$ ;
\item la \cle{période} $T$ est la durée d'un tour complet, en secondes : $T = \dfrac{1}{f}$.
\end{itemize}
Comme un tour correspond à un angle de $2\pi$ rad, on a :
\[ \omega = 2\pi f = \dfrac{2\pi}{T} \]
\end{propriete}

\begin{attention}[Le radian est une unité « sans dimension »]
Le radian est un rapport de deux longueurs (arc sur rayon) : il est \textbf{sans dimension}. Ainsi $\omega$ en rad/s est homogène à des $\text{s}^{-1}$, exactement comme la fréquence $f$ en Hz. Mais attention au piège classique du BEPC : \textbf{ne confonds pas rad/s et tr/s} ! Une fréquence de $f = 5$ tr/s (soit \SI{5}{Hz}) correspond à $\omega = 2\pi \times 5 \approx \SI{31,4}{rad/s}$, et non à \SI{5}{rad/s}. Le facteur $2\pi$ est indispensable.
\end{attention}

\courssub{Conversion tr/min $\leftrightarrow$ rad/s}

Dans l'industrie et la vie courante, les vitesses de rotation sont presque toujours données en \cle{tours par minute} (tr/min, noté aussi rpm en anglais) : moteur de moto à \SI{3000}{tr/min}, meuleuse à \SI{11000}{tr/min}, groupe électrogène à \SI{1500}{tr/min}. Or nos formules exigent des rad/s. Il faut donc savoir convertir.

\textbf{Établissons le facteur de conversion.} Une vitesse de $N$ tours par minute signifie :
\begin{itemize}
\item $N$ tours en \SI{60}{s}, soit une fréquence $f = \dfrac{N}{60}$ tours par seconde ;
\item chaque tour vaut $2\pi$ rad, donc $\omega = 2\pi f = \dfrac{2\pi N}{60} = \dfrac{\pi N}{30}$.
\end{itemize}

\begin{methode}[Convertir $N$ (tr/min) en $\omega$ (rad/s)]
\begin{enumerate}
\item Écris la valeur $N$ en tr/min.
\item Multiplie par $\dfrac{\pi}{30} \approx 0{,}105$ :
\[ \omega \;(\text{rad/s}) = N \;(\text{tr/min}) \times \dfrac{\pi}{30} \]
\item Pour la conversion inverse, multiplie par $\dfrac{30}{\pi} \approx 9{,}55$.
\end{enumerate}
\textbf{Exemple de référence à retenir :} $N = \SI{3000}{tr/min}$ donne $\omega = 3000 \times \dfrac{\pi}{30} = 100\pi \approx \SI{314}{rad/s}$.
\end{methode}

\begin{exemplebox}[Quelques conversions utiles]
\begin{center}
\begin{tabularx}{\linewidth}{|l|X|X|}\hline
\rowcolor{popblueL} \textbf{Machine} & \textbf{$N$ (tr/min)} & \textbf{$\omega$ (rad/s)} \\ \hline
Roue de voiture à \SI{90}{km/h} & $\approx 750$ & $\approx 78{,}5$ \\ \hline
Moteur de moto-taxi (régime moyen) & $3000$ & $\approx 314$ \\ \hline
Meuleuse d'atelier & $11000$ & $\approx 1152$ \\ \hline
\end{tabularx}
\end{center}
\end{exemplebox}

\courssub{Le vecteur vitesse angulaire et la règle de la main droite}

Comment indiquer à la fois l'axe, le sens et la rapidité d'une rotation en un seul objet mathématique ? Les physiciens ont inventé le \cle{vecteur vitesse angulaire} $\vec{\omega}$.

\begin{definition}[Vecteur vitesse angulaire]
Le vecteur vitesse angulaire $\vec{\omega}$ est un vecteur :
\begin{itemize}
\item dont la \cle{direction} est celle de l'axe de rotation $(\Delta)$ ;
\item dont le \cle{sens} est donné par la règle de la main droite ;
\item dont la \cle{norme} est la vitesse angulaire $\omega$ (en rad/s).
\end{itemize}
\end{definition}

\textbf{La règle de la main droite} (dite aussi règle de la \emph{vis de charpentier} ou du \emph{tire-bouchon}) : ferme le poing droit et laisse dépasser le pouce. Si les doigts repliés indiquent le sens de rotation du solide, alors le pouce indique le sens de $\vec{\omega}$. Autre image : une vis ordinaire avance dans le sens de $\vec{\omega}$ quand on la visse dans le sens de rotation.

\begin{popfigure}
\begin{tikzpicture}[scale=1.0]
% Poing stylisé
\fill[popgoldL] (0,0) circle (1.05);
\draw[popdark,thick] (0,0) circle (1.05);
% Doigts repliés : arcs
\foreach \a/\r in {150/0.62, 110/0.68, 70/0.68, 30/0.62}{
  \draw[popdark,thick,fill=popgold] (\a:\r) circle (0.26);
}
% Flèche de rotation autour du poing
\draw[-{Stealth[length=3mm]},popgreen,very thick] (1.5,-0.9) arc (-35:60:1.55);
\node[popgreen,align=center] at (2.6,0.35) {doigts : sens de};
\node[popgreen,align=center] at (2.6,-0.05) {la rotation};
% Pouce = vecteur omega
\draw[-{Stealth[length=3.5mm]},poporange,line width=2pt] (0,0.2) -- (0,2.6);
\node[poporange,right] at (0.05,2.5) {pouce $= \vec{\omega}$};
% Axe en pointillés
\draw[popdark,dashed,thick] (0,-1.6) -- (0,0.2);
\node[popdark] at (0,-1.9) {axe $(\Delta)$};
\end{tikzpicture}
\legende{Règle de la main droite : les doigts repliés suivent le sens de rotation (flèche verte), le pouce levé donne le sens du vecteur $\vec{\omega}$ (flèche orange), porté par l'axe de rotation.}
\end{popfigure}

\begin{exemplebox}[Sens de $\vec{\omega}$ pour une roue de vélo]
Le vélo avance vers la droite : vu depuis la gauche du vélo, la roue tourne dans le sens horaire. La main droite, doigts enroulés dans ce sens, pointe le pouce vers la gauche du vélo : $\vec{\omega}$ est dirigé le long de l'axe de la roue, vers la gauche. Si le vélo faisait demi-tour, $\vec{\omega}$ changerait de sens.
\end{exemplebox}

\courssub{Vitesse linéaire d'un point du solide : $v = \omega \times r$}

\textbf{Intuition.} Sur un manège de fête foraine, les enfants assis près du bord extérieur « filent » beaucoup plus vite que ceux assis près du centre, alors que le manège tourne d'un seul bloc : $\omega$ est le même pour tous. La différence vient de la \cle{distance à l'axe}. Plus on est loin de l'axe, plus le cercle parcouru est grand, et plus la vitesse est élevée.

\textbf{Établissement de la relation.} Considérons un point $M$ à la distance $r$ de l'axe. Pendant une période $T$ (un tour complet), le point $M$ parcourt le périmètre de son cercle, soit une distance $d = 2\pi r$. Sa vitesse est donc :
\begin{align*}
v &= \dfrac{d}{T} = \dfrac{2\pi r}{T} \\
v &= \dfrac{2\pi}{T} \times r \\
v &= \omega \times r
\end{align*}

\begin{propriete}[Relation fondamentale $v = \omega \times r$]
Pour tout point $M$ lié à un solide en rotation à la vitesse angulaire $\omega$, situé à la distance $r$ de l'axe :
\[ v = \omega \times r \]
avec $v$ en m/s, $\omega$ en rad/s et $r$ en \textbf{mètres}. Le vecteur vitesse $\vec{v}$ est \textbf{tangent} au cercle décrit par $M$, orienté dans le sens de la rotation.
\end{propriete}

\begin{attention}[Toujours convertir $r$ en mètres]
L'erreur la plus fréquente au BEPC : utiliser $r$ en centimètres ou en millimètres directement dans la formule. Si la meule a un rayon de \SI{10}{cm}, écris $r = \SI{0,10}{m}$ avant de calculer. Autre piège : confondre le \textbf{diamètre} et le \textbf{rayon}. Si l'énoncé donne le diamètre $d$, divise-le par deux : $r = \dfrac{d}{2}$.
\end{attention}

\begin{exoresolu}[Vitesse d'un point de la roue de vélo]
\textbf{Énoncé.} La roue du vélo du surveillant général a un rayon $r = \SI{0,35}{m}$. Elle tourne à $f = 5$ tours par seconde.
\begin{enumerate}
\item Calcule la vitesse angulaire $\omega$ de la roue.
\item Calcule la vitesse linéaire $v$ d'un point du pneu (à la périphérie). Convertis-la en km/h.
\end{enumerate}
\tcblower
\textbf{Solution détaillée.}
\begin{enumerate}
\item On applique $\omega = 2\pi f$ :
\[ \omega = 2\pi \times 5 = 10\pi \approx \SI{31,4}{rad/s}. \]
\item On applique $v = \omega \times r$ avec $r$ en mètres :
\[ v = 31{,}4 \times 0{,}35 \approx \SI{11}{m/s}. \]
Conversion : $v = 11 \times 3{,}6 \approx \SI{40}{km/h}$.
\end{enumerate}
\textbf{Conclusion.} La vitesse d'un point du pneu par rapport à l'axe est d'environ \SI{40}{km/h} : c'est cohérent avec la vitesse d'un vélo lancé dans une descente.
\end{exoresolu}

\begin{exoresolu}[Vitesse au bout d'une pale d'éolienne et d'une meuleuse]
\textbf{Énoncé.}
\begin{enumerate}
\item Une éolienne possède des pales de longueur $r = \SI{20}{m}$ et tourne à $N = \SI{15}{tr/min}$. Calcule la vitesse $v$ de l'extrémité d'une pale.
\item Une meuleuse d'atelier porte un disque de diamètre $d = \SI{115}{mm}$ et tourne à $N' = \SI{11000}{tr/min}$. Calcule la vitesse $v'$ d'un point du bord du disque.
\end{enumerate}
\tcblower
\textbf{Solution détaillée.}
\begin{enumerate}
\item Conversion : $\omega = N \times \dfrac{\pi}{30} = 15 \times \dfrac{\pi}{30} = \dfrac{\pi}{2} \approx \SI{1,57}{rad/s}$. \\
Puis : $v = \omega \times r = 1{,}57 \times 20 \approx \SI{31,4}{m/s}$, soit environ \SI{113}{km/h}. Le bout d'une pale d'éolienne file à la vitesse d'une voiture sur autoroute !
\item Rayon : $r' = \dfrac{d}{2} = \dfrac{115}{2} = \SI{57,5}{mm} = \SI{0,0575}{m}$. \\
Conversion : $\omega' = 11000 \times \dfrac{\pi}{30} \approx \SI{1152}{rad/s}$. \\
Puis : $v' = \omega' \times r' = 1152 \times 0{,}0575 \approx \SI{66}{m/s}$, soit environ \SI{238}{km/h}. C'est pourquoi un disque de meuleuse qui se brise est extrêmement dangereux : les éclats sont projetés à très grande vitesse. Porte toujours des lunettes de protection !
\end{enumerate}
\end{exoresolu}

\courssub{Démarrage d'une machine : de la rampe au régime permanent}

Quand tu mets en marche le moulin à maïs ou le groupe électrogène du quartier, la vitesse angulaire n'atteint pas instantanément sa valeur maximale. Elle croît progressivement pendant quelques secondes (c'est la \cle{phase de démarrage}), puis se stabilise : on dit que la machine atteint son \cle{régime permanent}. La courbe $\omega(t)$ ci-dessous illustre ces deux phases.

\begin{popfigure}
\begin{tikzpicture}
\begin{axis}[
  width=0.85\linewidth, height=6cm,
  xlabel={temps $t$ (s)}, ylabel={$\omega$ (rad/s)},
  xmin=0, xmax=10, ymin=0, ymax=350,
  xtick={0,2,4,6,8,10}, ytick={0,100,200,314},
  yticklabels={0,100,200,314},
  grid=both, grid style={popline!40},
  axis line style={popdark},
  label style={popdark},
]
% Rampe de demarrage : 0 a 4 s, omega monte lineairement jusqu'a 314
\addplot[popcurve,domain=0:4,samples=50]{78.5*x};
% Regime permanent : palier a 314 rad/s
\addplot[popcurve,domain=4:10,samples=10]{314};
% Annotations
\node[popblue,align=center] at (axis cs:2,120) {phase de démarrage};
\node[popblue,align=center] at (axis cs:2,80) {(rampe)};
\node[poporange,align=center] at (axis cs:7.2,270) {régime permanent};
\node[poporange,align=center] at (axis cs:7.2,230) {(palier : $\omega \approx \SI{314}{rad/s}$)};
\draw[popdark,dashed] (axis cs:4,0) -- (axis cs:4,314);
\end{axis}
\end{tikzpicture}
\legende{Évolution de la vitesse angulaire $\omega$ d'un moteur de moto-taxi au démarrage : rampe de montée en régime pendant environ \SI{4}{s}, puis palier au régime permanent ($N = \SI{3000}{tr/min}$, soit $\omega \approx \SI{314}{rad/s}$).}
\end{popfigure}

\begin{aretenir}[L'essentiel de la section]
\begin{itemize}
\item Un solide en \cle{rotation} tourne autour d'un \cle{axe fixe} ; chaque point décrit un cercle de rayon $r$ (distance à l'axe).
\item La \cle{vitesse angulaire} $\omega = \dfrac{\theta}{t}$ (en rad/s) mesure la rapidité de la rotation ; elle est la même pour tous les points du solide.
\item Relations fondamentales : $\omega = 2\pi f = \dfrac{2\pi}{T}$ et $\omega\,(\text{rad/s}) = N\,(\text{tr/min}) \times \dfrac{\pi}{30}$.
\item La \cle{vitesse linéaire} d'un point dépend de sa position : $v = \omega \times r$, avec $r$ en mètres ; $\vec{v}$ est tangente au cercle.
\item Le \cle{vecteur} $\vec{\omega}$ est porté par l'axe ; son sens est donné par la \cle{règle de la main droite}.
\end{itemize}
\end{aretenir}

\begin{exercice}[À toi de jouer]
\begin{enumerate}
\item Le tambour d'une machine à laver effectue un essorage à $N = \SI{1200}{tr/min}$. Calcule sa vitesse angulaire $\omega$ en rad/s, puis sa fréquence $f$ et sa période $T$.
\item Le tambour a un rayon $r = \SI{25}{cm}$. Calcule la vitesse linéaire d'un point du linge collé contre la paroi du tambour. Convertis le résultat en km/h.
\item Une hélice de ventilateur tourne dans le sens horaire vue de face. Dans quel sens est dirigé son vecteur $\vec{\omega}$ : vers toi ou vers le mur derrière le ventilateur ? Justifie avec la règle de la main droite.
\end{enumerate}
\end{exercice}

\begin{corrige}[Exercice : À toi de jouer]
\begin{enumerate}
\item $\omega = N \times \dfrac{\pi}{30} = 1200 \times \dfrac{\pi}{30} = 40\pi \approx \SI{125,7}{rad/s}$. Fréquence : $f = \dfrac{N}{60} = \SI{20}{Hz}$. Période : $T = \dfrac{1}{f} = \SI{0,05}{s}$.
\item Conversion du rayon : $r = \SI{25}{cm} = \SI{0,25}{m}$. Alors $v = \omega \times r = 125{,}7 \times 0{,}25 \approx \SI{31,4}{m/s}$, soit $31{,}4 \times 3{,}6 \approx \SI{113}{km/h}$.
\item Doigts de la main droite enroulés dans le sens horaire vu de face : le pouce pointe \emph{vers le mur}, c'est-à-dire à l'opposé de l'observateur. Le vecteur $\vec{\omega}$ est donc dirigé de l'observateur vers le ventilateur.
\end{enumerate}
\end{corrige}

Nous savons désormais décrire et mesurer une rotation. Reste la grande question de l'accroche : comment la rotation du moteur est-elle \emph{transmise} à la roue, et comment sa vitesse est-elle modifiée ? C'est le rôle du \cle{rapport de transmission}, que nous étudierons dans la section suivante.

\coursec{Rapport de transmission : définition, calcul et signe}

Dans la section précédente, tu as appris à caractériser un mouvement de rotation : la vitesse angulaire $\omega$ (en rad/s) et la fréquence de rotation $N$ (en tr/min). Mais dans une machine réelle — le moteur de la moto-taxi, la perceuse du menuisier de quartier, le moulin à maïs — le mouvement ne naît pas là où on en a besoin : il est produit par un moteur, puis \emph{transmis} jusqu'à l'organe qui travaille (la roue, le foret, la meule). Et presque toujours, la vitesse change en chemin. Comment comparer la vitesse d'entrée et la vitesse de sortie ? C'est exactement le rôle du \cle{rapport de transmission}, la notion centrale de cette leçon.

\courssub{Définition du rapport de transmission}

\textbf{Situation-problème.} Le moteur d'un groupe électrogène tourne à \SI{3000}{tr/min}, mais l'alternateur qui produit le courant n'a besoin que de \SI{1500}{tr/min}. Entre les deux, le mécanicien a interposé deux poulies reliées par une courroie. De combien la vitesse a-t-elle été divisée ? Réponse : par $2$. Ce nombre, $1/2$, est le rapport de transmission du système.

\begin{definition}[Rapport de transmission]
On appelle \cle{rapport de transmission} d'un système, noté $k$, le quotient de la vitesse de rotation de l'arbre de \textbf{sortie} par celle de l'arbre d'\textbf{entrée} :
\[ k = \frac{\omega_{\text{sortie}}}{\omega_{\text{entrée}}} = \frac{N_{\text{sortie}}}{N_{\text{entrée}}} \]
$k$ est un nombre \textbf{sans unité} (on divise une vitesse par une vitesse). L'entrée est l'arbre menant (celui du moteur), la sortie est l'arbre mené (celui qui entraîne la charge).
\end{definition}

\begin{attention}[Ne pas inverser entrée et sortie]
C'est l'erreur n°1 au BEPC : le rapport de transmission est toujours $k = \dfrac{\text{sortie}}{\text{entrée}}$, jamais l'inverse. Avant tout calcul, identifie clairement quel arbre est \textbf{menant} (entrée, côté moteur) et quel arbre est \textbf{mené} (sortie, côté charge). Un rapport $k = 3$ ou $k = 1/3$ ne racontent pas du tout la même histoire !
\end{attention}

\begin{attention}[Le sens du rapport : réducteur ou multiplicateur]
\begin{itemize}
\item Si $k < 1$ : la sortie tourne \textbf{moins vite} que l'entrée. Le système est un \cle{réducteur}. La vitesse diminue, mais le \textbf{couple} (l'« effort de rotation ») est multiplié : c'est ce qui permet à un petit moteur de soulever de lourdes charges (treuil, grue).
\item Si $k > 1$ : la sortie tourne \textbf{plus vite} que l'entrée. Le système est un \cle{multiplicateur}. La vitesse augmente, mais le couple diminue (éolienne, perceuse à grande vitesse).
\item Si $k = 1$ : la vitesse est transmise telle quelle (simple renvoi du mouvement).
\end{itemize}
Retenons : \textbf{on ne gagne jamais sur les deux tableaux} — ce qu'on gagne en vitesse, on le perd en force, et réciproquement.
\end{attention}

\courssub{Cas des engrenages (roues dentées)}

\textbf{Observation.} Prends deux roues dentées qui engrènent : une petite de 20 dents (menante) et une grande de 40 dents (menée). Quand la petite fait un tour complet, elle pousse 20 dents de la grande : la grande n'a donc fait qu'un \textbf{demi-tour}. La vitesse a été divisée par 2. De plus, en regardant les flèches, on constate que les deux roues tournent en \textbf{sens opposés}.

\begin{propriete}[Rapport de transmission d'un engrenage externe]
Pour deux roues dentées en prise directe (engrenage externe), le rapport de transmission est :
\[ k = \frac{N_{\text{sortie}}}{N_{\text{entrée}}} = \frac{Z_{\text{entrée}}}{Z_{\text{sortie}}} \]
où $Z$ désigne le \textbf{nombre de dents} de chaque roue. Le sens de rotation est \textbf{inversé} à chaque engrenage externe. Pour que deux roues puissent engrêner, elles doivent avoir le même \cle{module} $m = \dfrac{D}{Z}$ (rapport du diamètre primitif au nombre de dents) : c'est la « taille » des dents.
\end{propriete}

\begin{experience}[Vérification du rapport d'un engrenage]
\textbf{Matériel :} deux roues dentées de Lego ou de jouet mécanique ($Z_1 = 20$ et $Z_2 = 40$), un feutre, une feuille blanche.\par
\textbf{Protocole :} (1) Marque d'un trait de feutre une dent de chaque roue au point de contact. (2) Fais tourner lentement la petite roue d'exactement un tour. (3) Observe la position du repère de la grande roue.\par
\textbf{Observations :} la grande roue n'a effectué qu'un demi-tour, et elle tourne en sens inverse de la petite.\par
\textbf{Interprétation :} $N_{\text{sortie}} = N_{\text{entrée}} \times \dfrac{Z_1}{Z_2} = 1 \times \dfrac{20}{40} = 0{,}5$ tour. L'expérience confirme la formule $k = \dfrac{Z_{\text{entrée}}}{Z_{\text{sortie}}}$.
\end{experience}

\begin{popfigure}
\begin{tikzpicture}[scale=0.9]
% Roue 1 (petite) Z=20
\draw[thick, popblue, fill=popblueL] (0,0) circle (1);
\draw[thick, popblue] (0,0) circle (1.18);
\foreach \a in {0,18,...,342} {
  \draw[popblue, thick] (\a:1.12) -- (\a:1.28);
}
\draw[fill=popdark] (0,0) circle (0.08);
\node[popdark, font=\small\bfseries] at (0,-1.6) {Roue 1 : $Z_1 = 20$};
\node[popdark, font=\small] at (0,-2.0) {(menante)};
% Roue 2 (grande) Z=40
\draw[thick, poporange, fill=poporangeL] (3.2,0) circle (2);
\draw[thick, poporange] (3.2,0) circle (2.18);
\foreach \a in {0,9,...,351} {
  \draw[poporange, thick] (\a:2.12) -- (\a:2.28);
}
\draw[fill=popdark] (3.2,0) circle (0.08);
\node[popdark, font=\small\bfseries] at (3.2,-2.6) {Roue 2 : $Z_2 = 40$};
\node[popdark, font=\small] at (3.2,-3.0) {(menée)};
% Flèches de sens de rotation opposés (Engrenage externe = inversion)
% Roue 1 (gauche) : Flèche en haut vers la DROITE (Trigo 60->120 = CCW -> Tangente en haut = Gauche. On veut Droite -> CW -> 120->60)
% Roue 2 (droite) : Flèche en haut vers la GAUCHE (CW -> Tangente en haut = Gauche. 120->60)
\draw[-{Stealth[length=3mm]}, very thick, popgreen] (0.7,1.35) arc (120:60:1.55); % CW -> Fleche haut -> Droite
\node[popgreen, font=\small\bfseries] at (0,1.85) {Sens 1};
\draw[-{Stealth[length=3mm]}, very thick, poppink] (2.5,2.35) arc (60:120:1.55); % CCW -> Fleche haut -> Gauche
\node[poppink, font=\small\bfseries] at (3.2,2.85) {Sens 2 (inverse)};
\end{tikzpicture}
\legende{Deux roues dentées en prise : $Z_1 = 20$ (menante) et $Z_2 = 40$ (menée). Les sens de rotation sont opposés (flèches anti-parallèles en haut). Ici $k = \dfrac{Z_1}{Z_2} = \dfrac{20}{40} = \dfrac{1}{2}$ : c'est un réducteur.}
\end{popfigure}

\begin{exemplebox}[Boîte de vitesses de moto-taxi]
Le pignon de sortie moteur d'une moto a $Z_1 = 15$ dents et entraîne une couronne de $Z_2 = 45$ dents. Le rapport est $k = \dfrac{15}{45} = \dfrac{1}{3}$ : la roue tourne trois fois moins vite que le pignon moteur, mais avec trois fois plus de couple — indispensable pour grimper les côtes de Yaoundé en charge !
\end{exemplebox}

\courssub{Cas des poulies-courroies}

\textbf{Observation.} Sur le groupe électrogène, la courroie relie une petite poulie côté moteur à une grande poulie côté alternateur. Quand la petite fait un tour, une longueur de courroie égale à sa circonférence $\pi D_1$ défile ; cette même longueur fait tourner la grande poulie de $\dfrac{\pi D_1}{\pi D_2} = \dfrac{D_1}{D_2}$ tour. D'où la formule.

\begin{propriete}[Rapport de transmission d'une poulie-courroie]
Pour deux poulies reliées par une courroie, le rapport de transmission est :
\[ k = \frac{N_{\text{sortie}}}{N_{\text{entrée}}} = \frac{D_{\text{entrée}}}{D_{\text{sortie}}} \]
où $D$ est le \textbf{diamètre primitif} de chaque poulie. Le sens de rotation dépend du montage (convention de représentation schématique) :
\begin{itemize}
\item \textbf{Courroie ouverte} (non croisée) : les flèches de rotation sont \textbf{parallèles} (même sens schématique) ;
\item \textbf{Courroie croisée} : les flèches de rotation sont \textbf{anti-parallèles} (sens inversés schématique).
\end{itemize}
Contrairement aux engrenages, il peut exister un léger \textbf{glissement} entre la courroie et les poulies : le rapport réel est alors un peu différent du rapport théorique.
\end{propriete}

\begin{popfigure}
\begin{tikzpicture}[scale=0.85]
% --- Courroie ouverte : MÊME SENS (Flèches parallèles) ---
% Convention : Flèches en haut vers la DROITE pour les deux.
% Poulie 1 (Gauche) : CW (arc 120->60). Poulie 2 (Droite) : CCW (arc 60->120).
\begin{scope}
\draw[thick, popblue, fill=popblueL] (0,0) circle (0.7);
\draw[thick, poporange, fill=poporangeL] (3.5,0) circle (1.3);
\draw[fill=popdark] (0,0) circle (0.06);
\draw[fill=popdark] (3.5,0) circle (0.06);
% Courroie ouverte
\draw[very thick, popdark] (0,0.7) -- (3.5,1.3);
\draw[very thick, popdark] (0,-0.7) -- (3.5,-1.3);
% Flèches : Parallèles (toutes deux vers la droite en haut)
\draw[-{Stealth[length=2.5mm]}, very thick, popgreen] (0.45,0.95) arc (120:60:1.0); % CW -> Haut -> Droite
\draw[-{Stealth[length=2.5mm]}, very thick, popgreen] (3.05,1.55) arc (60:120:1.0);  % CCW -> Haut -> Droite
\node[popdark, font=\small\bfseries] at (1.75,-2.0) {Courroie ouverte : même sens};
\node[popdark, font=\small] at (0,-1.15) {$D_1$};
\node[popdark, font=\small] at (3.5,-1.75) {$D_2$};
\end{scope}
% --- Courroie croisée : SENS INVERSÉS (Flèches anti-parallèles) ---
% Convention : Flèche Gauche Haut -> Droite. Flèche Droite Haut -> Gauche.
% Poulie 1 (Gauche) : CW (120->60). Poulie 2 (Droite) : CW (120->60).
\begin{scope}[xshift=8.5cm]
\draw[thick, popblue, fill=popblueL] (0,0) circle (0.7);
\draw[thick, poporange, fill=poporangeL] (3.5,0) circle (1.3);
\draw[fill=popdark] (0,0) circle (0.06);
\draw[fill=popdark] (3.5,0) circle (0.06);
% Courroie croisée
\draw[very thick, popdark] (0,0.7) -- (3.5,-1.3);
\draw[very thick, popdark] (0,-0.7) -- (3.5,1.3);
% Flèches : Anti-parallèles (Gauche Haut->Droite, Droite Haut->Gauche)
\draw[-{Stealth[length=2.5mm]}, very thick, popgreen] (0.45,0.95) arc (120:60:1.0); % CW -> Haut -> Droite
\draw[-{Stealth[length=2.5mm]}, very thick, poppink] (3.95,1.55) arc (120:60:1.0);  % CW -> Haut -> Gauche
\node[popdark, font=\small\bfseries] at (1.75,-2.0) {Courroie croisée : sens inversés};
\node[popdark, font=\small] at (0,-1.15) {$D_1$};
\node[popdark, font=\small] at (3.5,-1.75) {$D_2$};
\end{scope}
\end{tikzpicture}
\legende{Transmission par poulies et courroie ($D_1 < D_2$, donc $k = \dfrac{D_1}{D_2} < 1$ : réducteur). À gauche, courroie ouverte : flèches parallèles (même sens schématique). À droite, courroie croisée : flèches anti-parallèles (sens inversés schématique).}
\end{popfigure}

\begin{attention}[Glissement de la courroie]
La transmission par courroie se fait par \textbf{adhérence} : si la courroie est détendue ou si la charge est trop brutale, elle \emph{patine}. C'est un défaut pour la précision (le rapport n'est plus exact), mais parfois un avantage : en cas de blocage, la courroie glisse au lieu de casser le moteur. C'est une sécurité mécanique naturelle.
\end{attention}

\courssub{Cas des chaînes et pignons}

Sur un vélo ou une moto, la transmission se fait par une \textbf{chaîne} qui relie deux pignons dentés. La chaîne combine les avantages des deux systèmes précédents : comme la courroie, elle transmet \emph{à distance} ; comme l'engrenage, ses dents s'engrènent dans les maillons, donc il n'y a \textbf{aucun glissement}.

\begin{propriete}[Rapport de transmission d'une chaîne]
Pour une transmission par chaîne et pignons :
\[ k = \frac{N_{\text{sortie}}}{N_{\text{entrée}}} = \frac{Z_{\text{pignon d'entrée}}}{Z_{\text{pignon de sortie}}} \]
Les deux pignons tournent dans le \textbf{même sens} (la chaîne joue le rôle d'une courroie ouverte : flèches parallèles), et le rapport est \textbf{exact} car il n'y a pas de glissement.
\end{propriete}

\begin{exemplebox}[Le vélo du quartier]
Sur un vélo, le plateau (entrée, côté pédales) a $Z_1 = 44$ dents et le pignon de la roue arrière a $Z_2 = 22$ dents. Alors $k = \dfrac{44}{22} = 2 > 1$ : c'est un \textbf{multiplicateur} — la roue fait deux tours pour un tour de pédalier. Le cycliste pédale lentement mais la roue tourne vite ; en contrepartie, il doit fournir un grand effort dans les montées.
\end{exemplebox}

\courssub{Enchaînement de plusieurs étages : rapport global}

Les machines réelles associent souvent plusieurs transmissions en cascade : c'est le cas d'une boîte de vitesses ou d'un réducteur de moulin. Comment obtenir le rapport global ?

\begin{methode}[Calcul du rapport de transmission global]
Pour un train de transmissions à $n$ étages :
\begin{enumerate}
\item Identifier l'entrée et la sortie de \textbf{chaque étage}.
\item Calculer le rapport de chaque étage : $k_1$, $k_2$, \dots, $k_n$ (avec la bonne formule selon le type : $Z_{\text{entrée}}/Z_{\text{sortie}}$ ou $D_{\text{entrée}}/D_{\text{sortie}}$).
\item \textbf{Multiplier} les rapports : $k_{\text{global}} = k_1 \times k_2 \times \dots \times k_n$.
\item Déterminer le \textbf{sens final} : compter le nombre d'inversions de sens (engrenages externes et courroies croisées). S'il est \textbf{pair}, la sortie tourne dans le même sens que l'entrée ; s'il est \textbf{impair}, le sens est inversé.
\item Conclure : si $k_{\text{global}} < 1$, le système est un réducteur ; si $k_{\text{global}} > 1$, un multiplicateur.
\end{enumerate}
\end{methode}

\begin{popfigure}
\begin{tikzpicture}[scale=0.8, font=\small]
% Étage 1 : engrenage (Inversion)
\draw[thick, popblue, fill=popblueL] (0,0) circle (0.55);
\draw[thick, poporange, fill=poporangeL] (1.75,0) circle (1.1);
\draw[fill=popdark] (0,0) circle (0.05);
\draw[fill=popdark] (1.75,0) circle (0.05);
\node[popdark] at (0,-0.95) {$Z_1$};
\node[popdark] at (1.75,-1.5) {$Z_2$};
\node[popblue, font=\small\bfseries] at (0.9,1.5) {Étage 1 : engrenage};
\node[popblue] at (0.9,1.1) {$k_1 = \dfrac{Z_1}{Z_2}$};
% Flèches Engrenage 1 : Inversion (Anti-parallèles)
% G1 Gauche : Haut -> Droite (CCW 60:120). G2 Droite : Haut -> Gauche (CW 120:60).
\draw[-{Stealth[length=2mm]}, thick, popgreen] (0.3,0.8) arc (60:120:0.85);
\draw[-{Stealth[length=2mm]}, thick, poppink] (1.45,1.35) arc (120:60:0.85);
% Liaison
\draw[very thick, popdark, -{Stealth[length=2.5mm]}] (2.9,0) -- (3.7,0);
% Étage 2 : engrenage (Inversion)
\draw[thick, popblue, fill=popblueL] (4.4,0) circle (0.55);
\draw[thick, poporange, fill=poporangeL] (6.15,0) circle (1.1);
\draw[fill=popdark] (4.4,0) circle (0.05);
\draw[fill=popdark] (6.15,0) circle (0.05);
\node[popdark] at (4.4,-0.95) {$Z_3$};
\node[popdark] at (6.15,-1.5) {$Z_4$};
\node[popgreen, font=\small\bfseries] at (5.3,1.5) {Étage 2 : engrenage};
\node[popgreen] at (5.3,1.1) {$k_2 = \dfrac{Z_3}{Z_4}$};
% Flèches Engrenage 2 : Inversion (Anti-parallèles)
% G3 Gauche (arbre commun G2) : Même sens que G2 -> Haut -> Gauche (CW 120:60).
% G4 Droite : Inversion -> Haut -> Droite (CCW 60:120).
\draw[-{Stealth[length=2mm]}, thick, popgreen] (4.1,0.8) arc (120:60:0.85);
\draw[-{Stealth[length=2mm]}, thick, poppink] (5.85,1.35) arc (60:120:0.85);
% Liaison
\draw[very thick, popdark, -{Stealth[length=2.5mm]}] (7.3,0) -- (8.1,0);
% Étage 3 : poulies Courroie OUVERTE (Même sens = Flèches Parallèles)
\draw[thick, popblue, fill=popblueL] (8.8,0) circle (0.55);
\draw[thick, poporange, fill=poporangeL] (10.55,0) circle (1.1);
\draw[fill=popdark] (8.8,0) circle (0.05);
\draw[fill=popdark] (10.55,0) circle (0.05);
\draw[very thick, popdark] (8.8,0.55) -- (10.55,1.1);
\draw[very thick, popdark] (8.8,-0.55) -- (10.55,-1.1);
\node[popdark] at (8.8,-0.95) {$D_5$};
\node[popdark] at (10.55,-1.5) {$D_6$};
\node[poppurple, font=\small\bfseries] at (9.7,1.5) {Étage 3 : courroie};
\node[poppurple] at (9.7,1.1) {$k_3 = \dfrac{D_5}{D_6}$};
% Flèches Courroie Ouverte : Parallèles (Haut -> Droite pour les deux)
% P5 Gauche (arbre commun G4) : Même sens que G4 -> Haut -> Droite (CCW 60:120).
% P6 Droite : Même sens -> Haut -> Droite (CCW 60:120).
\draw[-{Stealth[length=2mm]}, thick, popgreen] (9.1,0.8) arc (60:120:0.85);
\draw[-{Stealth[length=2mm]}, thick, popgreen] (10.25,1.35) arc (60:120:0.85);
% Entrée / Sortie
\draw[-{Stealth[length=3mm]}, very thick, poppink] (-1.6,0) -- (-0.7,0);
\node[poppink, font=\small\bfseries, align=center] at (-1.6,0.5) {Entrée\\ $N_e$};
\draw[-{Stealth[length=3mm]}, very thick, poppink] (11.7,0) -- (12.6,0);
\node[poppink, font=\small\bfseries, align=center] at (12.6,0.6) {Sortie\\ $N_s$};
% Formule
\node[popdark, font=\normalsize\bfseries] at (5.5,-2.6) {$k_{\text{global}} = k_1 \times k_2 \times k_3$ \quad et \quad $N_s = N_e \times k_{\text{global}}$};
\end{tikzpicture}
\legende{Train de transmission à trois étages (deux engrenages et une poulie-courroie ouverte). Le rapport global est le produit des rapports de chaque étage. Ici, deux engrenages externes inversent le sens deux fois, la courroie ouverte ne l'inverse pas : la sortie tourne finalement dans le même sens que l'entrée (flèches d'entrée et de sortie parallèles).}
\end{popfigure}

\begin{exoresolu}[Moteur, engrenage puis poulies : calcul complet]
\textbf{Énoncé.} Un moteur tourne à $N_e = \SI{1500}{tr/min}$. Il entraîne d'abord un engrenage de roues $Z_1 = 15$ dents (menante) et $Z_2 = 45$ dents (menée), puis un système poulies-courroie ouverte de diamètres $D_1 = \SI{50}{mm}$ (menante) et $D_2 = \SI{100}{mm}$ (menée).\par
1. Calculer le rapport de chaque étage puis le rapport global.\par
2. Calculer la vitesse de sortie $N_s$.\par
3. Le système est-il un réducteur ou un multiplicateur ? Comparer le sens de rotation de la sortie à celui de l'entrée.
\tcblower
\textbf{Solution détaillée.}\par
1. Rapports de chaque étage :
\begin{align*}
k_1 &= \frac{Z_1}{Z_2} = \frac{15}{45} = \frac{1}{3} \\
k_2 &= \frac{D_1}{D_2} = \frac{50}{100} = \frac{1}{2} \\
k_{\text{global}} &= k_1 \times k_2 = \frac{1}{3} \times \frac{1}{2} = \frac{1}{6}
\end{align*}
2. Vitesse de sortie :
\[ N_s = N_e \times k_{\text{global}} = 1500 \times \frac{1}{6} = \SI{250}{tr/min}. \]
3. Comme $k_{\text{global}} = \dfrac{1}{6} < 1$, le système est un \textbf{réducteur} : la vitesse est divisée par 6, et le couple disponible en sortie est multiplié (environ) par 6. Pour le sens : l'engrenage externe inverse le sens une fois, la courroie ouverte le conserve (même sens). Nombre d'inversions : 1 (impair), donc la sortie tourne en \textbf{sens inverse} de l'entrée.
\end{exoresolu}

\begin{exercice}[Entraînement : Réducteur de moulin à maïs]
Un moulin à maïs est entraîné par un moteur thermique tournant à $N_e = \SI{1800}{tr/min}$. La transmission comprend deux étages :
\begin{enumerate}
\item Un engrenage externe : pignon menant $Z_1 = 12$ dents, roue menée $Z_2 = 48$ dents.
\item Une transmission par chaîne : pignon d'entrée (sur l'arbre de la roue $Z_2$) $Z_3 = 18$ dents, pignon de sortie (sur l'arbre de la meule) $Z_4 = 36$ dents.
\end{enumerate}
\begin{enumerate}
\item Calculer le rapport de transmission global $k_{\text{global}}$.
\item Déterminer la vitesse de rotation de la meule $N_s$ (en tr/min).
\item Préciser si l'ensemble est un réducteur ou un multiplicateur.
\item Déterminer le sens de rotation de la meule par rapport au moteur (même sens ou sens inverse). Justifier en comptant les inversions.
\end{enumerate}
\end{exercice}

\begin{attention}[Pièges fréquents sur le rapport de transmission]
\begin{itemize}
\item \textbf{Ne jamais additionner} les rapports des étages : on les \textbf{multiplie}.
\item Pour les engrenages et les chaînes, on utilise les \textbf{nombres de dents} $Z$ ; pour les poulies-courroies, les \textbf{diamètres} $D$. Ne mélange pas les deux !
\item Une roue intermédiaire (roue « folle ») ne change pas la valeur de $k$, mais elle change le \textbf{sens} de rotation final.
\item Vérifie toujours la cohérence : si tu trouves $k > 1$ avec une petite roue menante et une grande roue menée, c'est forcément faux — une grande roue menée tourne toujours plus lentement.
\end{itemize}
\end{attention}

\begin{savaistu}[La boîte de vitesses de ta moto-taxi]
Quand le « bendskineur » passe les rapports, il sélectionne en réalité différents couples de pignons : en première, $k$ est très petit (forte réduction, beaucoup de couple pour démarrer en côte) ; en cinquième, $k$ se rapproche de 1 (grande vitesse sur le goudron). Le même moteur peut ainsi produire soit de la force, soit de la vitesse — jamais les deux à la fois. C'est la loi du rapport de transmission appliquée chaque jour sur les routes du Cameroun !
\end{savaistu}

\begin{aretenir}[L'essentiel de la section]
\begin{itemize}
\item Rapport de transmission : $k = \dfrac{N_{\text{sortie}}}{N_{\text{entrée}}}$, sans unité.
\item Engrenage : $k = \dfrac{Z_{\text{entrée}}}{Z_{\text{sortie}}}$, sens inversé à chaque engrenage externe.
\item Poulies-courroie : $k = \dfrac{D_{\text{entrée}}}{D_{\text{sortie}}}$, même sens (flèches parallèles) si courroie ouverte, sens inversés (flèches anti-parallèles) si croisée ; risque de glissement.
\item Chaîne-pignons : $k = \dfrac{Z_{\text{entrée}}}{Z_{\text{sortie}}}$, même sens (flèches parallèles), sans glissement.
\item Plusieurs étages : $k_{\text{global}} = k_1 \times k_2 \times \dots \times k_n$ ; le sens final dépend de la parité du nombre d'inversions (engrenages externes + courroies croisées).
\item $k < 1$ : réducteur (vitesse diminue, couple augmente) ; $k > 1$ : multiplicateur (vitesse augmente, couple diminue).
\end{itemize}
\end{aretenir}

\coursec{Les transmissions par contact direct : les engrenages}

Dans la section précédente, nous avons appris à calculer le \cle{rapport de transmission} $k$ entre deux roues en rotation : il suffit de comparer leurs nombres de dents ou leurs diamètres. Mais une question reste ouverte : \emph{comment} transmettre concrètement le mouvement d'un arbre moteur vers un arbre récepteur ? La réponse la plus directe, celle que l'on trouve dans la boîte de vitesses d'une moto-taxi, dans le réducteur d'un moulin à maïs ou dans une vieille horloge, c'est l'\cle{engrenage} : deux roues dentées dont les dents s'engrènent directement, sans courroie ni chaîne. On parle alors de \cle{transmission par contact direct}.

\courssub{Qu'est-ce qu'un engrenage ?}

\begin{definition}[Engrenage]
Un \cle{engrenage} est un ensemble de deux roues dentées (ou plus) dont les dents s'engrènent l'une dans l'autre. La plus petite roue s'appelle le \cle{pignon}, la plus grande s'appelle la \cle{roue} (ou la \cle{couronne}). Lorsque le pignon tourne, ses dents \emph{poussent} celles de la roue : le mouvement de rotation est transmis par contact direct.
\end{definition}

\begin{experience}[Observation : deux roues qui s'engrènent]
\textbf{Matériel :} deux engrenages en plastique de jouet (ou un vieux mécanisme d'horloge), un support.

\textbf{Protocole :} monte les deux roues dentées sur le support de façon que leurs dents s'engrènent. Fais tourner lentement le petit pignon à la main et observe la grande roue. Inverse ensuite les rôles.

\textbf{Observations :}
\begin{itemize}
\item les deux roues tournent en \textbf{sens opposés} ;
\item quand le petit pignon fait un tour complet, la grande roue n'a pas fini son tour : elle tourne \textbf{plus lentement} ;
\item il n'y a \textbf{aucun glissement} : la transmission est rigoureuse, comme si les deux roues étaient « cousues » l'une à l'autre.
\end{itemize}

\textbf{Interprétation :} les dents transmettent le mouvement sans glissement. Le rapport de transmission est donc parfaitement défini par les nombres de dents : $k = \dfrac{Z_1}{Z_2}$ (voir section 2). Le signe de $k$ est \textbf{négatif} pour un engrenage extérieur, car les sens de rotation sont opposés.
\end{experience}

\begin{aretenir}[Avantages et limites du contact direct]
\begin{itemize}
\item \textbf{Avantages :} aucun glissement (rapport exact), grande puissance transmissible, encombrement réduit, bon rendement.
\item \textbf{Limites :} les axes doivent être \textbf{proches} (petit entraxe), la fabrication est précise et coûteuse, il faut une \textbf{lubrification} et un carter de protection.
\end{itemize}
\end{aretenir}

\courssub{Le vocabulaire de la denture : module, cercle primitif, pas}

Pour fabriquer deux roues qui engrènent correctement, il ne suffit pas de « dessiner des dents » : il faut des dents \emph{normalisées}, dont toutes les dimensions se déduisent d'un seul nombre : le \cle{module}.

\begin{definition}[Module $m$]
Le \cle{module} $m$ (exprimé en millimètres) est la caractéristique de dimensionnement des dents d'un engrenage. Il est défini par :
\[ m = \dfrac{D}{Z} \]
où $D$ est le \cle{diamètre primitif} (diamètre du cercle primitif, en mm) et $Z$ le nombre de dents. Les modules sont \textbf{normalisés} (norme ISO) : 1 ; 1,25 ; 1,5 ; 2 ; 2,5 ; 3 ; 4 ; 5\dots{} Plus le module est grand, plus les dents sont grosses et résistantes.
\end{definition}

\begin{definition}[Cercle primitif et pas]
Le \cle{cercle primitif} est le cercle imaginaire sur lequel les deux roues « rouleraient sans glisser » si elles étaient de simples disques en contact. Le \cle{pas} $p$ est la distance entre deux dents consécutives, mesurée sur le cercle primitif. Il vaut :
\[ p = \pi \cdot m \]
Sur une dent, on distingue la \cle{tête} (partie au-dessus du cercle primitif), le \cle{pied} (partie en dessous) et les \cle{flancs} (profils latéraux en contact). Le \cle{jeu de fonctionnement} (ou \emph{backlash}) est le petit espace laissé entre les dents en prise : il est indispensable pour éviter le coincement, mais un jeu excessif trahit l'usure.
\end{definition}

\begin{popfigure}
\begin{center}
\begin{tikzpicture}[scale=1.1]
% Cercle primitif
\draw[popblue, dashed, thick] (0,0) circle (1.8);
\node[popblue, font=\small] at (2.6,1.5) {cercle primitif};
\draw[popblue, ->] (2.55,1.3) -- (1.35,1.15);
% Une dent en avant
\fill[poporangeL, draw=popdark, thick] (-0.55,1.62) -- (-0.45,2.55) -- (-0.15,2.75) -- (0.15,2.75) -- (0.45,2.55) -- (0.55,1.62) -- cycle;
% Dents voisines (simplifiées)
\fill[popblueL, draw=popdark] (0.75,1.5) -- (0.85,2.3) -- (1.1,2.45) -- (1.35,2.3) -- (1.4,1.35) -- cycle;
\fill[popblueL, draw=popdark] (-1.4,1.35) -- (-1.35,2.3) -- (-1.1,2.45) -- (-0.85,2.3) -- (-0.75,1.5) -- cycle;
% Annotations
\draw[<->, thick, popgreen] (-0.45,3.0) -- (0.45,3.0);
\node[popgreen, font=\small] at (0,3.25) {t\^ete};
\draw[<->, thick, poppurple] (0.6,1.05) -- (1.35,0.95);
\node[poppurple, font=\small] at (1.9,0.8) {pied};
\node[popdark, font=\small] at (-1.15,2.9) {flanc};
\draw[popdark, ->] (-1.05,2.7) -- (-0.5,2.3);
% Pas p
\draw[<->, thick, poppink] (0,2.9) -- (1.12,2.62);
\node[poppink, font=\small] at (1.35,3.05) {pas $p=\pi m$};
% Centre
\fill (0,0) circle (0.06);
\draw[<->, thick] (0,0) -- (1.27,-1.27);
\node[font=\small] at (0.75,-0.85) {$D/2$};
\end{tikzpicture}
\end{center}
\legende{Profil de dent d'engrenage : tête, pied, flancs, cercle primitif de diamètre $D$, pas $p = \pi m$. Toutes les dimensions de la dent se calculent à partir du module $m$.}
\end{popfigure}

\begin{propriete}[Condition de maillage]
Deux roues dentées ne peuvent engrèner correctement que si elles ont \textbf{le même module $m$} et donc \textbf{le même pas $p = \pi m$}. C'est la \cle{condition de maillage}. Deux engrenages de modules différents ne peuvent jamais fonctionner ensemble : les dents se coincent ou se brisent.
\end{propriete}

\begin{exemplebox}[Calculs autour d'un engrenage]
Un pignon de $Z_1 = 18$ dents a un module $m = \SI{2}{mm}$. Il engrène avec une roue de $Z_2 = 54$ dents.
\begin{itemize}
\item Diamètres primitifs : $D_1 = m \cdot Z_1 = 2 \times 18 = \SI{36}{mm}$ et $D_2 = m \cdot Z_2 = 2 \times 54 = \SI{108}{mm}$.
\item Rapport de transmission : $k = \dfrac{Z_1}{Z_2} = \dfrac{18}{54} = \dfrac{1}{3}$ : la roue tourne 3 fois moins vite que le pignon (c'est un \textbf{réducteur}).
\item Entraxe (distance entre les axes) : $a = \dfrac{D_1 + D_2}{2} = \dfrac{36 + 108}{2} = \SI{72}{mm}$.
\end{itemize}
\end{exemplebox}

\courssub{Les engrenages droits (cylindriques à denture droite)}

\begin{definition}[Engrenage droit]
Un \cle{engrenage droit} est un engrenage dont les deux axes sont \textbf{parallèles} et dont les dents sont \textbf{parallèles à l'axe} de rotation. C'est le type d'engrenage le plus simple et le plus répandu.
\end{definition}

L'engrenage droit est facile à fabriquer, économique et ne produit \textbf{aucun effort axial} (les dents ne poussent pas les roues le long de leurs axes). En revanche, chaque dent entre en contact \emph{brutalement} sur toute sa largeur : cela produit du \textbf{bruit} et des vibrations, surtout à \textbf{haute vitesse}. C'est pourquoi on réserve les engrenages droits aux vitesses modérées : boîtes de vitesses de machines-outils, réducteurs de moulins, mécanismes d'horlogerie.

\courssub{Les engrenages hélicoïdaux}

\begin{definition}[Engrenage hélicoïdal]
Un \cle{engrenage hélicoïdal} est un engrenage à axes \textbf{parallèles} dont les dents sont \textbf{inclinées} par rapport à l'axe, suivant une hélice. Le contact entre les dents se fait alors \textbf{progressivement}, d'un bout de la dent à l'autre.
\end{definition}

Grâce à ce contact progressif, l'engrenage hélicoïdal est beaucoup plus \textbf{silencieux} et plus doux que l'engrenage droit : c'est lui que l'on utilise dans les boîtes de vitesses des voitures et des motos modernes. Mais l'inclinaison des dents a un prix : elle crée une \cle{poussée axiale} $F_{ax}$, c'est-à-dire un effort qui pousse chaque roue le long de son axe.

\begin{attention}[Poussée axiale des engrenages hélicoïdaux]
La denture inclinée génère un \textbf{effort axial} $F_{ax}$ qui sollicite les \textbf{roulements} des arbres. Il faut donc prévoir des \textbf{butées} ou des roulements adaptés. Pour supprimer cet inconvénient, on utilise parfois des dentures \textbf{en chevrons} (double hélice en V) : les deux poussées axiales s'annulent mutuellement. Ne jamais monter un engrenage hélicoïdal sur un montage prévu pour un engrenage droit sans vérifier les butées !
\end{attention}

\begin{popfigure}
\begin{center}
\begin{tikzpicture}[scale=0.9]
% ===== Engrenage droit =====
\begin{scope}
% Roue 1 (vue de face simplifiée : cylindre)
\draw[thick, fill=popblueL] (-2.2,0) ellipse (0.35 and 1.0);
\draw[thick, fill=popblueL!60] (-1.2,-1.0) -- (-2.2,-1.0) arc (-90:90:0.35 and 1.0) -- (-1.2,1.0) arc (90:270:0.35 and 1.0);
\draw[thick] (-1.2,0) ellipse (0.35 and 1.0);
% dents droites (paralleles a l'axe)
\foreach \y in {-0.7,-0.35,0,0.35,0.7} {\draw[popblue, thick] (-2.2,\y) -- (-1.2,\y);}
% axe 1
\draw[popdark, dashed, thick] (-2.7,0) -- (-0.5,0);
% Roue 2
\draw[thick, fill=poporangeL] (0.2,1.6) ellipse (0.35 and 0.7);
\draw[thick, fill=poporangeL!60] (1.2,0.9) -- (0.2,0.9) arc (-90:90:0.35 and 0.7) -- (1.2,2.3) arc (90:270:0.35 and 0.7);
\draw[thick] (1.2,1.6) ellipse (0.35 and 0.7);
\foreach \y in {1.15,1.6,2.05} {\draw[poporange, thick] (0.2,\y) -- (1.2,\y);}
\draw[popdark, dashed, thick] (-0.3,1.6) -- (1.9,1.6);
\node[font=\small\bfseries] at (-0.5,-1.7) {Engrenage droit};
\node[font=\small, popdark] at (-0.5,-2.15) {axes parall\`eles, dents // axe};
\end{scope}
% ===== Engrenage helicoidal =====
\begin{scope}[xshift=6.5cm]
\draw[thick, fill=popgreenL] (-2.2,0) ellipse (0.35 and 1.0);
\draw[thick, fill=popgreenL!60] (-1.2,-1.0) -- (-2.2,-1.0) arc (-90:90:0.35 and 1.0) -- (-1.2,1.0) arc (90:270:0.35 and 1.0);
\draw[thick] (-1.2,0) ellipse (0.35 and 1.0);
% dents inclinees
\foreach \y in {-0.7,-0.35,0,0.35,0.7} {\draw[popgreen!70!black, thick] (-2.2,\y) -- (-1.2,{\y+0.3});}
\draw[popdark, dashed, thick] (-2.7,0) -- (-0.5,0);
\draw[thick, fill=poppinkL] (0.2,1.6) ellipse (0.35 and 0.7);
\draw[thick, fill=poppinkL!60] (1.2,0.9) -- (0.2,0.9) arc (-90:90:0.35 and 0.7) -- (1.2,2.3) arc (90:270:0.35 and 0.7);
\draw[thick] (1.2,1.6) ellipse (0.35 and 0.7);
\foreach \y in {1.15,1.6,2.05} {\draw[poppink!70!black, thick] (0.2,\y) -- (1.2,{\y+0.25});}
\draw[popdark, dashed, thick] (-0.3,1.6) -- (1.9,1.6);
% effort axial
\draw[->, very thick, poppurple] (-1.7,0) -- (-0.7,0);
\node[poppurple, font=\small] at (-0.55,0.3) {$F_{ax}$};
\node[font=\small\bfseries] at (-0.5,-1.7) {Engrenage h\'elico\"idal};
\node[font=\small, popdark] at (-0.5,-2.15) {dents inclin\'ees, pouss\'ee axiale $F_{ax}$};
\end{scope}
\end{tikzpicture}
\end{center}
\legende{Comparaison : engrenage droit (dents parallèles à l'axe, pas d'effort axial) et engrenage hélicoïdal (dents inclinées, contact progressif et silencieux, mais apparition d'une poussée axiale $F_{ax}$).}
\end{popfigure}

\courssub{Les engrenages coniques}

\begin{definition}[Engrenage conique]
Un \cle{engrenage conique} est un engrenage dont les axes sont \textbf{concourants} (ils se coupent), le plus souvent à \textbf{angle droit} (90°). Les dents sont taillées sur des cônes : c'est pourquoi on parle de pignon et de couronne \emph{coniques}.
\end{definition}

L'engrenage conique permet donc de transmettre un mouvement de rotation \textbf{entre deux arbres perpendiculaires} : il « tourne le coin ». On le trouve dans le \textbf{différentiel} des voitures et des camions, dans les perceuses d'angle, dans les renvois d'angle des machines agricoles. Sa fabrication est délicate : les deux cônes doivent être réglés avec précision pour que les dents portent correctement.

\begin{popfigure}
\begin{center}
\begin{tikzpicture}[scale=1.0]
% Axe horizontal (pignon)
\draw[popdark, dashed, thick] (-3.2,0) -- (2.2,0);
% Pignon conique (a droite)
\fill[poporangeL, draw=popdark, thick] (0.1,-0.85) -- (1.7,-0.35) -- (1.7,0.35) -- (0.1,0.85) -- cycle;
\draw[poporange, thick] (0.35,-0.72) -- (0.35,0.72);
\draw[poporange, thick] (0.7,-0.62) -- (0.7,0.62);
\draw[poporange, thick] (1.05,-0.52) -- (1.05,0.52);
\draw[poporange, thick] (1.4,-0.42) -- (1.4,0.42);
\node[font=\small, poporange!70!black] at (1.0,-1.2) {pignon conique};
% Axe vertical (couronne)
\draw[popdark, dashed, thick] (0,-2.6) -- (0,2.4);
% Couronne conique (en haut)
\fill[popblueL, draw=popdark, thick] (-0.85,-0.1) -- (-0.4,-1.6) -- (0.4,-1.6) -- (0.85,-0.1) -- cycle;
\draw[popblue, thick] (-0.72,-0.35) -- (0.72,-0.35);
\draw[popblue, thick] (-0.62,-0.7) -- (0.62,-0.7);
\draw[popblue, thick] (-0.52,-1.05) -- (0.52,-1.05);
\draw[popblue, thick] (-0.44,-1.4) -- (0.44,-1.4);
\node[font=\small, popblue!70!black] at (2.1,-1.5) {couronne conique};
% Angle 90
\draw[popgreen, thick] (0.55,0) arc (0:-90:0.55);
\node[popgreen, font=\small] at (0.85,-0.75) {$90^{\circ}$};
% Point de concours
\fill[popdark] (0,0) circle (0.07);
\node[font=\small, popdark] at (-1.15,0.28) {concours des axes};
% Fleches de rotation
\draw[->, very thick, poppurple] (-2.6,0.35) arc (120:60:0.8);
\draw[->, very thick, poppurple] (0.35,-2.2) arc (-30:30:0.8);
\end{tikzpicture}
\end{center}
\legende{Engrenage conique : le pignon et la couronne ont des axes \textbf{concourants à 90°}. Les dents sont taillées sur des cônes dont les sommets coïncident au point de concours des axes. Application : le différentiel automobile.}
\end{popfigure}

\begin{savaistu}[Le différentiel de ta voiture]
Quand une voiture tourne dans un virage, la roue extérieure parcourt un chemin \emph{plus long} que la roue intérieure : elle doit donc tourner \emph{plus vite}. Si les deux roues étaient solidaires du même arbre, elles patineraient et les pneus s'useraient très vite. La solution ? Le \cle{différentiel}, un mécanisme ingénieux qui utilise des \textbf{engrenages coniques} (appelés \emph{planétaires} et \emph{satellites}) : il transmet la puissance du moteur aux deux roues motrices tout en les autorisant à tourner à des \textbf{vitesses différentes}. Regarde sous un camion ou un pick-up à Douala : le « nez de pont » bombé entre les roues arrière abrite précisément ce mécanisme.
\end{savaistu}

\courssub{La roue et vis sans fin : un aperçu}

Il existe un quatrième cas de transmission par contact direct, que nous étudierons en détail dans la section 5 : le système \cle{roue et vis sans fin}. Ici, les axes sont \textbf{croisés} (généralement à 90°, sans se couper). Une \emph{vis} (semblable à une grosse vis filetée) engrène avec une roue dentée. Ce système permet d'obtenir un rapport de transmission \textbf{très élevé} en une seule étape (par exemple $k = 1/40$) et peut être \textbf{irréversible} : seule la vis peut entraîner la roue, jamais l'inverse. C'est le principe des treuils et de certains vérins.

\begin{methode}[Identifier le type d'engrenage]
Face à un mécanisme ou à un schéma, pose-toi deux questions dans l'ordre :
\begin{enumerate}
\item \textbf{Les axes sont-ils parallèles ?}
\begin{itemize}
\item Oui, et les dents sont parallèles à l'axe $\Rightarrow$ engrenage \textbf{droit}.
\item Oui, et les dents sont inclinées $\Rightarrow$ engrenage \textbf{hélicoïdal}.
\end{itemize}
\item \textbf{Les axes sont-ils perpendiculaires ?}
\begin{itemize}
\item Ils se coupent (concourants) $\Rightarrow$ engrenage \textbf{conique}.
\item Ils ne se coupent pas (croisés) $\Rightarrow$ système \textbf{roue et vis sans fin}.
\end{itemize}
\end{enumerate}
\end{methode}

\begin{center}
\begin{tabularx}{\linewidth}{|l|X|X|X|}
\hline
\rowcolor{popblueL} \textbf{Type} & \textbf{Position des axes} & \textbf{Avantage principal} & \textbf{Exemple d'usage} \\
\hline
Droit & Parallèles & Simple, économique, pas d'effort axial & Réducteur de moulin, horloge \\
\hline
Hélicoïdal & Parallèles & Silencieux, contact progressif & Boîte de vitesses de voiture \\
\hline
Conique & Concourants (souvent 90°) & Change la direction de l'axe & Différentiel, perceuse d'angle \\
\hline
Roue et vis sans fin & Croisés (sans se couper) & Très grand rapport, irréversible & Treuil, vérin \\
\hline
\end{tabularx}
\end{center}

\courssub{Lubrification et entretien des engrenages}

Des dents en acier qui frottent l'une contre l'autre à grande vitesse s'usent et chauffent vite. La \cle{lubrification} est donc vitale pour un engrenage :
\begin{itemize}
\item \textbf{Graisse} : pour les mécanismes ouverts ou tournant lentement (treuil, portail, presse manuelle) ; elle reste en place sans carter étanche.
\item \textbf{Huile} : pour les mécanismes fermés et rapides ou fortement chargés (boîte de vitesses de moto, pont de camion) ; l'huile circule, refroidit et évacue les particules d'usure. C'est pour cela qu'il faut \emph{vidanger} régulièrement la boîte et le pont d'un véhicule.
\end{itemize}

\begin{attention}[Signes d'usure à surveiller]
Un engrenage usé se détecte à l'oreille et à la main :
\begin{itemize}
\item un \textbf{bruit anormal} (claquement, sifflement, ronronnement rauque) ;
\item un \textbf{jeu excessif} (backlash trop grand) : la roue « bouge » avant d'entraîner l'autre ;
\item de la \textbf{limaille} métallique dans l'huile de vidange.
\end{itemize}
Continuer à utiliser un engrenage usé, c'est risquer la \textbf{rupture des dents} et l'arrêt complet de la machine. On remplace toujours un pignon usé \emph{avec} sa roue (les deux se sont usés ensemble), et jamais un seul élément neuf contre un élément usé.
\end{attention}

\begin{exoresolu}[Dimensionnement d'un engrenage de moulin]
Le moteur électrique d'un moulin à maïs tourne à $N_1 = \SI{1500}{tr/min}$. Il entraîne, par un engrenage droit de module $m = \SI{2,5}{mm}$, la meule qui doit tourner à $N_2 = \SI{500}{tr/min}$. Le pignon moteur a $Z_1 = 20$ dents.
\begin{enumerate}
\item Calcule le rapport de transmission $k$ nécessaire.
\item Détermine le nombre de dents $Z_2$ de la roue.
\item Calcule les diamètres primitifs $D_1$ et $D_2$, puis l'entraxe $a$.
\end{enumerate}
\tcblower
\textbf{Solution détaillée.}
\begin{enumerate}
\item Le rapport de transmission est le quotient des vitesses :
\[ k = \dfrac{N_2}{N_1} = \dfrac{500}{1500} = \dfrac{1}{3}. \]
C'est un réducteur ($k < 1$).
\item Pour un engrenage, $k = \dfrac{Z_1}{Z_2}$, donc :
\[ Z_2 = \dfrac{Z_1}{k} = \dfrac{20}{1/3} = 20 \times 3 = 60 \text{ dents}. \]
\item Les diamètres primitifs se calculent avec $D = m \cdot Z$ :
\begin{align*}
D_1 &= m \cdot Z_1 = 2{,}5 \times 20 = \SI{50}{mm},\\
D_2 &= m \cdot Z_2 = 2{,}5 \times 60 = \SI{150}{mm}.
\end{align*}
L'entraxe vaut alors :
\[ a = \dfrac{D_1 + D_2}{2} = \dfrac{50 + 150}{2} = \SI{100}{mm}. \]
\end{enumerate}
\textbf{Vérification :} $k = \dfrac{D_1}{D_2} = \dfrac{50}{150} = \dfrac{1}{3}$ : le résultat est cohérent.
\end{exoresolu}

\begin{attention}[Pièges fréquents au BEPC]
\begin{itemize}
\item \textbf{Ne pas inverser le rapport :} $k = \dfrac{Z_{\text{menant}}}{Z_{\text{menée}}}$. Si le pignon mène la roue, $k = \dfrac{Z_1}{Z_2}$ ; si c'est la roue qui mène, le rapport s'inverse.
\item \textbf{Ne pas oublier la condition de maillage :} deux roues de modules différents ne peuvent pas engrèner, même si les nombres de dents « semblent » convenir.
\item \textbf{Unités :} le module $m$ s'exprime en mm ; $D = m \cdot Z$ donne directement un diamètre en mm. Ne mélange pas mm et cm dans un même calcul.
\item \textbf{Sens de rotation :} dans un engrenage extérieur simple, les deux roues tournent en \textbf{sens opposés} ; le rapport $k$ est donc \emph{négatif} si l'on tient compte du signe.
\end{itemize}
\end{attention}

\begin{exercice}[À toi de jouer]
Une perceuse d'angle utilise un engrenage conique dont le pignon (côté moteur) a $Z_1 = 12$ dents et la couronne (côté mandrin) a $Z_2 = 36$ dents.
\begin{enumerate}
\item De quel type d'engrenage s'agit-il ? Justifie à partir de la position des axes.
\item Calcule le rapport de transmission $k$.
\item Le moteur tourne à $N_1 = \SI{3000}{tr/min}$. Calcule la vitesse de rotation $N_2$ du mandrin.
\item Le module est $m = \SI{1,5}{mm}$. Calcule le diamètre primitif de la couronne.
\end{enumerate}
\end{exercice}

\begin{corrige}[Exercice : perceuse d'angle]
\begin{enumerate}
\item Il s'agit d'un \textbf{engrenage conique} : dans une perceuse d'angle, l'axe du moteur et l'axe du mandrin sont \textbf{concourants à 90°}, ce qui est la signature des engrenages coniques.
\item $k = \dfrac{Z_1}{Z_2} = \dfrac{12}{36} = \dfrac{1}{3}$.
\item $N_2 = k \cdot N_1 = \dfrac{1}{3} \times 3000 = \SI{1000}{tr/min}$. Le mandrin tourne 3 fois moins vite que le moteur (réducteur), ce qui augmente le couple de perçage.
\item $D_2 = m \cdot Z_2 = 1{,}5 \times 36 = \SI{54}{mm}$.
\end{enumerate}
\end{corrige}

\begin{aretenir}[L'essentiel de la section]
\begin{itemize}
\item Un \textbf{engrenage} transmet le mouvement par \textbf{contact direct} entre des dents : pas de glissement, rapport exact $k = \dfrac{Z_1}{Z_2}$.
\item Le \textbf{module} $m = \dfrac{D}{Z}$ dimensionne les dents ; le \textbf{pas} vaut $p = \pi m$. Deux roues n'engrènent que si elles ont le \textbf{même module} (condition de maillage).
\item \textbf{Droit} : axes parallèles, dents parallèles à l'axe, simple mais bruyant à haute vitesse.
\item \textbf{Hélicoïdal} : axes parallèles, dents inclinées, silencieux, mais \textbf{poussée axiale} à compenser (butées, chevrons).
\item \textbf{Conique} : axes concourants (souvent 90°) ; applications : différentiel, perceuse d'angle.
\item \textbf{Roue et vis sans fin} : axes croisés, très grand rapport, irréversibilité possible (détaillé en section 5).
\item Tout engrenage exige une \textbf{lubrification} (graisse ou huile) et une surveillance de l'usure (jeu, bruit, limaille).
\end{itemize}
\end{aretenir}

\coursec{Les transmissions par liaison flexible : courroies et chaînes}

\courssub{Transition depuis les engrenages}

Dans la section précédente, nous avons vu que les \textbf{engrenages} transmettent le mouvement par \emph{contact direct} entre dents. Cela impose une distance très courte entre les axes (entraxe imposé par le module et le nombre de dents) et nécessite une lubrification soignée. Or, dans de nombreux systèmes réels — moto-taxi, groupe électrogène, ventilateur de cuisine, compresseur de climatisation —, les arbres sont \textbf{éloignés} de plusieurs dizaines de centimètres, voire de mètres. Les engrenages deviennent alors trop encombrants, trop coûteux ou trop bruyants.

C'est là qu'interviennent les \textbf{liaisons flexibles} : une courroie ou une chaîne enroule les poulies (ou pignons) et assure la transmission à distance. Nous allons étudier leurs principes, leurs variantes, leurs avantages et leurs limites, en gardant toujours en tête le \textbf{rapport de transmission} $k = \frac{\omega_{\text{mené}}}{\omega_{\text{moteur}}}$ et le \textbf{rendement} $\eta$.

\begin{experience}[Observation de courroies et chaînes du quotidien]
\textbf{Matériel :} courroie plate (vieille courroie de machine à coudre), courroie trapézoïdale (récupérée sur un ventilateur ou un compresseur), courroie dentée (distribution de moto), chaîne à rouleaux + pignon (moto-taxi), règle, loupe.

\textbf{Protocole :}
\begin{enumerate}
    \item Examine chaque élément à la loupe. Note la forme de la section, la présence de dents, l'état de surface.
    \item Pour la courroie trapézoïdale, mesure l'angle du trapèze (environ \SI{38}{\degree}--\SI{40}{\degree}).
    \item Pour la chaîne, identifie : plaque extérieure, plaque intérieure, goupille, rouleau. Compte le pas (distance entre deux goupilles successives).
    \item Essaie de faire glisser une courroie plate sur une poulie lisse, puis une courroie trapézoïdale dans sa gorge : sens-tu la différence d'adhérence ?
\end{enumerate}

\textbf{Observations :}
\begin{itemize}
    \item Courroie plate : section rectangulaire, surface lisse, glisse facilement sur poulie lisse.
    \item Courroie trapézoïdale : section en V, flancs légèrement rugueux, \emph{coince} dans la gorge de la poulie.
    \item Courroie dentée : dents trapézoïdales ou arrondies sur la face intérieure, s'engrènent dans les gorges de la poulie dentée.
    \item Chaîne : articulation rouleau/goupille, pas constant, s'engrène sur pignons à dents.
\end{itemize}

\textbf{Interprétation :} La géométrie de la section détermine le mode de transmission de l'effort : \textbf{adhérence par frottement} (plate, trapézoïdale) ou \textbf{engrènement formel} (dentée, chaîne). Cela conditionne le \textbf{glissement}, le \textbf{rendement} et la \textbf{précision du rapport $k$}.
\legende{Observation comparée des quatre types de liaisons flexibles.}
\end{experience}

\courssub{Courroies plates : simplicité et grand entraxe}

\begin{definition}[Courroie plate]
Une \cle{courroie plate} est une bande flexible de section rectangulaire (largeur $l$, épaisseur $e$) qui s'enroule sur des \textbf{poulies lisses} (cylindriques ou légèrement bombées). La transmission de l'effort repose uniquement sur le \textbf{frottement} entre la courroie et la poulie.
\end{definition}

\begin{propriete}[Caractéristiques principales]
\begin{itemize}
    \item \textbf{Axes parallèles}, distance grande (entraxe $E$ jusqu'à \SIrange{5}{10}{\meter}).
    \item \textbf{Sens de rotation} : même sens (brins non croisés) ou sens opposés (brins croisés).
    \item \textbf{Glissement inévitable} : sous l'effet du couple, la courroie s'allonge élastiquement sur le brin tendu et se raccourcit sur le brin relâché $\Rightarrow$ vitesse de la courroie $\neq$ vitesse périphérique des poulies.
    \item \textbf{Rendement} $\eta \approx 95\,\%$ (pertes par frottement, glissement, hystérésis).
    \item \textbf{Puissance transmise} limitée (quelques \SI{}{\kilo\watt}) : risque de patinage si couple trop élevé.
    \item \textbf{Coût faible}, \textbf{silencieuse}, \textbf{pas de lubrification}.
\end{itemize}
\end{propriete}

\begin{attention}[Glissement $\neq$ patinage]
\begin{itemize}
    \item \textbf{Glissement (élastique)} : différence de vitesse \emph{intrinsèque} due à l'allongement élastique de la courroie sous la tension. Existe toujours, même en charge nominale. Rendement $< 100\,\%$.
    \item \textbf{Patinage (rupture d'adhérence)} : la courroie \emph{glisse franchement} sur la poulie parce que le couple dépasse la capacité de frottement (surcharge, courroie détendue, poulie mouillée/huilée). \textbf{Arrêt immédiat nécessaire}.
\end{itemize}
Au BEPC, on parle de \textbf{glissement} pour expliquer pourquoi $k_{\text{réel}} \neq k_{\text{théorique}}$.
\end{attention}

\begin{methode}[Calcul du rapport de transmission réel (avec glissement)]
Pour une courroie plate ou trapézoïdale, le glissement $g$ (en \%) se définit par :
\[
g = \frac{v_{\text{poulie motrice}} - v_{\text{courroie}}}{v_{\text{poulie motrice}}} \times 100
\]
Le rapport de transmission \textbf{réel} devient :
\[
k_{\text{réel}} = \frac{\omega_2}{\omega_1} = \frac{D_1}{D_2} \cdot (1 - g)
\]
où $D_1, D_2$ sont les \textbf{diamètres primitifs} (au niveau du neutre de la courroie).
\end{methode}

\courssub{Courroies trapézoïdales (V) : l'adhérence par coinçage}

\begin{definition}[Courroie trapézoïdale]
Une \cle{courroie trapézoïdale} (ou \emph{courroie V}) a une section en trapèze isocèle (angle au sommet $\alpha \approx \SI{38}{\degree}$--\SI{40}{\degree}). Elle s'enroule dans une \textbf{poulie à gorge} de même angle. L'effet de coinçage multiplie la force de frottement utile par $\frac{1}{\sin(\alpha/2)} \approx 3$ (pour $\alpha=\SI{38}{\degree}$) par rapport à une courroie plate de même tension.
\end{definition}

\begin{propriete}[Avantages sur la plate]
\begin{itemize}
    \item \textbf{Adhérence multipliée} $\Rightarrow$ transmission de \textbf{puissance bien supérieure} (jusqu'à \SIrange{50}{100}{\kilo\watt} selon sections).
    \item \textbf{Glissement réduit} $\Rightarrow$ rendement $\eta \approx 97\,\%$.
    \item \textbf{Auto-centrage} dans la gorge $\Rightarrow$ moins sensible au désalignement léger.
    \item \textbf{Entraxe} moyen (\SIrange{0.5}{3}{\meter} typique).
\end{itemize}
\end{propriete}

\begin{attention}[Usure et remplacement]
Une courroie trapézoïdale usée s'enfonce au fond de la gorge (contact sur le fond au lieu des flancs) $\Rightarrow$ perte totale de l'effet de coinçage $\Rightarrow$ patinage immédiat. \textbf{Ne jamais graisser une courroie V} (perte d'adhérence). Remplacer par jeu complet si plusieurs brins (courroies jumelées).
\end{attention}

\courssub{Courroies dentées (synchrone) : la précision sans glissement}

\begin{definition}[Courroie dentée]
Une \cle{courroie dentée} (ou \emph{courroie synchrone}) porte des dents sur sa face intérieure qui s'engrènent exactement dans les gorges d'une \textbf{poulie dentée}. Le \textbf{pas} $p$ (distance entre deux dents successives) est identique pour la courroie et la poulie.
\end{definition}

\begin{propriete}[Liaison kinématique stricte]
\begin{itemize}
    \item \textbf{Pas de glissement} : $k = \frac{Z_1}{Z_2} = \frac{D_1}{D_2}$ \textbf{exact} (comme un engrenage).
    \item \textbf{Rendement} $\eta \approx 98\,\%$ (frottement roulement uniquement).
    \item \textbf{Silencieuse} (pas de choc dent/dent si profil bien usiné).
    \item \textbf{Pas de lubrification}.
    \item \textbf{Entraxe} imposé par la longueur de courroie (standardisée).
\end{itemize}
\end{propriete}

\begin{aretenir}[Courroie dentée vs Chaîne]
\begin{center}
\begin{tabularx}{\linewidth}{|l|X|X|}\hline
\textbf{Critère} & \textbf{Courroie dentée} & \textbf{Chaîne à rouleaux} \\ \hline
Glissement & \textbf{Nul} ($k$ exact) & \textbf{Nul} ($k$ exact) \\ \hline
Bruit & Très faible & Modéré (claquement rouleaux) \\ \hline
Lubrification & \textbf{Aucune} & \textbf{Obligatoire} \\ \hline
Entraxe & Figé (longueur standard) & Réglable (maillons) \\ \hline
Puissance max & Moyenne ($\sim \SIrange{10}{20}{\kilo\watt}$) & Élevée ($\sim \SI{100}{\kilo\watt}+$) \\ \hline
Coût & Moyen & Moyen à élevé \\ \hline
\end{tabularx}
\end{center}
\end{aretenir}

\begin{savaistu}[La courroie de distribution du moteur thermique]
Dans tout moteur 4 temps (moto, voiture, groupe électrogène), une \textbf{courroie dentée} (ou parfois une chaîne) synchronise le \textbf{vilebrequin} (arbre moteur) et l'\textbf{arbre à cames} (commande des soupapes). Le rapport est \textbf{$k = 1/2$} : l'arbre à cames tourne \textbf{2 fois moins vite} que le vilebrequin (1 tour arbre à cames pour 2 tours vilebrequin = 1 cycle 4 temps).
\begin{itemize}
    \item Si la courroie casse $\Rightarrow$ perte de synchronisation $\Rightarrow$ chocs soupapes/pistons $\Rightarrow$ \textbf{casse moteur grave} (souvent irréparable).
    \item \textbf{Entretien :} remplacement préventif tous les \SIrange{60000}{100000}{\kilo\meter} ou 5 ans (selon constructeur). Au Cameroun, avec la poussière et la chaleur, on raccourcit l'intervalle.
\end{itemize}
\end{savaistu}

\courssub{Chaînes à rouleaux et pignons : robustesse et puissance}

\begin{definition}[Chaîne à rouleaux (type roller chain)]
Une \cle{chaîne à rouleaux} est une succession d'articulations \textbf{plaques extérieures -- goupille -- rouleau -- plaque intérieure}. Elle s'engrène sur des \textbf{pignons à dents} (profil adapté au pas $p$). Le pas standardisé (ex. 1/2", 5/8", 3/4") doit être identique pour chaîne et pignons.
\end{definition}

\begin{experience}[Démontage/observation d'une chaîne de moto]
\textbf{Matériel :} chaîne de moto usagée, pignon, maillon rapide, dérive-chaîne (ou marteau + pointeau), chiffon, lubrifiant chaîne.
\textbf{Protocole :}
\begin{enumerate}
    \item Nettoie la chaîne. Identifie un maillon complet : 2 plaques ext., 2 plaques int., 2 goupilles, 2 rouleaux.
    \item Mesure le \textbf{pas} $p$ : distance entre axes de deux goupilles successives (ex. \SI{15.875}{\milli\meter} = 5/8").
    \item Compte les dents du pignon moteur ($Z_1$) et de la couronne ($Z_2$).
    \item Vérifie l'alignement : pose une règle contre les flancs des deux pignons.
    \item Vérifie la tension : battement au milieu de l'entraxe.
    \item (Optionnel) Démaillonne un maillon rapide, observe l'usure des goupilles (ovalisation).
\end{enumerate}
\textbf{Interprétation :} La chaîne transforme le mouvement de rotation en translation (brin tendu) puis en rotation. Le \textbf{rouleau} roule sur la dent du pignon (frottement roulement $\ll$ frottement glissement). L'articulation goupille/rouleau subit des cycles de charge $\Rightarrow$ \textbf{usure par allongement} (pas qui augmente) $\Rightarrow$ risque de saut de dent.
\legende{Anatomie d'une chaîne à rouleaux et engrènement sur pignon.}
\end{experience}

\begin{propriete}[Caractéristiques principales]
\begin{itemize}
    \item \textbf{Pas de glissement} : $k = \frac{Z_1}{Z_2}$ exact (liaison forme).
    \item \textbf{Rendement} $\eta \approx 98\,\%$ (très bon).
    \item \textbf{Puissance élevée} : jusqu'à plusieurs centaines de \SI{}{\kilo\watt} (industriel).
    \item \textbf{Entraxe moyen} (\SIrange{0.5}{2}{\meter}), \textbf{axes //}.
    \item \textbf{Contraintes :} \textbf{Lubrification obligatoire} (graisse ou bain d'huile), \textbf{tension rigoureuse}, \textbf{alignement parfait} (tolérance $\pm \SI{0.5}{\milli\meter\per\meter}$).
    \item \textbf{Bruit} : claquements rouleaux/dents, vibrations.
    \item \textbf{Allongement progressif} (usure goupilles) $\Rightarrow$ tension à surveiller.
\end{itemize}
\end{propriete}

\begin{methode}[Vérifier et régler la tension d'une chaîne (moto, vélo, groupe)]
\begin{enumerate}
    \item Place le véhicule sur béquille centrale (roue arrière libre).
    \item Repère le point le plus tendu (tourne la roue).
    \item Mesure le \textbf{battement} $b$ (flèche verticale) au milieu de l'entraxe $E$.
    \item \textbf{Règle pratique :} $b \approx 1\,\% \text{ à } 2\,\% \text{ de } E$ (ex. $E=\SI{300}{\milli\meter} \Rightarrow b=\SIrange{3}{6}{\milli\meter}$).
    \item Si $b$ trop grand : desserre l'axe de roue, repousse la roue par les tirettes (également des deux côtés !), resserre l'axe au couple.
    \item \textbf{Vérifie l'alignement} : les repères de tirettes doivent être identiques des deux côtés.
\end{enumerate}
\end{methode}

\begin{attention}[Alignement des pignons : la cause n°1 de casse chaîne]
Un désalignement de \SI{1}{\degree} suffit à :
\begin{itemize}
    \item User les flancs des dents (pignon « en biais »).
    \item Pousser la chaîne hors des dents (saut de dent, déraillement).
    \item Créer des à-coups, vibrations, bruit anormal.
\end{itemize}
\textbf{Vérification :} Règle métallique ou laser posé sur les flancs des deux pignons. L'écart doit être $< \SI{0.5}{\milli\meter}$ sur \SI{300}{\milli\meter}. Sur moto, les \textbf{repères de tirettes} (graduations sur le bras oscillant) doivent coïncider exactement.
\end{attention}

\courssub{Tendeurs (galets) : indispensables pour la durée de vie}

\begin{definition}[Galet tendeur (ou tensionneur)]
Un \cle{galet tendeur} est une poulie libre (sur roulement) montée sur un bras mobile (ressort ou vis) qui appuie sur le \textbf{brin relâché} (courroie) ou le \textbf{brin tendu} (chaîne) pour maintenir la tension nominale.
\end{definition}

\begin{propriete}[Rôles du tendeur]
\begin{itemize}
    \item \textbf{Compenser l'allongement} progressif (courroie : fluage élastique/visqueux ; chaîne : usure goupilles).
    \item \textbf{Augmenter l'angle d'enroulement} sur la poulie/pignon moteur $\Rightarrow$ meilleure adhérence (courroies) ou engrènement sur plus de dents (chaîne).
    \item \textbf{Absorber les vibrations} et chocs de couple (galet à ressort).
    \item \textbf{Guider} la courroie/chaîne (galet lisse ou cannelé).
\end{itemize}
\end{propriete}

\begin{attention}[Position du tendeur]
\begin{itemize}
    \item \textbf{Courroie :} galet sur \textbf{brin relâché} (côté sortie poulie menée), \textbf{extérieur} de la courroie (pousse vers l'intérieur). Augmente l'enroulement sur les deux poulies.
    \item \textbf{Chaîne :} galet sur \textbf{brin tendu} (côté entrée pignon mené), \textbf{intérieur} de la chaîne (pousse vers l'extérieur). Évite que la chaîne ne « monte » sur les dents du pignon mené sous l'effet du couple.
    \item \textbf{Courroie dentée :} galet \textbf{cannelé} (denté) obligatoire pour ne pas abîmer les dents de la courroie.
\end{itemize}
\end{attention}

\courssub{Tableau comparatif et choix technologique}

\begin{popfigure}
\begin{tikzpicture}[scale=0.9, every node/.style={font=\small}]
% En-têtes
\matrix (tab) [matrix of nodes, nodes in empty cells,
    nodes={anchor=center, minimum height=8mm, text width=2.2cm, align=center},
    column 1/.style={text width=2.8cm, font=\bfseries, fill=popblue!20},
    column 2/.style={fill=popgreen!10},
    column 3/.style={fill=poporange!10},
    column 4/.style={fill=poppurple!10},
    column 5/.style={fill=poppink!10},
    column 6/.style={fill=popgold!10},
    row 1/.style={font=\bfseries\small, fill=popdark, text=white, minimum height=10mm},
    row sep=-\pgflinewidth, column sep=-\pgflinewidth,
    draw=popline, line width=0.5pt
]{
Critère & Courroie plate & Courroie V & Courroie dentée & Chaîne à rouleaux & Engrenages \\
Glissement & Oui (élastique) & Faible & \textbf{Nul} & \textbf{Nul} & \textbf{Nul} \\
Précision $k$ & Approchée & Approchée & \textbf{Exacte} & \textbf{Exacte} & \textbf{Exacte} \\
Rendement $\eta$ & $\sim 95\,\%$ & $\sim 97\,\%$ & $\sim 98\,\%$ & $\sim 98\,\%$ & $\sim 98-99\,\%$ \\
Puissance max & Faible ($<\SI{5}{\kilo\watt}$) & Moyenne ($<\SI{50}{\kilo\watt}$) & Moyenne ($<\SI{20}{\kilo\watt}$) & \textbf{Élevée} ($>\SI{100}{\kilo\watt}$) & \textbf{Très élevée} \\
Bruit & Très faible & Faible & Très faible & Modéré & Variable \\
Entretien & Tension occasionnelle & Tension + état flancs & Tension + alignement & \textbf{Lubrification + tension + alignement} & Lubrification \\
Entraxe & Grand ($>\SI{1}{\meter}$) & Moyen (\SIrange{0.5}{3}{\meter}) & Figé (longueur std) & Moyen (\SIrange{0.5}{2}{\meter}) & Court (imposé) \\
Coût & Faible & Faible & Moyen & Moyen/Élevé & Élevé \\
Lubrification & Non & Non & Non & \textbf{Oui} & \textbf{Oui} \\
};
% Lignes horizontales
\foreach \i in {1,...,10} {
    \draw[popline, line width=0.5pt] (tab-\i-1.north west) -- (tab-\i-6.north east);
}
\draw[popline, line width=0.8pt] (tab-1-1.north west) -- (tab-1-6.north east);
\draw[popline, line width=0.8pt] (tab-10-1.south west) -- (tab-10-6.south east);
% Lignes verticales
\foreach \j in {1,...,6} {
    \draw[popline, line width=0.5pt] (tab-1-\j.north west) -- (tab-10-\j.south west);
}
\draw[popline, line width=0.8pt] (tab-1-6.north east) -- (tab-10-6.south east);
\end{tikzpicture}
\legende{Tableau comparatif des systèmes de transmission par liaison flexible et engrenages. Le choix dépend de la puissance, de la précision requise, de l'entraxe, du budget et de la maintenance possible.}
\end{popfigure}

\begin{exemplebox}[Choix pour un ventilateur industriel]
\textbf{Cahier des charges :} Moteur \SI{4}{\kilo\watt}, \SI{1500}{tr/min} $\rightarrow$ ventilateur \SI{500}{tr/min}. Entraxe \SI{1.2}{\meter}. Environnement poussiéreux (scierie). Maintenance limitée.
\begin{enumerate}
    \item \textbf{Rapport} $k = 500/1500 = 1/3$.
    \item \textbf{Puissance \SI{4}{\kilo\watt}} : courroie plate limite ($\sim\SI{5}{\kilo\watt}$), risqué. Courroie V confortable.
    \item \textbf{Poussière} : chaîne à éviter (usure accélérée, lubrification piège poussière). Engrenages impossibles (entraxe).
    \item \textbf{Choix :} \textbf{Courroie trapézoïdale} (section SPB ou B selon norme). Poulie moteur $D_1 \approx \SI{100}{\milli\meter}$, poulie menée $D_2 \approx \SI{300}{\milli\meter}$. Galet tendeur sur brin relâché. Protection par carter grillagé.
\end{enumerate}
\end{exemplebox}

\courssub{Exemple chiffré complet : transmission moto 125 cm³}

\begin{exoresolu}[Transmission secondaire d'une moto 125 cm³]
\textbf{Données :} Pignon moteur $Z_1 = 15$ dents. Couronne roue arrière $Z_2 = 45$ dents. Régime moteur max $N_1 = \SI{8000}{tr/min}$. Roue arrière : diamètre $D_{\text{roue}} = 17$ pouces ($= \SI{431.8}{\milli\meter}$). Pas de chaîne $p = \SI{15.875}{\milli\meter}$ (5/8"). Rendement chaîne $\eta = 0,98$.

\textbf{Questions :}
\begin{enumerate}
    \item Calculer le rapport de transmission $k$ et le régime de la roue $N_2$.
    \item Calculer la vitesse théorique $v$ (km/h) à \SI{8000}{tr/min}.
    \item Si la chaîne s'allonge de 1\,\% (usure), que se passe-t-il sur le rapport $k$ ?
    \item Pourquoi un galet tendeur est-il indispensable ?
\end{enumerate}

\textbf{Solution détaillée :}

\textbf{1. Rapport et régime roue}
\[
k = \frac{Z_1}{Z_2} = \frac{15}{45} = \frac{1}{3}
\]
\[
N_2 = k \cdot N_1 = \frac{1}{3} \times 8000 = \SI{2667}{tr/min}
\]
\emph{(Pas de glissement : $k$ exact car chaîne = liaison forme).}

\textbf{2. Vitesse théorique}
Périmètre roue $P = \pi \cdot D_{\text{roue}} = \pi \times 0,4318 = \SI{1,356}{\meter}$.
Distance parcourue par minute $= N_2 \times P = 2667 \times 1,356 = \SI{3616}{\meter\per\minute}$.
\[
v = \frac{3616}{1000} \times 60 = \SI{217}{km/h}
\]
\emph{Attention : c'est la vitesse \textbf{théorique} (sans glissement pneu/route, sans résistance air/roulement). En réel, une 125 cm³ atteint $\sim \SIrange{100}{110}{km/h}$. La différence vient du \textbf{glissement pneu/route} (patinage) et des pertes aérodynamiques.}

\textbf{3. Effet de l'allongement 1\,\%}
Le pas effectif devient $p' = 1,01 \times p$. La chaîne « monte » sur les dents du pignon mené $\Rightarrow$ risque de \textbf{saut de dent} sous couple. Le rapport $k$ reste \textbf{théoriquement exact} (liaison forme) tant que la chaîne ne saute pas, mais l'engrènement se dégrade (bruit, usure accélérée, rupture possible). \textbf{D'où l'importance du tendeur.}

\textbf{4. Rôle du galet tendeur}
Compense l'allongement (maintient tension), augmente l'angle d'enroulement sur le pignon mené (évite le saut de dent), absorbe les à-coups de couple (accélération/freinage moteur). Sur les motos modernes, le tendeur est souvent \textbf{automatique à ressort} (ou hydraulique) + galet guide fixe.
\end{exoresolu}

\courssub{Synthèse : comment choisir ?}

\begin{methode}[Arbre de décision simplifié pour le choix d'une transmission flexible]
\begin{enumerate}
    \item \textbf{Précision $k$ exigée ?} (horloge, distribution moteur, robotique) $\rightarrow$ \textbf{Courroie dentée} ou \textbf{Chaîne}.
    \item \textbf{Puissance $> \SI{20}{\kilo\watt}$ ?} $\rightarrow$ \textbf{Chaîne} (ou courroies V multiples).
    \item \textbf{Entraxe $> \SI{3}{\meter}$ ?} $\rightarrow$ \textbf{Courroie plate} (ou V si puissance modérée).
    \item \textbf{Environnement sale/poussièreux, maintenance difficile ?} $\rightarrow$ \textbf{Courroie V} (pas de lubrification, tolère légère saleté).
    \item \textbf{Silence absolu requis ?} $\rightarrow$ \textbf{Courroie dentée} ou \textbf{Courroie plate}.
    \item \textbf{Budget serré, puissance faible, entraxe moyen ?} $\rightarrow$ \textbf{Courroie V} (standard industriel).
\end{enumerate}
\end{methode}

\begin{exercice}[Application : ventilateur de cuisine]
Un ventilateur de cuisine (hôtel) est entraîné par un moteur électrique $N_1 = \SI{1450}{tr/min}$ via une courroie trapézoïdale. Poulie moteur $D_1 = \SI{80}{\milli\meter}$, poulie ventilateur $D_2 = \SI{240}{\milli\meter}$. Entraxe $E = \SI{600}{\milli\meter}$. Glissement estimé $g = 2\,\%$.
\begin{enumerate}
    \item Calculer le rapport théorique $k_{\text{th}}$ et le régime théorique $N_{2,\text{th}}$.
    \item Calculer le régime réel $N_{2,\text{réel}}$ avec glissement.
    \item Quel type de galet tendeur préconises-tu ? Où le placer ?
    \item Si on veut supprimer le glissement pour une régulation précise de débit, quelle courroie choisir ? Quelles modifications sur les poulies ?
\end{enumerate}
\end{exercice}

\begin{corrige}[Exercice : ventilateur de cuisine]
\begin{enumerate}
    \item $k_{\text{th}} = D_1/D_2 = 80/240 = 1/3$. $N_{2,\text{th}} = 1450/3 = \SI{483}{tr/min}$.
    \item $N_{2,\text{réel}} = N_{2,\text{th}} \times (1 - g) = 483 \times 0,98 = \SI{473}{tr/min}$.
    \item Galet lisse sur \textbf{brin relâché}, extérieur, à ressort (auto-tendeur). Augmente l'enroulement sur les deux poulies.
    \item \textbf{Courroie dentée} + \textbf{poulies dentées} de même pas. $k$ devient exact. Attention : entraxe doit correspondre à une longueur standard de courroie dentée (ou utiliser un tendeur denté pour ajuster).
\end{enumerate}
\end{corrige}

\begin{savaistu}[La courroie plate... en cuir !]
Avant l'invention du caoutchouc vulcanisé (Goodyear, 1839) et des fibres synthétiques, les \textbf{courroies plates étaient en cuir} (cuir de bœuf, tanné, graissé). Elles entraînaient les machines-outils, métiers à tisser, scieries, depuis la machine à vapeur centrale par des arbres de transmission aériens (les fameux « poulies et courroies » des usines du XIXe siècle). Au Cameroun, on en voit encore sur d'anciennes décortiqueuses à riz ou presses à huile artisanales. Le cuir glisse beaucoup ($g \approx 5-10\,\%$), mais il est increvable, se répare par couture, et ne craint pas l'humidité comme le caoutchouc naturel non vulcanisé.
\end{savaistu}

\begin{attention}[Pièges BEPC à éviter absolument]
\begin{itemize}
    \item Confondre \textbf{glissement} (élastique, normal, $\sim 1-3\,\%$) et \textbf{patinage} (rupture d'adhérence, anormal, arrêt machine).
    \item Oublier que $k = D_1/D_2$ n'est \textbf{valable qu'en l'absence de glissement} (courroie dentée, chaîne, engrenages). Pour courroies plate/V : $k_{\text{réel}} = \frac{D_1}{D_2}(1-g)$.
    \item Placer le tendeur du \textbf{mauvais côté} (brin tendu pour courroie, brin relâché pour chaîne).
    \item Négliger l'\textbf{alignement} des pignons (chaîne) ou des poulies (courroies V/dentées) : cause n°1 de casse prématurée.
    \item Graisser une courroie (V, plate, dentée) : \textbf{perte d'adhérence garantie}.
    \item Confondre \textbf{pas de chaîne} (distance axes goupilles) et \textbf{pas de dent de pignon} (identiques par construction).
\end{itemize}
\end{attention}

\coursec{Systèmes particuliers, rendement et sécurité}

Dans la section précédente, nous avons étudié les transmissions par \cle{liaison flexible} : la courroie et la chaîne permettent de transmettre un mouvement de rotation entre deux arbres \emph{parallèles et éloignés}. Mais que faire lorsque les arbres sont \emph{perpendiculaires}, lorsqu'il faut transformer une rotation en \emph{translation} (portail coulissant), ou lorsque l'angle entre les arbres \emph{varie} pendant le fonctionnement (arbre de transmission d'une voiture) ? Et surtout : une machine réelle restitue-t-elle toute l'énergie qu'on lui fournit ? C'est l'objet de cette dernière section : découvrir trois systèmes particuliers (vis sans fin, crémaillère, joint de Cardan), comprendre la notion de \cle{rendement}, et apprendre les règles de \cle{sécurité} qui sauvent des vies autour des machines.

\courssub{La vis sans fin et la roue hélicoïdale}

\textbf{Situation-problème.} Le portail électrique de la maison s'ouvre lentement grâce à un petit moteur. Une fois fermé, impossible de le pousser à la main, même en forçant ! Pourtant, aucun verrou n'est visible. Quel mécanisme bloque le portail ?

\textbf{Observation et description.} Le système \cle{vis sans fin -- roue hélicoïdale} associe :
\begin{itemize}
\item une \cle{vis sans fin} : une tige cylindrique portant un filet hélicoïdal (comme une grosse vis), placée sur l'arbre \emph{menant} ;
\item une \cle{roue hélicoïdale} : une roue dentée dont les dents sont inclinées et creusées pour épouser le filet de la vis, placée sur l'arbre \emph{mené}.
\end{itemize}
Les deux axes sont \cle{croisés à 90\textdegree} : ils ne sont ni parallèles ni concourants, ils sont perpendiculaires sans se couper. Le nombre de filets de la vis (souvent 1, parfois 2 ou 3) se note $n$ ; le nombre de dents de la roue se note $Z$.

\begin{popfigure}
\begin{center}
\begin{tikzpicture}[scale=1.0]
% Roue hélicoïdale (vue de face)
\draw[fill=popblueL,draw=popblue,thick] (0,0) circle (1.5);
\draw[fill=white,draw=popblue,thick] (0,0) circle (0.25);
\foreach \a in {0,45,...,315}{
  \draw[popblue,thick] (\a:1.5) -- (\a:1.75);
}
\foreach \a in {22.5,67.5,...,337.5}{
  \draw[popblue!60] (\a:1.5) -- (\a:1.68);
}
\node[popblue] at (0,-2.1) {\textbf{Roue hélicoïdale (menée)}};
% axe de la roue
\draw[popdark,thick,->] (0,0) -- (0,0.9) node[midway,right,font=\small]{axe vertical};
% Vis sans fin (cylindre horizontal avec filet)
\draw[fill=poporangeL,draw=poporange,thick] (-2.6,1.9) rectangle (2.6,2.5);
\foreach \x in {-2.4,-2.0,...,2.4}{
  \draw[poporange,thick] (\x,1.9) -- (\x+0.25,2.5);
}
\draw[popdark,thick] (-3.4,2.2) -- (3.4,2.2);
\draw[popdark,thick,->] (3.4,2.2) -- (3.9,2.2) node[right,font=\small]{axe horizontal};
\node[poporange] at (0,2.9) {\textbf{Vis sans fin (menante)}};
% flèches de rotation
\draw[popgreen,thick,->] (2.0,1.55) arc (-60:60:0.35);
\node[popgreen,font=\small] at (2.6,1.55) {$\omega_e$};
\draw[popgreen,thick,->] (-1.9,0.4) arc (150:270:0.4);
\node[popgreen,font=\small] at (-2.4,-0.2) {$\omega_s$};
% angle 90°
\node[popdark,font=\small] at (1.6,1.15) {axes croisés à $90°$};
\end{tikzpicture}
\end{center}
\legende{Schéma d'un réducteur à vis sans fin : la vis (orange), entraînée par le moteur, engrène avec la roue hélicoïdale (bleu). Les deux axes sont croisés à 90\textdegree. Un tour de vis (à un filet) fait avancer la roue d'une seule dent.}
\end{popfigure}

\textbf{Fonctionnement et rapport de transmission.} À chaque tour complet de la vis, le filet pousse la roue de $n$ dents. Si la roue possède $Z$ dents, il faut $Z/n$ tours de vis pour un tour de roue. Le rapport de transmission est donc :
\[ k = \frac{\omega_s}{\omega_e} = \frac{n}{Z} \]
Pour une vis à un filet ($n=1$), $k = 1/Z$. Avec des roues de 20 à 300 dents, on obtient des rapports \emph{très petits} : $k = 1/20$ à $1/300$. C'est le système le plus \cle{compact} pour obtenir une très grande réduction de vitesse dans un faible volume.

\begin{propriete}[Irréversibilité de la vis sans fin]
Si l'angle d'inclinaison de l'hélice de la vis est inférieur à l'angle de frottement entre les matériaux en contact, la transmission est \cle{irréversible} : la vis peut entraîner la roue, mais la roue \emph{ne peut pas} entraîner la vis. Le système se comporte comme un \textbf{frein naturel}.
\begin{itemize}
\item \textbf{Conséquence pratique :} le portail électrique ne peut pas être rouvert à la main ; la charge d'un treuil ne redescend pas toute seule quand on arrête le moteur.
\item \textbf{Contrepartie :} les frottements importants dégradent le rendement, souvent faible : $\eta \approx 0{,}4$ à $0{,}7$.
\end{itemize}
\end{propriete}

\begin{exemplebox}[Applications de la vis sans fin]
Vérins électriques, motorisations de portails, treuils de levage, anciennes directions de voiture (vis et galet), tables tournantes de machines-outils. Partout où l'on veut une grande réduction de vitesse et un blocage automatique en position.
\end{exemplebox}

\courssub{Le pignon et la crémaillère : de la rotation à la translation}

\textbf{Situation-problème.} Quand le chauffeur d'une voiture tourne le volant, les roues avant pivotent. Mais dans un portail \emph{coulissant}, le moteur tourne et le portail se déplace \emph{en ligne droite}. Comment une rotation peut-elle produire un déplacement rectiligne ?

\begin{definition}[Crémaillère]
Une \cle{crémaillère} est une barre rectiligne dentée. On peut la considérer comme une roue à dents droites de \textbf{rayon infini} : sa denture est développée « à plat ». Associée à un petit pignon qui engrène sur elle, elle transforme :
\begin{itemize}
\item une \textbf{rotation en translation} (le pignon moteur fait avancer la crémaillère) ;
\item ou une \textbf{translation en rotation} (on pousse la crémaillère, le pignon tourne).
\end{itemize}
C'est une transformation de mouvement \emph{réversible} : Rotation $\leftrightarrow$ Translation.
\end{definition}

\begin{popfigure}
\begin{center}
\begin{tikzpicture}[scale=1.0]
% Crémaillère
\draw[fill=popgreenL,draw=popgreen,thick] (-3.5,0) rectangle (3.5,0.5);
\foreach \x in {-3.3,-2.9,...,3.3}{
  \draw[fill=popgreenL,draw=popgreen,thick] (\x,0.5) -- (\x+0.1,0.8) -- (\x+0.3,0.8) -- (\x+0.4,0.5) -- cycle;
}
\node[popgreen] at (0,-0.4) {\textbf{Crémaillère (barre dentée)}};
% Pignon
\draw[fill=popblueL,draw=popblue,thick] (0,1.75) circle (0.75);
\foreach \a in {0,30,...,330}{
  \draw[popblue,thick] (\a:0.75) -- (\a:0.95);
}
\draw[fill=white,draw=popblue,thick] (0,1.75) circle (0.15);
\node[popblue] at (0,2.9) {\textbf{Pignon}};
% flèches
\draw[poporange,thick,->] (-1.35,1.75) arc (180:300:0.6);
\node[poporange,font=\small] at (-1.9,1.3) {$\omega$};
\draw[poporange,thick,->] (1.2,0.25) -- (2.6,0.25) node[right,font=\small]{$v$ (translation)};
% détail denture
\draw[popmuted,thick] (4.1,0.4) -- (4.1,2.4);
\node[popmuted,font=\small,align=left] at (5.6,1.4) {Détail : les dents\\du pignon pénètrent\\entre les dents\\de la crémaillère.};
\end{tikzpicture}
\end{center}
\legende{Ensemble pignon--crémaillère : la rotation du pignon (bleu) provoque la translation de la barre dentée (vert). Si le pignon a un rayon primitif $r$, un tour de pignon déplace la crémaillère de $2\pi r$.}
\end{popfigure}

\textbf{Relation vitesse de rotation -- vitesse de translation.} Si le pignon de rayon primitif $r$ tourne à la vitesse angulaire $\omega$ (en rad/s), la crémaillère se déplace à la vitesse :
\[ v = r \times \omega \]
Un tour complet du pignon fait avancer la crémaillère d'une distance égale au périmètre primitif $2\pi r$.

\begin{exemplebox}[Applications du pignon-crémaillère]
\begin{itemize}
\item \textbf{Direction à crémaillère} des voitures : le volant fait tourner un pignon qui déplace latéralement la crémaillère, laquelle pousse les roues.
\item \textbf{Portails coulissants} motorisés : le pignon du moteur engrène sur une crémaillère fixée au portail.
\item \textbf{Trains à crémaillère} (lignes de montagne) : la crémaillère posée entre les rails empêche le train de patiner dans les fortes pentes.
\end{itemize}
\end{exemplebox}

\courssub{Le joint de Cardan : transmettre entre des arbres à angle variable}

\textbf{Situation-problème.} Sur un camion ou un tracteur, le moteur est fixe, mais l'essieu arrière monte et descend avec les cahots de la route. Comment transmettre la rotation entre deux arbres dont l'angle \emph{change sans cesse} ?

\begin{definition}[Joint de Cardan]
Le \cle{joint de Cardan} (ou joint universel) est une transmission entre deux arbres \emph{concourants} dont l'angle $\alpha$ peut varier en fonctionnement. Il est constitué de deux \textbf{fourchettes} (une par arbre) reliées par une \textbf{croix} pivotante appelée croisillon.
\end{definition}

\begin{popfigure}
\begin{center}
\begin{tikzpicture}[scale=1.0]
% arbre entrée (horizontal)
\draw[popdark,thick] (-3.2,0) -- (-1.1,0);
\node[popdark,font=\small] at (-2.6,0.35) {arbre menant};
% fourchette entrée
\draw[fill=popblueL,draw=popblue,thick] (-1.1,-0.55) rectangle (-0.35,-0.15);
\draw[fill=popblueL,draw=popblue,thick] (-1.1,0.15) rectangle (-0.35,0.55);
\node[popblue,font=\small] at (-1.5,0.95) {fourchette};
% croisillon
\draw[poporange,very thick] (0,-0.85) -- (0,0.85);
\draw[poporange,very thick] (-0.55,0) -- (0.85,0.35);
\draw[fill=poporange,draw=poporange] (0.15,0.09) circle (0.14);
\node[poporange,font=\small] at (0.15,-1.15) {croisillon};
% fourchette sortie (inclinée)
\draw[fill=popgreenL,draw=popgreen,thick,rotate around={25:(0.85,0.35)}] (0.85,0.1) rectangle (1.6,0.6);
\draw[fill=popgreenL,draw=popgreen,thick,rotate around={25:(0.85,0.35)}] (0.85,-0.35) rectangle (1.6,0.15);
% arbre sortie incliné
\draw[popdark,thick] (1.55,0.72) -- (3.3,1.55);
\node[popdark,font=\small] at (3.0,1.95) {arbre mené};
% angle alpha
\draw[popmuted,thick] (1.3,0) arc (0:25:1.3);
\node[popmuted,font=\small] at (1.75,0.28) {$\alpha$};
\draw[popmuted,dashed] (0.2,0) -- (2.6,0);
\end{tikzpicture}
\end{center}
\legende{Joint de Cardan simple : deux fourchettes (bleu et vert) articulées autour d'un croisillon (orange) transmettent la rotation entre deux arbres formant un angle $\alpha$.}
\end{popfigure}

\textbf{Propriété essentielle.} En un tour complet, l'arbre mené fait bien un tour : le rapport \emph{moyen} vaut 1. Mais si l'angle $\alpha$ n'est pas nul, la vitesse de sortie \emph{n'est pas constante} : elle oscille deux fois par tour entre un maximum et un minimum :
\[ \omega_e \cos\alpha \;\le\; \omega_s \;\le\; \frac{\omega_e}{\cos\alpha} \]

\begin{attention}[Le Cardan simple n'est PAS homocinétique]
Avec un joint simple et un angle $\alpha > 0$, la vitesse de sortie est \cle{pulsatoire} : elle présente \textbf{2 pics et 2 creux par tour}. Au-delà de 15\textdegree{} à 20\textdegree{} d'angle, cela provoque des \textbf{vibrations} gênantes et une usure rapide. Pour obtenir une vitesse de sortie rigoureusement constante (transmission \emph{homocinétique}), il faut un \textbf{double joint de Cardan} (deux joints montés tête-bêche dont les défauts se compensent) ou un joint homocinétique spécial. \emph{Piège classique au BEPC : ne jamais écrire qu'un cardan simple transmet à vitesse constante !}
\end{attention}

\begin{exemplebox}[Applications du joint de Cardan]
Arbre de transmission des voitures à propulsion, prise de force des tracteurs (qui entraîne la remorque ou la pulpeuse même en virage et sur terrain bosselé), volant de direction (colonne coudée), perceuses à renvoi d'angle.
\end{exemplebox}

\courssub{Le rendement d'une transmission}

\textbf{Situation-problème.} Un groupe électrogène fournit une puissance mécanique de \SI{1000}{W} à un alternateur... mais à l'arbre de sortie d'un réducteur, on ne récupère jamais toute la puissance d'entrée. Où passe l'énergie perdue ?

\textbf{Démarche d'investigation.} Faisons tourner à la main un réducteur à vis sans fin : après quelques minutes, le carter est \emph{chaud}. L'énergie mécanique perdue a été transformée en \emph{chaleur} par les frottements. Hypothèse : toute transmission réelle perd de l'énergie. Vérifions par la mesure : on mesure la puissance d'entrée $P_e$ (moteur) et la puissance de sortie $P_s$ (charge) ; on constate toujours $P_s < P_e$. Conclusion : une transmission ne fait que \emph{transférer} l'énergie, avec des pertes.

\begin{definition}[Rendement]
Le \cle{rendement} d'une transmission, noté $\eta$ (lettre grecque « êta »), est le rapport de la puissance utile de sortie à la puissance fournie en entrée :
\[ \eta = \frac{P_{\text{sortie}}}{P_{\text{entrée}}} \qquad \text{avec toujours } \eta < 1 \]
Il s'exprime sans unité, ou en pourcentage. Les pertes proviennent des \textbf{frottements} (dents d'engrenages, roulements), du \textbf{glissement} (courroies) et de la \textbf{ventilation et du barbotage de l'huile}.
\end{definition}

\begin{aretenir}[Ordres de grandeur des rendements]
\begin{center}
\begin{tabularx}{\linewidth}{|l|X|}\hline
\rowcolor{popblueL} \textbf{Transmission} & \textbf{Rendement typique} \\ \hline
Engrenages bien lubrifiés (un étage) & $0{,}95$ à $0{,}98$ \\ \hline
Courroies et chaînes & $0{,}90$ à $0{,}97$ \\ \hline
Vis sans fin & $0{,}4$ à $0{,}8$ (frottements élevés) \\ \hline
\end{tabularx}
\end{center}
Pour plusieurs étages en cascade, les rendements se \emph{multiplient} : $\eta_{\text{total}} = \eta_1 \times \eta_2 \times \dots$
\end{aretenir}

\begin{methode}[Calculer le couple de sortie d'une transmission]
La puissance se conserve aux pertes près : $P_s = \eta \times P_e$. Comme $P = C \times \omega$ (couple $\times$ vitesse angulaire), on obtient :
\[ C_{\text{sortie}} = \eta \times \frac{\omega_e}{\omega_s} \times C_{\text{entrée}} = \frac{\eta}{k} \times C_{\text{entrée}} \]
où $k = \omega_s/\omega_e$ est le rapport de transmission. Étapes :
\begin{enumerate}
\item Calculer le rapport $k = \omega_s/\omega_e$ (ou $N_s/N_e$).
\item Retenir : dans un \textbf{réducteur} ($k < 1$), la vitesse diminue mais le \textbf{couple est multiplié} par $1/k$.
\item Appliquer le rendement : multiplier encore par $\eta$.
\end{enumerate}
\emph{C'est pourquoi un réducteur qui ralentit beaucoup permet de soulever de lourdes charges avec un petit moteur.}
\end{methode}

\begin{exoresolu}[Treuil de levage motorisé]
Un moteur de treuil fournit un couple $C_e = \SI{5}{N.m}$ à $N_e = \SI{1500}{tr/min}$. Il entraîne un réducteur à vis sans fin de rapport $k = 1/50$ et de rendement $\eta = 0{,}6$. Calculer la vitesse et le couple de sortie.
\tcblower
\textbf{Solution.} Vitesse de sortie :
\[ N_s = k \times N_e = \frac{1}{50} \times 1500 = \SI{30}{tr/min} \]
Couple de sortie :
\[ C_s = \frac{\eta}{k} \times C_e = \frac{0{,}6}{1/50} \times 5 = 0{,}6 \times 50 \times 5 = \SI{150}{N.m} \]
\textbf{Conclusion :} la vitesse est divisée par 50, mais le couple est multiplié par 30 (à cause du rendement de 0,6, on ne récupère pas le facteur 50 complet). Le treuil soulève de lourdes charges, lentement, et la charge reste bloquée moteur arrêté grâce à l'irréversibilité.
\end{exoresolu}

\courssub{La sécurité autour des transmissions mécaniques}

Chaque année, des accidents graves se produisent dans les ateliers et les champs : une main happée par une courroie, un vêtement enroulé autour d'un arbre de prise de force. Les transmissions présentent des \cle{zones de pincement} (point d'entrée des dents d'engrenages, contact courroie-poulie, chaîne-pignon) et des \cle{parties tournantes accessibles} qui attrapent doigts, cheveux, pagne ou vêtements amples.

\begin{popfigure}
\begin{center}
\begin{tikzpicture}[scale=0.85]
% Picto 1 : danger pièces tournantes
\draw[popgold,very thick,fill=popgoldL] (0,1.1) -- (-1.15,-0.85) -- (1.15,-0.85) -- cycle;
\draw[popdark,thick] (0,0.35) circle (0.28);
\foreach \a in {90,210,330}{ \draw[popdark,thick] (\a:0.28) -- (\a:0.55); }
\node[font=\scriptsize,align=center] at (0,-1.25) {Danger : pièces\\tournantes};
% Picto 2 : interdiction mains
\draw[popink,very thick,fill=white] (3.4,0.1) circle (0.95);
\draw[popink,very thick] (2.75,-0.55) -- (4.05,0.75);
\draw[popdark,thick,fill=popcream] (3.4,-0.15) ellipse (0.22 and 0.35);
\foreach \dx in {-0.18,-0.06,0.06,0.18}{ \draw[popdark,thick] (3.4+\dx,0.15) -- (3.4+\dx,0.5); }
\node[font=\scriptsize,align=center] at (3.4,-1.25) {Ne pas approcher\\les mains};
% Picto 3 : port protections
\draw[popblue,very thick,fill=popblueL] (6.8,0.1) circle (0.95);
\draw[popdark,thick,fill=white] (6.55,0.25) circle (0.22);
\draw[popdark,thick,fill=white] (7.15,0.25) circle (0.22);
\draw[popdark,thick] (6.77,0.25) -- (6.93,0.25);
\draw[popdark,thick,fill=popcream] (6.85,-0.35) circle (0.3);
\node[font=\scriptsize,align=center] at (6.8,-1.25) {Port obligatoire :\\lunettes, gants};
% Picto 4 : consignation
\draw[popgreen,very thick,fill=popgreenL] (10.2,0.1) circle (0.95);
\draw[popdark,very thick] (10.2,0.35) arc (90:270:0.25) -- cycle;
\draw[popdark,very thick,fill=popgold] (9.95,-0.5) rectangle (10.45,0.1);
\draw[popgold,fill=popdark] (10.2,-0.2) circle (0.07);
\node[font=\scriptsize,align=center] at (10.2,-1.25) {Consignation :\\cadenas sur\\la source d'énergie};
\end{tikzpicture}
\end{center}
\legende{Pictogrammes de sécurité autour des machines tournantes : signalisation du danger, interdiction d'approcher les mains, port obligatoire des équipements de protection, consignation avant intervention.}
\end{popfigure}

\textbf{Les moyens de protection.} Les constructeurs et les utilisateurs doivent mettre en place :
\begin{itemize}
\item des \cle{carters fixes} qui recouvrent engrenages, courroies et chaînes (comme le carter de chaîne d'une moto) ;
\item des \cle{capots verrouillés} : la machine s'arrête automatiquement si le capot est ouvert ;
\item des \cle{arrêts d'urgence} : gros bouton rouge coupant immédiatement l'alimentation.
\end{itemize}

\begin{attention}[Consignes avant tout démontage ou réglage]
Avant toute intervention sur une machine (réglage, graissage, démontage), respecter dans l'ordre :
\begin{enumerate}
\item \textbf{Arrêter} la machine à l'aide de la commande normale.
\item \textbf{Consigner} (Lockout/Tagout) : placer un \textbf{cadenas personnel} sur l'interrupteur ou la source d'énergie, avec une étiquette signalant l'intervention. Personne d'autre ne peut redémarrer.
\item \textbf{Attendre l'arrêt total} des parties tournantes (une roue peut tourner encore plusieurs minutes par inertie).
\item \textbf{Vérifier l'absence de tension} (VAT) pour les machines électriques.
\end{enumerate}
Ne jamais porter de vêtements amples, gants près d'un arbre en rotation, ni laisser les cheveux longs détachés.
\end{attention}

\begin{savaistu}[La direction assistée électrique]
Sur les voitures modernes, la direction assistée \emph{hydraulique} (une pompe entraînée par le moteur) est remplacée par la direction assistée \emph{électrique} : un petit moteur électrique aide le conducteur seulement quand il tourne le volant, ce qui économise du carburant. Mais attention : le lien mécanique final entre le volant et les roues reste un bon vieux \textbf{pignon-crémaillère} ! Même en cas de panne totale d'assistance, le conducteur garde le contrôle du véhicule. Au Cameroun, les mécaniciens de quartier savent bien que la crémaillère usée provoque un « jeu » au volant : les roues répondent avec retard.
\end{savaistu}

\begin{exercice}[Pour s'entraîner]
Un portail coulissant est motorisé par un pignon de rayon primitif $r = \SI{3}{cm}$ tournant à $N = \SI{20}{tr/min}$.
\begin{enumerate}
\item Calculer la vitesse angulaire $\omega$ du pignon en rad/s.
\item En déduire la vitesse de déplacement $v$ du portail en m/s.
\item Le réducteur à vis sans fin a un rendement $\eta = 0{,}5$. Le moteur fournit $\SI{200}{W}$. Quelle puissance utile arrive au pignon ?
\end{enumerate}
\end{exercice}

\begin{corrige}[Exercice]
\begin{enumerate}
\item $\omega = \dfrac{2\pi N}{60} = \dfrac{2\pi \times 20}{60} \approx \SI{2,09}{rad/s}$.
\item $v = r \times \omega = 0{,}03 \times 2{,}09 \approx \SI{0,063}{m/s}$, soit environ \SI{6,3}{cm/s} : le portail de 4 m met un peu plus d'une minute à s'ouvrir.
\item $P_{\text{utile}} = \eta \times P_e = 0{,}5 \times 200 = \SI{100}{W}$. La moitié de la puissance est perdue en chaleur dans la vis sans fin.
\end{enumerate}
\end{corrige}

\begin{aretenir}[L'essentiel de la section]
\begin{itemize}
\item La \cle{vis sans fin} : axes croisés à 90\textdegree, très grande réduction ($k = n/Z$), souvent \textbf{irréversible} (frein naturel), rendement faible.
\item Le \cle{pignon-crémaillère} : transforme rotation $\leftrightarrow$ translation, avec $v = r\,\omega$.
\item Le \cle{joint de Cardan} : arbres concourants à angle variable ; simple = non homocinétique, double = homocinétique.
\item Le \cle{rendement} $\eta = P_s/P_e < 1$ mesure les pertes ; le couple de sortie vaut $C_s = \dfrac{\eta}{k}\,C_e$.
\item La \cle{sécurité} impose carters, arrêts d'urgence et \textbf{consignation} avant toute intervention.
\end{itemize}
\end{aretenir}

\newpage

\coursec{Activité d'intégration}

\courssub{Objectifs de l'activité}

À la fin de cette activité d'intégration, tu seras capable de :
\begin{itemize}
\item calculer une \cle{vitesse angulaire} à partir d'une vitesse linéaire et du rayon d'une roue ;
\item déterminer un \cle{rapport de transmission} global nécessaire et le décomposer en étages ;
\item choisir et justifier un \cle{système de transmission} (chaîne, courroie, engrenages) adapté à un contexte réel ;
\item dessiner un \cle{schéma cinématique} simple avec les sens de rotation ;
\item identifier les \cle{zones de danger} d'une machine et proposer des protections ;
\item estimer un \cle{couple} transmis en tenant compte du \cle{rendement} et vérifier une condition de traction en pente.
\end{itemize}

\courssub{Situation}

\textbf{Le défi du tricycle « Kéké » de M. Njoya.}

M. Njoya est transporteur de marchandises au marché Mokolo à Yaoundé. Son tricycle à pédales, chargé de sacs d'oignons et de cartons, peine dans les côtes du quartier. Il décide de le motoriser et te sollicite, toi l'élève de 3ème passionné de mécanique, pour l'aider à concevoir la transmission.

\textbf{Données du problème :}
\begin{itemize}
\item Charge utile maximale : \SI{200}{kg} ; masse totale roulante estimée (tricycle + moteur + conducteur + charge) : $m = \SI{250}{kg}$.
\item Moteur thermique 4 temps essence (type moteur de groupe électrogène) : puissance $P = \SI{3,5}{kW}$ à \SI{3600}{tr/min} ; couple maximal $C_m = \SI{9}{N.m}$. On prendra pour les calculs le régime nominal $N_m = \SI{3600}{tr/min}$ et on admettra que le couple de \SI{9}{N.m} reste disponible à ce régime (hypothèse simplificatrice, prudente car le couple réel y est voisin).
\item Roue arrière motrice : diamètre $D = \SI{500}{mm}$, donc rayon $R = \SI{0,25}{m}$.
\item Vitesse souhaitée : $v = \SI{25}{km/h}$.
\item Pente maximale à gravir : 15 \% (soit $\tan\alpha = 0,15$, d'où $\sin\alpha \approx 0,148$).
\item Rendement global estimé de la transmission : $\eta = 0,85$.
\item Intensité de la pesanteur : $g = \SI{10}{N/kg}$.
\end{itemize}

\textbf{Documents mis à disposition :} fiche technique du moteur ; photos de trois transmissions existantes (chaîne de moto-taxi « benskin », courroie de scooter, engrenages d'un réducteur de moulin à maïs) ; extrait de catalogue (pignons de 12 à 15 dents, couronnes de 48 à 60 dents, chaînes pas 428, poulies de 60 à 200 mm, courroies trapézoïdales) ; extraits de consignes de sécurité des machines.

\textbf{Tâche finale :} rédiger une fiche technique argumentée d'une page, accompagnée du schéma cinématique annoté, que M. Njoya pourra présenter à son mécanicien.

\courssub{Consignes}

\begin{enumerate}
\item Calcule la vitesse angulaire $\omega_r$ de la roue arrière lorsque le tricycle roule à \SI{25}{km/h}. Exprime le résultat en rad/s puis en tr/min ($N_r$).
\item Détermine le rapport de transmission global $k = \dfrac{\omega_r}{\omega_m}$ nécessaire entre le moteur et la roue. Ce rapport est-il réalisable en un seul étage de chaîne du catalogue ? Justifie.
\item Propose une solution technique complète : type(s) de transmission, nombre de dents des pignons et couronnes (ou diamètres des poulies). Justifie ton choix en comparant chaîne, courroie et engrenages selon quatre critères : coût, maintenance, rendement, sécurité.
\item Dessine le schéma cinématique complet (moteur $\rightarrow$ réducteur $\rightarrow$ roue) en indiquant par des flèches le sens de rotation de chaque arbre.
\item Identifie trois zones de danger sur le tricycle motorisé et propose pour chacune une protection adaptée (carter, capot, écran).
\item Estime le couple disponible à la roue $C_r$ avec $\eta = 0,85$, puis la force de traction $F = \dfrac{C_r}{R}$. Vérifie que cette force suffit à gravir la pente de 15 \% en la comparant à la composante du poids le long de la pente, $F_p = m \cdot g \cdot \sin\alpha$.
\end{enumerate}

\courssub{Ressources à mobiliser}

\begin{aretenir}[Formules utiles]
\begin{itemize}
\item Vitesse linéaire et vitesse angulaire : $v = R \cdot \omega$, avec $v$ en m/s, $R$ en m et $\omega$ en rad/s.
\item Conversion : $N\ (\text{tr/min}) = \omega\ (\text{rad/s}) \times \dfrac{60}{2\pi}$.
\item Rapport de transmission : $k = \dfrac{\omega_s}{\omega_e} = \dfrac{N_s}{N_e}$ ; pour un engrenage ou une chaîne, $k = \dfrac{Z_e}{Z_s}$ ($Z$ : nombre de dents ; $e$ : entrée, $s$ : sortie).
\item Rapports en cascade : $k_{\text{global}} = k_1 \times k_2$.
\item Couple transmis : $C_s = C_e \times \dfrac{1}{k} \times \eta$ (le couple augmente quand la vitesse diminue).
\item Force de traction à la roue : $F = \dfrac{C_r}{R}$.
\item Force nécessaire en pente : $F_p = m \cdot g \cdot \sin\alpha$.
\end{itemize}
\end{aretenir}

\courssub{Démarche et corrigé commenté}

\begin{exoresolu}[Le défi du tricycle « Kéké » — corrigé détaillé]
Énoncé : concevoir la transmission du tricycle de M. Njoya (données ci-dessus) et répondre aux six consignes.
\tcblower
\textbf{Étape 1 : vitesse angulaire de la roue.}

On convertit d'abord la vitesse dans l'unité du système international :
\[
v = \SI{25}{km/h} = \frac{25}{3,6} \approx \SI{6,94}{m/s}
\]

La roue roule sans glisser : un point de sa circonférence avance à la vitesse $v$ du véhicule. D'où :
\[
\omega_r = \frac{v}{R} = \frac{6,94}{0,25} \approx \SI{27,8}{rad/s}
\]

Conversion en tr/min :
\[
N_r = 27,8 \times \frac{60}{2\pi} \approx \SI{265}{tr/min}
\]

\textbf{Étape 2 : rapport de transmission global.}

\[
k = \frac{N_r}{N_m} = \frac{265}{3600} \approx 0,074 \approx \frac{1}{13,6}
\]

Il faut donc réduire la vitesse environ 13,6 fois. Une chaîne seule avec le plus petit pignon du catalogue (12 dents) exigerait une couronne de $12 \times 13,6 \approx 163$ dents : énorme, coûteuse, encombrante et hors catalogue (maximum 60 dents). Il faut donc \textbf{deux étages} de réduction.

\textbf{Étape 3 : solution technique proposée.}

On choisit :
\begin{itemize}
\item \textbf{Étage 1 — réducteur à engrenages} (carter fermé, type réducteur de moulin) : pignon menant $Z_1 = 15$ dents, roue menée $Z_2 = 60$ dents, soit $k_1 = \dfrac{Z_1}{Z_2} = \dfrac{15}{60} = \dfrac{1}{4}$.
\item \textbf{Étage 2 — chaîne} (pas 428, comme les motos) : pignon $Z_3 = 14$ dents, couronne $Z_4 = 48$ dents, soit $k_2 = \dfrac{Z_3}{Z_4} = \dfrac{14}{48} \approx \dfrac{1}{3,43}$.
\end{itemize}

Vérification :
\[
k_{\text{global}} = k_1 \times k_2 = \frac{1}{4} \times \frac{1}{3,43} \approx \frac{1}{13,7}
\]
La roue tourne alors à :
\[
N_r = \frac{3600}{13,7} \approx \SI{263}{tr/min}
\]
soit $\omega_r \approx \SI{27,5}{rad/s}$ et $v = R\,\omega_r \approx 0,25 \times 27,5 \approx \SI{6,9}{m/s} \approx \SI{24,8}{km/h}$ : l'objectif de \SI{25}{km/h} est atteint.

\textbf{Justification du choix :} les engrenages en carter sont placés près du moteur : à cet endroit le couple est encore faible (il n'a pas encore été multiplié), l'encombrement est réduit, le rendement est excellent et le graissage est permanent. La chaîne est préférée à la courroie pour l'étage final : elle ne patine pas sous la charge de \SI{200}{kg}, supporte la poussière et les nids-de-poule, et les mécaniciens de motos du quartier savent l'entretenir et trouvent les pièces facilement. La courroie, moins chère et plus silencieuse mais glissante sous forte charge et fragile à la chaleur, est écartée.

\textbf{Étape 4 : schéma cinématique.}

\begin{popfigure}
\begin{center}
\begin{tikzpicture}[scale=1, every node/.style={font=\small}]
% Moteur
\draw[fill=poporangeL, thick] (0,0) rectangle (1.6,1.2);
\node at (0.8,0.6) {Moteur};
\node[below, font=\footnotesize] at (0.8,0) {\SI{3600}{tr/min}};
% Engrenages
\draw[fill=popblueL, thick] (2.4,0.35) circle (0.25);
\node[font=\footnotesize] at (2.4,0.35) {$Z_1$};
\draw[fill=popblue, thick] (2.4,1.15) circle (0.5);
\node[font=\footnotesize, white] at (2.4,1.15) {$Z_2$};
\draw[thick] (1.6,0.35) -- (2.15,0.35);
% Chaine
\draw[fill=popgreenL, thick] (4.2,0.9) circle (0.28);
\node[font=\footnotesize] at (4.2,0.9) {$Z_3$};
\draw[fill=popgreen, thick] (6.6,0.9) circle (0.75);
\node[font=\footnotesize, white] at (6.6,0.9) {$Z_4$};
\draw[thick] (2.4,1.15) -- (3.92,0.94);
\draw[thick, popdark] (4.2,1.18) -- (6.6,1.65);
\draw[thick, popdark] (4.2,0.62) -- (6.6,0.15);
% Fleches sens
\draw[->, thick, poporange] (0.8,1.35) arc (120:60:0.5);
\draw[->, thick, popblue] (4.2,1.45) arc (120:60:0.5);
\draw[->, thick, popgreen!60!black] (6.6,1.9) arc (120:60:0.6);
\node[below, font=\footnotesize] at (6.6,0.1) {Roue \SI{263}{tr/min}};
\end{tikzpicture}
\end{center}
\legende{Schéma cinématique : moteur $\rightarrow$ engrenages $Z_1/Z_2$ $\rightarrow$ chaîne $Z_3/Z_4$ $\rightarrow$ roue. Les flèches indiquent les sens de rotation : la chaîne conserve le sens de rotation, l'engrenage à contact extérieur l'inverse.}
\end{popfigure}

\textbf{Étape 5 : zones de danger et protections.}
\begin{enumerate}
\item \textbf{La chaîne et la couronne} : risque de happement des doigts, des vêtements ou des lacets $\rightarrow$ \textbf{carter de chaîne} métallique fixe, comme sur les motos.
\item \textbf{Les engrenages du réducteur} : écrasement des doigts, projection d'huile chaude $\rightarrow$ \textbf{capot fermé} boulonné, ouvrable seulement avec un outil.
\item \textbf{Le moteur (pot d'échappement brûlant, pièces en rotation)} : brûlure et contact accidentel $\rightarrow$ \textbf{écran thermique} et grille de protection.
\end{enumerate}
Consigne associée : arrêter le moteur et attendre son refroidissement avant toute intervention.

\textbf{Étape 6 : couple à la roue et vérification en pente.}

Le réducteur divise la vitesse par 13,7 : il multiplie le couple par 13,7, moins les pertes (rendement $\eta = 0,85$) :
\[
C_r = C_m \times \frac{1}{k} \times \eta = 9 \times 13,7 \times 0,85 \approx \SI{105}{N.m}
\]

Force de traction disponible au contact de la roue et du sol :
\[
F = \frac{C_r}{R} = \frac{105}{0,25} \approx \SI{420}{N}
\]

Force nécessaire pour gravir la pente (composante du poids le long de la pente) :
\[
F_p = m \cdot g \cdot \sin\alpha = 250 \times 10 \times 0,148 \approx \SI{370}{N}
\]

Comparaison : $F \approx \SI{420}{N} > F_p \approx \SI{370}{N}$. La traction est suffisante, avec une marge d'environ 13 \% pour vaincre les frottements et assurer les démarrages en côte. Le projet de M. Njoya est réalisable.
\end{exoresolu}

\begin{attention}[Pièges fréquents]
\begin{itemize}
\item Oublier de convertir km/h en m/s (diviser par 3,6) avant d'utiliser $v = R\,\omega$.
\item Inverser le rapport : ici $k < 1$ car on \emph{réduit} la vitesse (et on augmente le couple). Vérifie toujours : la sortie tourne-t-elle plus vite ou moins vite que l'entrée ?
\item Additionner les rapports au lieu de les multiplier : $k_{\text{global}} = k_1 \times k_2$.
\item Oublier le rendement $\eta$ dans le calcul du couple : le couple réel est toujours inférieur au couple idéal.
\item Confondre le rayon et le diamètre de la roue dans $F = \dfrac{C_r}{R}$ : ici $R = \SI{0,25}{m}$, pas \SI{0,50}{m}.
\item Dans $k = \dfrac{Z_e}{Z_s}$, inverser le nombre de dents de l'entrée et de la sortie : c'est la roue \emph{menée} qui est au dénominateur.
\end{itemize}
\end{attention}

\courssub{Bilan de l'activité}

Cette activité t'a permis de mobiliser l'ensemble des savoirs de la leçon dans une situation authentique de la vie camerounaise. Tu as constaté que :
\begin{itemize}
\item la vitesse angulaire relie le monde du moteur (tr/min) au monde de la route (km/h) par la relation $v = R\,\omega$ ;
\item le rapport de transmission $k$ est la grandeur centrale : il fixe à la fois la vitesse de sortie \emph{et} le couple disponible, car réduire la vitesse multiplie le couple ;
\item un rapport important se décompose en plusieurs étages, et le choix entre engrenages, chaîne et courroie résulte d'un compromis entre coût, rendement, robustesse et maintenance locale ;
\item aucune transmission motorisée n'est acceptable sans protections (carter, capot, écran) : la sécurité fait partie intégrante de la conception.
\end{itemize}
Comme M. Njoya, tu sais désormais transformer une fiche technique de moteur en une solution mécanique complète, chiffrée et sécurisée : c'est exactement le métier d'un technicien en mécanique.

\newpage

\coursec{Méthodes et exercices résolus}

\coursec{Généralités sur la transmission du mouvement de rotation}

\courssub{Rotation et grandeurs caractéristiques du mouvement}

\begin{methode}[Reconnaître et caractériser une rotation]
Pour étudier le mouvement d'un solide (roue de moto-taxi, hélice de ventilateur, tambour de machine à laver) :
\begin{enumerate}
\item \textbf{Observer} le solide : tous ses points décrivent-ils des cercles autour d'une même droite fixe ? Si oui, c'est un \cle{mouvement de rotation} autour d'un \cle{axe de rotation}.
\item \textbf{Repérer l'axe de rotation} : c'est la droite dont tous les points restent immobiles (l'axe de la roue, l'axe du moteur).
\item \textbf{Identifier le sens de rotation} : sens horaire (sens des aiguilles d'une montre) ou sens antihoraire, en précisant toujours le côté d'où l'on regarde.
\item \textbf{Mesurer la vitesse de rotation} $N$ : compter le nombre de tours effectués pendant une durée connue, puis calculer
\[ N = \frac{\text{nombre de tours}}{\text{durée}} \quad \text{en tr/s ou tr/min.} \]
\item \textbf{Convertir si besoin} : $1\ \text{tr/s} = 60\ \text{tr/min}$ ; pour passer des tr/min aux tr/s, diviser par 60.
\end{enumerate}
\textbf{Réflexe :} toujours préciser l'axe, le sens et la vitesse : ce sont les trois éléments qui caractérisent complètement une rotation.
\end{methode}

\begin{popfigure}
\begin{tikzpicture}[scale=1.2, >=Stealth]
% Wheel
\draw[popblue, thick] (0,0) circle (2cm);
% Axis
\draw[popdark, thick] (-2.5,0) -- (2.5,0);
\draw[popdark, fill=popdark] (0,0) circle (2pt);
% Radius to point M
\draw[poporange, thick] (0,0) -- (1.5,1.5);
\fill[poporange] (1.5,1.5) circle (3pt) node[above right] {$M$};
% Trajectory circle (dashed)
\draw[popmuted, dashed] (0,0) circle (2.12cm);
% Rotation arrow
\draw[popgreen, thick, ->] (1.5,0) arc (0:90:1.5cm);
\node[popgreen] at (1.8,1.8) {Sens de rotation};
% Labels
\node[popdark] at (-2.5,-0.5) {Axe de rotation $\Delta$};
\node[popdark] at (2.5,-0.5) {Centre $O$};
\end{tikzpicture}
\legende{Mouvement de rotation : tout point $M$ du solide décrit un cercle de centre $O$ sur l'axe $\Delta$.}
\end{popfigure}

\begin{exoresolu}[Le ventilateur de la maison]
Un ventilateur de salon effectue 900 tours en 30 secondes.
\begin{enumerate}
\item Justifier que les pales du ventilateur sont animées d'un mouvement de rotation.
\item Déterminer la vitesse de rotation $N$ des pales en tr/s, puis en tr/min.
\end{enumerate}
\tcblower
\textbf{1. Nature du mouvement.}

Observons une pale du ventilateur : chacun de ses points (sauf ceux situés sur l'axe du moteur) décrit un cercle dont le centre appartient à l'axe du moteur. Cet axe reste fixe pendant le fonctionnement.

\textbf{Conclusion :} les pales sont animées d'un \cle{mouvement de rotation} autour de l'axe du moteur, qui est l'\cle{axe de rotation}.

\textbf{2. Calcul de la vitesse de rotation.}

La vitesse de rotation est le quotient du nombre de tours par la durée :
\[ N = \frac{\text{nombre de tours}}{\text{durée}} = \frac{900\ \text{tours}}{30\ \text{s}} = 30\ \text{tr/s}. \]

Conversion en tours par minute : en une minute il y a 60 secondes, donc
\[ N = 30 \times 60 = 1800\ \text{tr/min}. \]

\textbf{Conclusion :} les pales du ventilateur tournent à la vitesse $N = 30\ \text{tr/s}$, c'est-à-dire $N = \SI{1800}{\text{tr/min}}$.
\end{exoresolu}

\courssub{Rapport de transmission : définition, calcul et signe}

\begin{methode}[Exploiter le rapport de transmission]
Dans une transmission entre une roue \cle{menante} (celle qui reçoit le mouvement du moteur) et une roue \cle{menée} (celle qui est entraînée) :
\begin{enumerate}
\item \textbf{Identifier} la roue menante (indice 1) et la roue menée (indice 2).
\item \textbf{Écrire le rapport de transmission} :
\[ r = \frac{N_2}{N_1} = \frac{Z_1}{Z_2} \quad \text{(engrenages : } Z = \text{nombre de dents)} \]
ou, pour les poulies :
\[ r = \frac{N_2}{N_1} = \frac{D_1}{D_2} \quad \text{(} D = \text{diamètre de la poulie).} \]
\item \textbf{Isoler l'inconnue} : par exemple $N_2 = r \times N_1$ ou $Z_2 = \dfrac{Z_1}{r}$.
\item \textbf{Interpréter le résultat} :
\begin{itemize}
\item si $r < 1$ : la roue menée tourne \textbf{moins vite} : c'est un \cle{réducteur} de vitesse ;
\item si $r > 1$ : la roue menée tourne \textbf{plus vite} : c'est un \cle{multiplicateur} de vitesse ;
\item si $r = 1$ : les deux roues tournent à la même vitesse.
\end{itemize}
\item \textbf{Préciser le sens de rotation} : deux roues dentées en contact direct tournent en \textbf{sens contraires} ; deux poulies reliées par une courroie non croisée tournent dans le \textbf{même sens}.
\end{enumerate}
\textbf{Réflexe :} dans un engrenage, la \textbf{petite} roue tourne toujours \textbf{plus vite} que la grande.
\end{methode}

\begin{popfigure}
\begin{tikzpicture}[scale=1.2, >=Stealth]
% Gear 1 (Driving)
\draw[popblue, thick] (-2.5,0) circle (1.5cm);
\draw[popblue, fill=popblue!20] (-2.5,0) circle (0.3cm);
% Teeth suggestion
\foreach \a in {0,30,...,330} \draw[popblue, thick] (-2.5,0) -- ++(\a:1.5cm) -- ++(\a:0.3cm);
% Gear 2 (Driven)
\draw[poporange, thick] (1.0,0) circle (3cm);
\draw[poporange, fill=poporange!20] (1.0,0) circle (0.3cm);
\foreach \a in {15,45,...,345} \draw[poporange, thick] (1.0,0) -- ++(\a:3cm) -- ++(\a:0.3cm);
% Contact point
\draw[popdark, dashed] (-1.0,0) -- (1.0,0) node[midway, above] {Contact};
% Arrows
\draw[popblue, thick, ->] (-2.5,1.5) arc (90:180:1.5cm) node[left, midway] {$N_1$};
\draw[poporange, thick, ->] (1.0,3) arc (90:0:3cm) node[right, midway] {$N_2$};
% Labels
\node[popblue] at (-2.5,-2) {Menante ($Z_1$)};
\node[poporange] at (1.0,-3.5) {Menée ($Z_2 > Z_1$)};
\node[popdark] at (-2.5,2.2) {Sens horaire};
\node[popdark] at (1.0,3.7) {Sens antihoraire};
\end{tikzpicture}
\legende{Engrenage simple : la roue menante (petite) et la roue menée (grande) tournent en sens contraires. $r = Z_1/Z_2 < 1$ (réducteur).}
\end{popfigure}

\begin{exoresolu}[Le rapport de transmission d'un engrenage]
Le pignon d'un moteur de moulin à maïs possède $Z_1 = 20$ dents et tourne à $N_1 = \SI{1500}{\text{tr/min}}$. Il engrène avec une roue menée de $Z_2 = 60$ dents.
\begin{enumerate}
\item Calculer le rapport de transmission $r$ de cet engrenage.
\item Calculer la vitesse de rotation $N_2$ de la roue menée.
\item S'agit-il d'un réducteur ou d'un multiplicateur ? Justifier.
\item Comparer les sens de rotation des deux roues.
\end{enumerate}
\tcblower
\textbf{1. Rapport de transmission.}

Pour un engrenage, le rapport de transmission est le quotient du nombre de dents de la roue menante par celui de la roue menée :
\[ r = \frac{Z_1}{Z_2} = \frac{20}{60} = \frac{1}{3} \approx 0,33. \]

\textbf{2. Vitesse de la roue menée.}

On sait que $r = \dfrac{N_2}{N_1}$, donc $N_2 = r \times N_1$ :
\[ N_2 = \frac{1}{3} \times 1500 = 500\ \text{tr/min}. \]

\textbf{3. Réducteur ou multiplicateur ?}

On a $r = \dfrac{1}{3} < 1$ : la roue menée tourne moins vite que la roue menante ($500 < 1500$ tr/min).

\textbf{Conclusion :} il s'agit d'un \cle{réducteur} de vitesse. C'est logique : la roue menée (60 dents) est plus grande que le pignon menant (20 dents), or la grande roue tourne toujours plus lentement.

\textbf{4. Sens de rotation.}

Les deux roues dentées sont en \textbf{contact direct} : elles tournent donc en \textbf{sens contraires}. Si le pignon du moteur tourne dans le sens horaire, la roue menée tourne dans le sens antihoraire.
\end{exoresolu}

\courssub{Les transmissions par contact direct : les engrenages}

\begin{methode}[Analyser un train d'engrenages]
Pour étudier une transmission par engrenages (boîte de vitesses, perceuse, réducteur de moulin) :
\begin{enumerate}
\item \textbf{Vérifier la condition d'engrenement} : deux roues dentées ne peuvent engrener que si elles ont le \textbf{même module} (même taille de dents).
\item \textbf{Compter les dents} de chaque roue : $Z_1$ (menante), $Z_2$ (menée).
\item \textbf{Compter le nombre de contacts} (d'engrenements) entre l'entrée et la sortie :
\begin{itemize}
\item nombre \textbf{pair} de contacts $\Rightarrow$ entrée et sortie tournent dans le \textbf{même sens} ;
\item nombre \textbf{impair} de contacts $\Rightarrow$ entrée et sortie tournent en \textbf{sens contraires}.
\end{itemize}
\item \textbf{Calculer le rapport global} en multipliant les rapports intermédiaires :
\[ r = \frac{N_{\text{sortie}}}{N_{\text{entrée}}} = \frac{\text{produit des } Z \text{ des roues menantes}}{\text{produit des } Z \text{ des roues menées}}. \]
\item \textbf{Conclure} sur la vitesse de sortie et sur le rôle du mécanisme (réducteur ou multiplicateur).
\end{enumerate}
\textbf{Réflexe :} à chaque contact entre deux roues, le sens de rotation s'inverse.
\end{methode}

\begin{popfigure}
\begin{tikzpicture}[scale=1.1, >=Stealth]
% Gear A
\draw[popblue, thick] (0,0) circle (1cm);
\draw[popblue, fill=popblue!20] (0,0) circle (0.2cm);
\foreach \a in {0,45,...,315} \draw[popblue, thick] (0,0) -- ++(\a:1cm) -- ++(\a:0.2cm);
\node at (0,-1.5) {A ($Z_A=15$)};
\node at (0,1.8) {$N_A$};
\draw[popblue, thick, ->] (0,1) arc (90:180:1cm);
% Gear B
\draw[poporange, thick] (2.5,0) circle (2cm);
\draw[poporange, fill=poporange!20] (2.5,0) circle (0.2cm);
\foreach \a in {22.5,67.5,...,337.5} \draw[poporange, thick] (2.5,0) -- ++(\a:2cm) -- ++(\a:0.2cm);
\node at (2.5,-2.5) {B ($Z_B=45$)};
% Axis B-C (solidary)
\draw[popdark, dashed, thick] (2.5,0) -- (5.5,0);
% Gear C (on same axis as B)
\draw[poppurple, thick] (5.5,0) circle (1cm);
\draw[poppurple, fill=poppurple!20] (5.5,0) circle (0.2cm);
\foreach \a in {0,45,...,315} \draw[poppurple, thick] (5.5,0) -- ++(\a:1cm) -- ++(\a:0.2cm);
\node at (5.5,-1.5) {C ($Z_C=15$)};
\node at (5.5,1.8) {$N_C = N_B$};
\draw[poppurple, thick, ->] (5.5,1) arc (90:180:1cm);
% Gear D
\draw[popgreen, thick] (8.5,0) circle (2.5cm);
\draw[popgreen, fill=popgreen!20] (8.5,0) circle (0.2cm);
\foreach \a in {18,54,...,342} \draw[popgreen, thick] (8.5,0) -- ++(\a:2.5cm) -- ++(\a:0.2cm);
\node at (8.5,-3) {D ($Z_D=60$)};
\node at (8.5,3.2) {$N_D$};
\draw[popgreen, thick, ->] (8.5,2.5) arc (90:0:2.5cm);
% Contact labels
\node at (1.25, 2.3) {Contact 1};
\node at (7.0, 2.3) {Contact 2};
% Solidary label
\node[popdark, rotate=90] at (4.0, 0.5) {Solidaires (même axe)};
\end{tikzpicture}
\legende{Train d'engrenages de la perceuse : deux contacts (pair) $\Rightarrow$ même sens entrée/sortie. Roues B et C solidaires : $N_B = N_C$.}
\end{popfigure}

\begin{exoresolu}[Train d'engrenages d'une perceuse]
Le moteur d'une perceuse entraîne un pignon A de $Z_A = 15$ dents à $N_A = \SI{3000}{\text{tr/min}}$. Le pignon A engrène avec une roue B de $Z_B = 45$ dents, solidaire (fixée sur le même axe) d'un pignon C de $Z_C = 15$ dents. Le pignon C engrène avec la roue de sortie D de $Z_D = 60$ dents, qui porte le mandrin.
\begin{enumerate}
\item Calculer la vitesse $N_B$ de la roue B.
\item En déduire la vitesse $N_D$ de la roue de sortie D.
\item Calculer le rapport global $r = \dfrac{N_D}{N_A}$ de deux façons différentes.
\item Le mandrin tourne-t-il dans le même sens que le moteur ?
\end{enumerate}
\tcblower
\textbf{1. Vitesse de la roue B.}

A engrène avec B (contact direct) :
\[ \frac{N_B}{N_A} = \frac{Z_A}{Z_B} \quad \Rightarrow \quad N_B = N_A \times \frac{Z_A}{Z_B} = 3000 \times \frac{15}{45} = 3000 \times \frac{1}{3} = 1000\ \text{tr/min}. \]

\textbf{2. Vitesse de la roue D.}

B et C sont solidaires du même axe, donc $N_C = N_B = 1000$ tr/min. C engrène avec D :
\[ N_D = N_C \times \frac{Z_C}{Z_D} = 1000 \times \frac{15}{60} = 1000 \times \frac{1}{4} = 250\ \text{tr/min}. \]

\textbf{3. Rapport global, par deux méthodes.}

\textit{Méthode 1 : à partir des vitesses.}
\[ r = \frac{N_D}{N_A} = \frac{250}{3000} = \frac{1}{12} \approx 0,083. \]

\textit{Méthode 2 : produit des rapports intermédiaires.}
\[ r = \frac{Z_A}{Z_B} \times \frac{Z_C}{Z_D} = \frac{15}{45} \times \frac{15}{60} = \frac{1}{3} \times \frac{1}{4} = \frac{1}{12}. \]
Les deux méthodes donnent le même résultat : $r = \dfrac{1}{12}$. C'est un réducteur puissant ($r \ll 1$) : le mandrin tourne 12 fois moins vite que le moteur, ce qui donne plus de force pour percer.

\textbf{4. Sens de rotation.}

Il y a \textbf{deux contacts} (A--B puis C--D), c'est-à-dire un nombre \textbf{pair} : le sens s'inverse deux fois, donc le mandrin D tourne dans le \textbf{même sens} que le moteur A.
\end{exoresolu}

\courssub{Les transmissions par liaison flexible : courroies et chaînes}

\begin{methode}[Analyser une transmission poulie-courroie]
Pour une transmission par courroie (groupe électrogène, machine à coudre, moto-pompe) :
\begin{enumerate}
\item \textbf{Identifier} la poulie menante (côté moteur) et la poulie menée, puis mesurer ou relever leurs diamètres $D_1$ et $D_2$.
\item \textbf{Appliquer la relation} (la courroie a la même vitesse partout) :
\[ r = \frac{N_2}{N_1} = \frac{D_1}{D_2}. \]
\item \textbf{Calculer l'inconnue} : $N_2 = N_1 \times \dfrac{D_1}{D_2}$.
\item \textbf{Déterminer le sens de rotation} :
\begin{itemize}
\item courroie \textbf{non croisée} (ouverte) : les deux poulies tournent dans le \textbf{même sens} ;
\item courroie \textbf{croisée} : les deux poulies tournent en \textbf{sens contraires}.
\end{itemize}
\item \textbf{Citer les défauts possibles} : la courroie peut \cle{patiner} (glisser) si elle est détendue ou graissée : la vitesse réelle est alors inférieure à la vitesse calculée.
\end{enumerate}
\textbf{Réflexe :} pour \textbf{augmenter} la vitesse de sortie, mettre la \textbf{grande} poulie côté moteur ; pour la \textbf{réduire}, mettre la \textbf{petite} poulie côté moteur.
\end{methode}

\begin{popfigure}
\begin{tikzpicture}[scale=1.1, >=Stealth]
% Open belt (top)
\draw[popblue, thick] (0,0) circle (1cm); % Pulley 1
\draw[popblue, fill=popblue!20] (0,0) circle (0.2cm);
\draw[popblue, thick, ->] (0,1) arc (90:-90:1cm); % Rotation arrow left
\node[popblue] at (0,-1.5) {Menante};
\node[popblue] at (0,1.5) {Sens horaire};

\draw[popblue, thick] (5,0) circle (2cm); % Pulley 2
\draw[popblue, fill=popblue!20] (5,0) circle (0.2cm);
\draw[popblue, thick, ->] (5,2) arc (90:-90:2cm); % Rotation arrow right
\node[popblue] at (5,-2.5) {Menée};
\node[popblue] at (5,2.5) {Sens horaire};

% Belt lines (open)
\draw[popdark, thick] (1,0) -- (3,0);
\draw[popdark, thick] (1,0) -- (3,0); % duplicate for thickness visual? no, just two lines top/bottom
\draw[popdark, thick] (0,1) -- (5,2); % Top tangent approx
\draw[popdark, thick] (0,-1) -- (5,-2); % Bottom tangent approx
% Correct tangents for circles (0,0) r=1 and (5,0) r=2
% Top tangent: y = 1 to y = 2? No, external tangents.
% Let's draw simple lines for schematic
\draw[popdark, thick] (1,0) -- (3,0); % Top horizontal approx
\draw[popdark, thick] (1,0) -- (3,0); % Bottom horizontal approx
% Actually use coordinates
\draw[popdark, thick] (0,1) -- (5,2);
\draw[popdark, thick] (0,-1) -- (5,-2);
\node at (2.5, 2.5) {Courroie non croisée};

% Crossed belt (bottom)
\begin{scope}[yshift=-6cm]
\draw[poporange, thick] (0,0) circle (1cm);
\draw[poporange, fill=poporange!20] (0,0) circle (0.2cm);
\draw[poporange, thick, ->] (0,1) arc (90:-90:1cm);
\node[poporange] at (0,-1.5) {Menante};
\node[poporange] at (0,1.5) {Sens horaire};

\draw[poporange, thick] (5,0) circle (2cm);
\draw[poporange, fill=poporange!20] (5,0) circle (0.2cm);
\draw[poporange, thick, ->] (5,-2) arc (-90:90:2cm); % Opposite sense
\node[poporange] at (5,-2.5) {Menée};
\node[poporange] at (5,-2.5) {Sens antihoraire}; % Corrected position
\node[poporange] at (5,2.5) {Sens antihoraire};

% Crossed belt lines
\draw[popdark, thick] (0,1) -- (5,-2);
\draw[popdark, thick] (0,-1) -- (5,2);
\node at (2.5, -2.5) {Courroie croisée};
\end{scope}
\end{tikzpicture}
\legende{Haut : courroie non croisée (même sens). Bas : courroie croisée (sens contraires).}
\end{popfigure}

\begin{exoresolu}[Le groupe électrogène du quartier]
Sur un groupe électrogène, la poulie du moteur a un diamètre $D_1 = \SI{100}{\text{mm}}$ et tourne à $N_1 = \SI{3000}{\text{tr/min}}$. Elle entraîne par une courroie non croisée la poulie de l'alternateur de diamètre $D_2 = \SI{200}{\text{mm}}$.
\begin{enumerate}
\item Calculer le rapport de transmission $r$.
\item Calculer la vitesse de rotation $N_2$ de l'alternateur.
\item L'alternateur tourne-t-il dans le même sens que le moteur ? Justifier.
\item On souhaite que l'alternateur tourne à \SI{1500}{\text{tr/min}} sans changer le moteur : quel diamètre $D_2'$ faut-il choisir ?
\end{enumerate}
\tcblower
\textbf{1. Rapport de transmission.}
\[ r = \frac{D_1}{D_2} = \frac{100}{200} = \frac{1}{2} = 0,5. \]

\textbf{2. Vitesse de l'alternateur.}
\[ N_2 = r \times N_1 = 0,5 \times 3000 = 1500\ \text{tr/min}. \]

\textbf{3. Sens de rotation.}

La courroie est \textbf{non croisée} (ouverte) : les deux poulies tournent donc dans le \textbf{même sens}. L'alternateur tourne dans le même sens que le moteur.

\textbf{4. Nouveau diamètre de la poulie menée.}

On veut $N_2' = 1500$ tr/min avec $N_1 = 3000$ tr/min. Or :
\[ \frac{N_2'}{N_1} = \frac{D_1}{D_2'} \quad \Rightarrow \quad D_2' = D_1 \times \frac{N_1}{N_2'} = 100 \times \frac{3000}{1500} = 200\ \text{mm}. \]

\textbf{Conclusion :} il faut une poulie menée de diamètre $D_2' = \SI{200}{\text{mm}}$. C'est exactement le montage actuel : la question 2 confirme que l'alternateur tourne déjà à \SI{1500}{\text{tr/min}}.
\end{exoresolu}

\begin{methode}[Analyser une transmission par chaîne]
Pour une transmission par chaîne (bicyclette, moto-taxi) :
\begin{enumerate}
\item \textbf{Identifier} le plateau (roue menante, côté pédalier ou moteur) et le pignon (roue menée, côté roue arrière), puis compter leurs nombres de dents $Z_1$ et $Z_2$.
\item \textbf{Appliquer la relation des engrenages} (la chaîne impose le même nombre de dents passant par seconde) :
\[ r = \frac{N_2}{N_1} = \frac{Z_1}{Z_2}. \]
\item \textbf{Conclure sur le sens} : plateau et pignon tournent toujours dans le \textbf{même sens} (la chaîne ne croise pas).
\item \textbf{Interpréter physiquement} :
\begin{itemize}
\item grand plateau + petit pignon $\Rightarrow$ $r > 1$ : la roue arrière tourne vite (vitesse élevée, mais effort grand) ;
\item petit plateau + grand pignon $\Rightarrow$ $r < 1$ : la roue tourne lentement (montées, effort réduit).
\end{itemize}
\item \textbf{Comparer avec la courroie} : la chaîne \textbf{ne patine pas} (transmission précise), mais elle demande un graissage régulier.
\end{enumerate}
\textbf{Réflexe :} une chaîne se comporte comme un engrenage « à distance » : même formule, mais mêmes sens de rotation.
\end{methode}

\begin{popfigure}
\begin{tikzpicture}[scale=1.1, >=Stealth]
% Sprocket 1 (Plateau) - Large
\draw[poppurple, thick] (0,0) circle (2.5cm);
\draw[poppurple, fill=poppurple!20] (0,0) circle (0.3cm);
% Teeth
\foreach \a in {0,15,...,345} \draw[poppurple, thick] (0,0) -- ++(\a:2.5cm) -- ++(\a:0.3cm);
\node[poppurple] at (0,-3) {Plateau ($Z_1$)};
\node[poppurple] at (0,3.2) {Sens horaire};
\draw[poppurple, thick, ->] (0,2.5) arc (90:180:2.5cm);

% Sprocket 2 (Pignon) - Small
\draw[poppurple, thick] (6,0) circle (1cm);
\draw[poppurple, fill=poppurple!20] (6,0) circle (0.3cm);
\foreach \a in {0,30,...,330} \draw[poppurple, thick] (6,0) -- ++(\a:1cm) -- ++(\a:0.3cm);
\node[poppurple] at (6,-1.8) {Pignon ($Z_2$)};
\node[poppurple] at (6,1.8) {Sens horaire};
\draw[poppurple, thick, ->] (6,1) arc (90:180:1cm);

% Chain (simplified as lines + arcs)
\draw[popdark, thick] (2.5,0) -- (5,0); % Top straight
\draw[popdark, thick] (2.5,0) -- (5,0); % Bottom straight (overlap)
% Better: top and bottom tangents
\draw[popdark, thick] (0,2.5) -- (6,1); % Top tangent approx
\draw[popdark, thick] (0,-2.5) -- (6,-1); % Bottom tangent approx
% Chain links suggestion
\foreach \x in {0.5,1.5,...,5.5} \draw[popmuted, thick] (\x, 2.5 - 0.3*\x) -- ++(0.8, -0.24);
\foreach \x in {0.5,1.5,...,5.5} \draw[popmuted, thick] (\x, -2.5 + 0.3*\x) -- ++(0.8, 0.24);

\node at (3, 3.5) {Chaîne tendue : même sens de rotation, pas de patinage};
\end{tikzpicture}
\legende{Transmission par chaîne (moto-taxi, bicyclette) : plateau et pignon tournent dans le même sens.}
\end{popfigure}

\begin{exoresolu}[La bicyclette du coursier]
Sur une bicyclette, le plateau du pédalier comporte $Z_1 = 44$ dents et le pignon de la roue arrière comporte $Z_2 = 11$ dents. Le cycliste pédale à $N_1 = 60$ tr/min.
\begin{enumerate}
\item Calculer le rapport de transmission $r$.
\item Calculer la vitesse de rotation $N_2$ de la roue arrière.
\item Pour gravir une côte, le cycliste passe sur un pignon de $Z_2' = 22$ dents. Calculer la nouvelle vitesse $N_2'$ de la roue arrière et expliquer l'intérêt de ce changement.
\end{enumerate}
\tcblower
\textbf{1. Rapport de transmission.}
\[ r = \frac{Z_1}{Z_2} = \frac{44}{11} = 4. \]

\textbf{2. Vitesse de la roue arrière.}
\[ N_2 = r \times N_1 = 4 \times 60 = 240\ \text{tr/min}. \]
La roue arrière fait 4 tours pendant que le pédalier en fait 1 : c'est un \cle{multiplicateur} de vitesse ($r > 1$).

\textbf{3. Avec le grand pignon.}
\[ N_2' = N_1 \times \frac{Z_1}{Z_2'} = 60 \times \frac{44}{22} = 60 \times 2 = 120\ \text{tr/min}. \]

\textbf{Interprétation :} la roue arrière tourne deux fois moins vite ($120 < 240$ tr/min) : le cycliste avance moins vite, mais l'effort à fournir sur les pédales est deux fois plus faible. C'est exactement ce qu'il faut pour gravir une côte sans se fatiguer : on \textbf{sacrifie la vitesse pour gagner en force}.

\textbf{Conclusion :} $r = 4$ ; $N_2 = 240$ tr/min ; avec le pignon de 22 dents, $N_2' = 120$ tr/min, ce qui facilite la montée.
\end{exoresolu}

\courssub{Systèmes particuliers, rendement et sécurité}

\begin{methode}[Choisir une transmission et travailler en sécurité]
Pour choisir le bon système de transmission dans une situation concrète :
\begin{enumerate}
\item \textbf{Repérer la position des axes} :
\begin{itemize}
\item axes \textbf{proches et parallèles}, transmission précise et puissante $\Rightarrow$ \textbf{engrenages} ;
\item axes \textbf{éloignés}, fonctionnement silencieux, protection en cas de blocage $\Rightarrow$ \textbf{poulies et courroie} ;
\item axes \textbf{éloignés}, transmission \textbf{sans glissement} et grande puissance $\Rightarrow$ \textbf{chaîne} (moto, bicyclette).
\end{itemize}
\item \textbf{Tenir compte du sens de rotation} souhaité : engrenage direct = sens contraires ; courroie ouverte ou chaîne = même sens ; courroie croisée = sens contraires.
\item \textbf{Tenir compte de l'entretien} : engrenages et chaînes demandent un graissage ; la courroie s'use et doit être tendue puis remplacée.
\item \textbf{Appliquer les consignes de sécurité} :
\begin{itemize}
\item ne jamais approcher les mains, les cheveux ou les vêtements amples d'une courroie, d'une chaîne ou d'un engrenage en mouvement ;
\item arrêter le moteur et attendre l'arrêt complet avant toute intervention (tension de courroie, graissage) ;
\item remettre en place les \cle{carters de protection} après chaque réglage ;
\item ne jamais retirer une courroie en la faisant sauter d'une poulie en marche.
\end{itemize}
\end{enumerate}
\textbf{Réflexe :} au BEPC, justifie toujours ton choix par au moins \textbf{deux arguments} (distance des axes, précision, sens de rotation, entretien ou sécurité).
\end{methode}

\begin{exoresolu}[Le moulin à manioc du village]
Un menuisier-mécanicien veut installer un moulin à manioc. Le moteur thermique (axe A) se trouve à \SI{1,5}{\text{m}} de l'arbre du moulin (axe B). Le moteur tourne à $N_A = \SI{2400}{\text{tr/min}}$ et le moulin doit tourner à $N_B = \SI{800}{\text{tr/min}}$, dans le même sens que le moteur.
\begin{enumerate}
\item Quel système de transmission conseilles-tu entre A et B ? Donner deux arguments.
\item Calculer le rapport de transmission nécessaire.
\item La poulie du moteur a un diamètre $D_A = \SI{120}{\text{mm}}$. Calculer le diamètre $D_B$ de la poulie du moulin.
\item Citer deux consignes de sécurité à respecter lors de l'utilisation de ce moulin.
\end{enumerate}
\tcblower
\textbf{1. Choix du système.}

Je conseille une transmission par \textbf{poulies et courroie non croisée}. Arguments :
\begin{itemize}
\item les deux axes sont \textbf{éloignés} (\SI{1,5}{\text{m}}) : un engrenage direct est impossible, la courroie convient parfaitement ;
\item la courroie non croisée garantit que le moulin tourne dans le \textbf{même sens} que le moteur, et elle protège le moteur en \textbf{patinant} si le moulin se bloque (morceau de manioc coincé).
\end{itemize}

\textbf{2. Rapport de transmission.}
\[ r = \frac{N_B}{N_A} = \frac{800}{2400} = \frac{1}{3} \approx 0,33. \]
C'est un réducteur ($r < 1$), ce qui est cohérent : le moulin tourne moins vite que le moteur.

\textbf{3. Diamètre de la poulie du moulin.}
\[ r = \frac{D_A}{D_B} \quad \Rightarrow \quad D_B = \frac{D_A}{r} = \frac{120}{\frac{1}{3}} = 120 \times 3 = 360\ \text{mm}. \]
La poulie du moulin doit avoir un diamètre $D_B = \SI{360}{\text{mm}}$ (trois fois plus grande que celle du moteur, donc trois fois plus lente : vérification cohérente).

\textbf{4. Consignes de sécurité.}
\begin{itemize}
\item Ne jamais approcher les mains ou les vêtements de la courroie et des poulies en mouvement ; mettre en place un \textbf{carter de protection} autour de la transmission.
\item \textbf{Arrêter le moteur} et attendre l'arrêt complet des poulies avant toute intervention (tension de la courroie, déblocage du moulin, graissage).
\end{itemize}
\end{exoresolu}

\begin{savaistu}[La transmission au Cameroun : du terrain à l'industrie]
\begin{itemize}
\item \textbf{Moto-taxi (\"Benskin\") :} La transmission finale est toujours par \textbf{chaîne} (ou cardan pour les grosses cylindrées). Le kit chaîne (plateau, pignon, chaîne) s'use vite sur nos routes ; un graissage tous les 500 km et une tension correcte (2-3 cm de jeu) évitent la casse.
\item \textbf{Groupe électrogène ENEO / Quartier :} Le moteur thermique (essence/gasoil SONARA) entraîne l'alternateur par \textbf{courroie trapézoïdale}. Le patinage au démarrage protège le moteur du pic de couple. Vérifier l'usure des gorges et la tension (déflexion 10-15 mm sous 10 kg).
\item \textbf{Moulin à maïs/manioc :} Souvent par \textbf{courroie plate} ou \textbf{engrenages} si le moteur est accolé. Le réducteur (grand rapport) transforme la vitesse moteur (3000 tr/min) en couple puissant (500-800 tr/min) pour broyer.
\item \textbf{Pompage eau (SODECA / Forage) :} Transmission par \textbf{courroie} ou \textbf{accouplement direct} (alignement parfait requis). La sécurité impose un carter sur la courroie et une mise à la terre du châssis.
\end{itemize}
\end{savaistu}

\begin{attention}[Les 5 erreurs qui coûtent des points au BEPC]
\begin{enumerate}
\item \textbf{Inverser la formule du rapport de transmission.} Pour un engrenage, $r = \dfrac{N_2}{N_1} = \dfrac{Z_1}{Z_2}$ : les vitesses et les nombres de dents sont « croisés » (la roue qui a le plus de dents tourne le moins vite). Écrire $r = \dfrac{Z_2}{Z_1}$ est l'erreur la plus fréquente.
\item \textbf{Oublier de préciser le sens de rotation.} Un rapport de transmission ne suffit pas : il faut toujours conclure sur les sens (engrenage direct $\Rightarrow$ sens contraires ; courroie ouverte ou chaîne $\Rightarrow$ même sens ; courroie croisée $\Rightarrow$ sens contraires).
\item \textbf{Se tromper dans la conversion tr/min $\leftrightarrow$ tr/s.} Pour passer des tr/s aux tr/min, on \textbf{multiplie} par 60 ; pour passer des tr/min aux tr/s, on \textbf{divise} par 60. Vérifier la cohérence : une valeur en tr/min est toujours 60 fois plus grande qu'en tr/s.
\item \textbf{Oublier les roues solidaires dans un train d'engrenages.} Deux roues fixées sur le même axe ont la \textbf{même vitesse} de rotation ($N_B = N_C$) : ne pas appliquer la formule du rapport entre elles.
\item \textbf{Négliger les unités et la conclusion.} Une vitesse sans unité (tr/min) est sanctionnée ; de même, il faut toujours terminer par une phrase : « il s'agit d'un réducteur car $r < 1$ » ou « la roue menée tourne plus vite car $r > 1$ ». Enfin, toute question de sécurité exige au moins une consigne concrète (carter, arrêt du moteur avant intervention).
\end{enumerate}
\end{attention}

\newpage

\coursec{Exercices}

\courssub{Je vérifie mes connaissances}

\begin{exercice}[QCM : les bons réflexes]
Pour chaque question, entoure la bonne réponse.
\begin{enumerate}
\item En technologie, l'unité la plus utilisée pour exprimer la vitesse de rotation d'un arbre de machine est :
\begin{enumerate}
\item le m/s \quad b) le tr/min \quad c) le N.m \quad d) le watt
\end{enumerate}
\item Une poulie motrice de diamètre $D_1 = \SI{100}{mm}$ entraîne une poulie menée de diamètre $D_2 = \SI{200}{mm}$. Le rapport de transmission $k = \dfrac{N_2}{N_1}$ vaut :
\begin{enumerate}
\item 2 \quad b) 0,5 \quad c) 100 \quad d) 20
\end{enumerate}
\item Dans un engrenage à deux roues dentées extérieures, les deux roues tournent :
\begin{enumerate}
\item dans le même sens \quad b) en sens contraires \quad c) cela dépend du lubrifiant
\end{enumerate}
\item Une courroie croisée fait tourner les deux poulies :
\begin{enumerate}
\item dans le même sens \quad b) en sens contraires \quad c) elle ne transmet pas le mouvement
\end{enumerate}
\end{enumerate}
\end{exercice}

\begin{exercice}[Vrai ou faux, en justifiant]
Réponds par VRAI ou FAUX, puis justifie ta réponse en une ou deux phrases.
\begin{enumerate}
\item Si une roue menée possède plus de dents que la roue menante, la vitesse de sortie est plus grande que la vitesse d'entrée.
\item Une chaîne de moto peut glisser sur les pignons comme une courroie lisse glisse sur les poulies.
\item Un système roue et vis sans fin est irréversible : c'est la vis qui entraîne la roue, jamais l'inverse.
\item Le rendement d'une transmission réelle est toujours supérieur à 1.
\end{enumerate}
\end{exercice}

\begin{exercice}[Définitions à compléter]
Recopie et complète les phrases suivantes avec les mots : \emph{rapport de transmission}, \emph{vitesse angulaire}, \emph{rotation}, \emph{engrenage}, \emph{courroie}.
\begin{enumerate}
\item Un solide est animé d'un mouvement de \ldots{} lorsque tous ses points décrivent des cercles autour d'un axe fixe.
\item La \ldots{} d'un arbre, notée $N$ ou $\omega$, exprime le nombre de tours effectués par unité de temps.
\item Le \ldots{} est le quotient de la vitesse de sortie par la vitesse d'entrée : $k = \dfrac{N_s}{N_e}$.
\item Deux roues dentées en contact direct forment un \ldots{} ; deux poulies reliées par un organe flexible forment une transmission par \ldots{}.
\end{enumerate}
\end{exercice}

\begin{exercice}[Calculs directs]
\begin{enumerate}
\item Un moteur de groupe électrogène tourne à $N_1 = \SI{3000}{tr/min}$. Le rapport de transmission vers l'alternateur est $k = 0,5$. Calcule la vitesse de rotation $N_2$ de l'alternateur.
\item Un pignon de $Z_1 = 20$ dents entraîne une roue de $Z_2 = 60$ dents. Calcule le rapport de transmission $k = \dfrac{Z_1}{Z_2}$.
\item Convertis $\SI{1500}{tr/min}$ en tr/s.
\end{enumerate}
\end{exercice}

\courssub{Je m'entraîne}

\begin{exercice}[Pompe entraînée par poulies-courroie]
Un moteur tourne à $N_m = \SI{1400}{tr/min}$. Il est couplé à une pompe à eau par deux poulies et une courroie : $D_{\text{moteur}} = \SI{80}{mm}$ et $D_{\text{pompe}} = \SI{240}{mm}$.
\begin{enumerate}
\item Calcule le rapport de transmission $k$.
\item Calcule la vitesse de rotation $N_p$ de la pompe.
\item Les deux arbres tournent dans le même sens : la courroie est-elle ouverte ou croisée ? Justifie.
\end{enumerate}
\end{exercice}

\begin{exercice}[Réducteur à deux étages]
Un réducteur comporte deux étages d'engrenages : étage 1 ($Z_1 = 16$ dents menant $Z_2 = 48$ dents), étage 2 ($Z_3 = 20$ dents menant $Z_4 = 60$ dents). La vitesse d'entrée est $N_e = \SI{3000}{tr/min}$.
\begin{enumerate}
\item Calcule le rapport de transmission de chaque étage, puis le rapport global $k$.
\item Calcule la vitesse de sortie $N_s$.
\item Précise le sens de rotation de l'arbre de sortie par rapport à l'arbre d'entrée. Justifie.
\end{enumerate}
\end{exercice}

\begin{exercice}[Le vélo du coursier]
Sur un vélo, le plateau (pédalier) possède $Z_1 = 44$ dents et le pignon arrière $Z_2 = 11$ dents. Le cycliste pédale à $N_1 = \SI{60}{tr/min}$.
\begin{enumerate}
\item Calcule le rapport de transmission de la chaîne.
\item Calcule la vitesse de rotation de la roue arrière en tr/min.
\item S'agit-il d'un multiplicateur ou d'un réducteur de vitesse ? Justifie.
\end{enumerate}
\end{exercice}

\begin{exercice}[Choix technique : le ventilateur de plafond]
Pour équiper un ventilateur de plafond (vitesse lente, couple faible, fonctionnement silencieux, axe vertical), on hésite entre quatre solutions : engrenages droits, courroie trapézoïdale, courroie dentée, roue et vis sans fin.
\begin{enumerate}
\item Choisis la solution la plus adaptée.
\item Justifie ton choix par trois arguments techniques (bruit, glissement possible en sécurité, coût et entretien).
\end{enumerate}
\end{exercice}

\courssub{Je me prépare au Bac}

\begin{exercice}[Le treuil de chantier (type BEPC)]
Sur un chantier à Douala, un treuil est entraîné par un moteur électrique à travers un réducteur roue et vis sans fin de rapport $k = \dfrac{1}{40}$ et de rendement $\eta = 0,6$. Le moteur fournit un couple $C_m = \SI{5}{N.m}$ et tourne à $N_m = \SI{1500}{tr/min}$. La charge à lever exerce un effort $F = \SI{500}{N}$ sur un câble enroulé sur un tambour de rayon $r = \SI{0,1}{m}$.
\begin{enumerate}
\item Calcule la vitesse de rotation $N_t$ du tambour.
\item Calcule le couple disponible $C_t$ en sortie du réducteur, sachant que $C_t = \dfrac{C_m \times \eta}{k}$.
\item Calcule le couple résistant $C_r$ exercé par la charge sur le tambour ($C_r = F \times r$).
\item Le système peut-il lever la charge ? Justifie par une comparaison.
\item Cite deux raisons qui rendent le système roue et vis sans fin intéressant pour un treuil (sécurité, irréversibilité).
\end{enumerate}
\end{exercice}

\begin{exercice}[Analyse de document : le différentiel d'un camion]
Sur la photo d'un pont arrière de camion ouvert, on observe : une couronne dentée, un pignon d'attaque, deux pignons coniques appelés satellites et deux planétaires reliés aux roues.
\begin{enumerate}
\item Identifie, sur le document, les engrenages coniques, la couronne et le pignon d'attaque.
\item Quel est le rôle du couple pignon-couronne sur la direction de l'axe de rotation ?
\item Explique en trois phrases le rôle du différentiel lorsque le camion prend un virage (pense aux distances parcourues par la roue intérieure et la roue extérieure).
\item Pourquoi dit-on que les engrenages coniques permettent une transmission entre arbres concourants ?
\end{enumerate}
\end{exercice}

\begin{exercice}[Sécurité autour d'une transmission à chaîne (type BEPC)]
Dans un atelier de menuiserie à Yaoundé, une scie circulaire est entraînée par une transmission à chaîne laissée sans protection, à hauteur des mains de l'ouvrier.
\begin{enumerate}
\item Sur le schéma fourni, entoure en rouge trois zones de pincement dangereuses (points de contact chaîne-pignon, brin tendu, sortie de pignon).
\item Dessine le carter de protection normé : précise sa forme, son mode de fixation et le matériau adapté.
\item Cite deux consignes de sécurité à respecter avant toute intervention sur la transmission (arrêt, verrouillage, EPI).
\item La réglementation du travail impose la protection des organes en mouvement des machines : explique en deux phrases pourquoi cette obligation protège à la fois l'ouvrier et l'employeur.
\end{enumerate}
\end{exercice}



\newpage

\coursec{Corrigés des exercices}

\courssub{Je vérifie mes connaissances}

\begin{corrige}[Exercice 1 — QCM : les bons réflexes]
\begin{enumerate}
\item Réponse b) : le \cle{tr/min} (tours par minute) est l'unité usuelle en technologie. Le m/s est une unité de vitesse linéaire, le N.m une unité de couple, le watt une unité de puissance.
\item Réponse b) : pour une transmission par poulies-courroie sans glissement, $k = \dfrac{N_2}{N_1} = \dfrac{D_1}{D_2} = \dfrac{100}{200} = 0,5$. La grande poulie tourne deux fois moins vite que la petite.
\item Réponse b) : au point de contact, les dents de deux roues extérieures se déplacent en sens opposés : les roues tournent en \cle{sens contraires}.
\item Réponse b) : croiser la courroie \cle{inverse} le sens de rotation ; une courroie ouverte (non croisée) conserve le sens.
\end{enumerate}
\end{corrige}

\begin{corrige}[Exercice 2 — Vrai ou faux, en justifiant]
\begin{enumerate}
\item FAUX. Pour un engrenage, $k = \dfrac{N_s}{N_e} = \dfrac{Z_{\text{menante}}}{Z_{\text{menée}}}$. Si la roue menée a plus de dents, alors $k < 1$ : la vitesse de sortie est plus \emph{faible} (réducteur), mais le couple est plus grand.
\item FAUX. La chaîne engrène sur les dents des pignons : la transmission se fait par \cle{obstacle}, sans glissement possible. C'est un avantage de la chaîne sur la courroie lisse.
\item VRAI. Dans un système roue et vis sans fin, seule la vis peut entraîner la roue ; l'inverse est impossible (\cle{irréversibilité}), ce qui constitue une sécurité, par exemple pour un treuil.
\item FAUX. À cause des frottements, une partie de l'énergie est perdue en chaleur : le rendement $\eta = \dfrac{P_{\text{sortie}}}{P_{\text{entrée}}}$ est toujours \emph{inférieur} à 1 (soit moins de \SI{100}{\percent}).
\end{enumerate}
\end{corrige}

\begin{corrige}[Exercice 3 — Définitions à compléter]
\begin{enumerate}
\item Un solide est animé d'un mouvement de \cle{rotation} lorsque tous ses points décrivent des cercles autour d'un axe fixe.
\item La \cle{vitesse angulaire} d'un arbre, notée $N$ ou $\omega$, exprime le nombre de tours effectués par unité de temps.
\item Le \cle{rapport de transmission} est le quotient de la vitesse de sortie par la vitesse d'entrée : $k = \dfrac{N_s}{N_e}$.
\item Deux roues dentées en contact direct forment un \cle{engrenage} ; deux poulies reliées par un organe flexible forment une transmission par \cle{courroie}.
\end{enumerate}
\end{corrige}

\begin{corrige}[Exercice 4 — Calculs directs]
\begin{enumerate}
\item $N_2 = k \times N_1 = 0,5 \times 3000 = \SI{1500}{tr/min}$.
\item $k = \dfrac{Z_1}{Z_2} = \dfrac{20}{60} = \dfrac{1}{3} \approx 0,33$. Comme $k < 1$, c'est un réducteur.
\item $N = \dfrac{1500}{60} = \SI{25}{tr/s}$, car une minute contient 60 secondes.
\end{enumerate}
\end{corrige}

\courssub{Je m'entraîne}

\begin{corrige}[Exercice 5 — Pompe entraînée par poulies-courroie]
\begin{enumerate}
\item Pour une transmission par courroie sans glissement : $k = \dfrac{N_p}{N_m} = \dfrac{D_{\text{moteur}}}{D_{\text{pompe}}} = \dfrac{80}{240} = \dfrac{1}{3} \approx 0,33$.
\item $N_p = k \times N_m = \dfrac{1}{3} \times 1400 \approx \SI{467}{tr/min}$. La pompe tourne environ trois fois moins vite que le moteur : c'est un réducteur.
\item La courroie est \cle{ouverte} (non croisée) : une courroie ouverte conserve le sens de rotation, alors qu'une courroie croisée l'inverse.
\end{enumerate}
\end{corrige}

\begin{corrige}[Exercice 6 — Réducteur à deux étages]
\begin{enumerate}
\item Étage 1 : $k_1 = \dfrac{Z_1}{Z_2} = \dfrac{16}{48} = \dfrac{1}{3}$. Étage 2 : $k_2 = \dfrac{Z_3}{Z_4} = \dfrac{20}{60} = \dfrac{1}{3}$. Le rapport global est le produit des rapports : $k = k_1 \times k_2 = \dfrac{1}{3} \times \dfrac{1}{3} = \dfrac{1}{9} \approx 0,11$.
\item $N_s = k \times N_e = \dfrac{3000}{9} \approx \SI{333}{tr/min}$.
\item Chaque étage d'engrenages extérieurs inverse le sens de rotation. Deux inversions successives redonnent le sens initial : l'arbre de sortie tourne dans le \textbf{même sens} que l'arbre d'entrée.
\end{enumerate}
\end{corrige}

\begin{corrige}[Exercice 7 — Le vélo du coursier]
\begin{enumerate}
\item Pour une transmission par chaîne : $k = \dfrac{N_2}{N_1} = \dfrac{Z_1}{Z_2} = \dfrac{44}{11} = 4$.
\item $N_2 = k \times N_1 = 4 \times 60 = \SI{240}{tr/min}$. La roue arrière fait 4 tours pour 1 tour de pédalier.
\item Comme $k = 4 > 1$, la vitesse de sortie est plus grande que la vitesse d'entrée : c'est un \cle{multiplicateur} de vitesse. Le cycliste fournit en contrepartie un effort plus grand sur les pédales.
\end{enumerate}
\end{corrige}

\begin{corrige}[Exercice 8 — Choix technique : le ventilateur de plafond]
\begin{enumerate}
\item La solution la plus adaptée est la \cle{courroie trapézoïdale}.
\item Trois arguments techniques :
\begin{itemize}
\item \textbf{Bruit} : la courroie trapézoïdale est souple et silencieuse, contrairement aux engrenages droits qui sont bruyants ; c'est essentiel pour un appareil domestique.
\item \textbf{Sécurité} : en cas de blocage des pales (un enfant qui accroche un objet), la courroie peut \emph{glisser} sur les poulies, ce qui protège le moteur ; une courroie dentée ou un engrenage transmettrait tout l'effort.
\item \textbf{Coût et entretien} : la courroie est bon marché, ne demande pas de lubrification et se remplace facilement ; le système roue et vis sans fin, coûteux et de faible rendement, est inutile ici car le couple à transmettre est faible.
\end{itemize}
\end{enumerate}
\end{corrige}

\courssub{Je me prépare au Bac}

\begin{corrige}[Exercice 9 — Le treuil de chantier (type BEPC)]
\begin{enumerate}
\item $N_t = k \times N_m = \dfrac{1}{40} \times 1500 = \SI{37,5}{tr/min}$.
\item $C_t = \dfrac{C_m \times \eta}{k} = \dfrac{5 \times 0,6}{\frac{1}{40}} = 3 \times 40 = \SI{120}{N.m}$.
\item $C_r = F \times r = 500 \times 0,1 = \SI{50}{N.m}$.
\item On compare : $C_t = \SI{120}{N.m} > C_r = \SI{50}{N.m}$. Le couple moteur disponible est supérieur au couple résistant : le système \textbf{peut lever la charge}, avec une marge de sécurité confortable.
\item Deux avantages du système roue et vis sans fin pour un treuil :
\begin{itemize}
\item \textbf{Irréversibilité} : si le moteur s'arrête, la charge ne peut pas entraîner la vis en arrière ; la charge reste bloquée en l'air, ce qui évite les chutes.
\item \textbf{Grand rapport de réduction} en un seul étage (ici $1/40$) : le couple est fortement multiplié, ce qui permet de lever de lourdes charges avec un petit moteur.
\end{itemize}
\end{enumerate}
\textbf{Barème indicatif (10 points)} : question 1 : 1,5 pt ; question 2 : 2 pts ; question 3 : 1,5 pt ; question 4 : 2 pts (comparaison chiffrée exigée) ; question 5 : 3 pts (1,5 pt par avantage justifié).\par
\textbf{Points de vigilance} : écrire l'application numérique avant le résultat ; ne pas oublier les unités (tr/min, N.m) ; à la question 4, une simple réponse « oui » sans comparaison chiffrée ne rapporte aucun point ; le rendement $\eta$ s'applique au couple de sortie, pas à la vitesse.
\end{corrige}

\begin{corrige}[Exercice 10 — Analyse de document : le différentiel d'un camion]
\begin{enumerate}
\item Le pignon d'attaque est la petite roue conique au bout de l'arbre de transmission (arrivée du moteur) ; la couronne est la grande roue conique qu'il entraîne ; les satellites et les planétaires sont les petits pignons coniques logés dans le boîtier du différentiel.
\item Le couple pignon-couronne est un engrenage conique : il \textbf{change la direction de l'axe de rotation de \SI{90}{\degree}}, en faisant passer le mouvement de l'arbre longitudinal (le long du camion) aux arbres transversaux des roues arrière.
\item Dans un virage, la roue extérieure parcourt une distance plus grande que la roue intérieure : elle doit donc tourner plus vite. Les satellites du différentiel se mettent à tourner sur eux-mêmes et répartissent le mouvement : ils accélèrent la roue extérieure et ralentissent la roue intérieure d'autant. Ainsi, les deux roues motrices restent entraînées sans patiner ni user les pneus.
\item Les axes du pignon d'attaque et de la couronne se coupent (ils sont \emph{concourants}) : les dents coniques permettent de transmettre le mouvement entre deux arbres dont les axes se rencontrent, généralement à angle droit.
\end{enumerate}
\textbf{Barème indicatif (10 points)} : question 1 : 2 pts ; question 2 : 2 pts ; question 3 : 4 pts ; question 4 : 2 pts.\par
\textbf{Points de vigilance} : à la question 3, la justification doit mentionner explicitement la différence de distances parcourues par les deux roues ; employer le vocabulaire exact (couronne, satellites, planétaires, arbres concourants).
\end{corrige}

\begin{corrige}[Exercice 11 — Sécurité autour d'une transmission à chaîne (type BEPC)]
\begin{enumerate}
\item Les trois zones à entourer sont : le point d'entrée de la chaîne sur le pignon moteur, le point d'entrée de la chaîne sur le pignon mené (sortie de pignon), et le brin tendu de la chaîne entre les deux pignons, où un doigt ou un vêtement peut être happé.
\item Le carter est un capot fermé qui épouse la forme de la transmission (profil allongé couvrant les deux pignons et la chaîne), en tôle d'acier ou en matériau rigide résistant aux chocs. Il est fixé au bâti de la machine par des vis ou des boulons, de façon à ne pouvoir être retiré qu'avec un outil.
\item Deux consignes avant toute intervention : arrêter la machine, la débrancher et \textbf{verrouiller} l'alimentation électrique (consignation) pour empêcher tout redémarrage intempestif ; porter les équipements de protection individuelle (gants ajustés, lunettes, vêtements serrés, pas de bijoux).
\item Protéger les organes en mouvement évite à l'ouvrier les accidents graves (happement, coupure, amputation) et préserve sa santé et sa capacité de travail. Elle protège aussi l'employeur, qui évite les arrêts de travail, les poursuites et les indemnisations liées aux accidents du travail.
\end{enumerate}
\textbf{Barème indicatif (10 points)} : question 1 : 3 pts (1 pt par zone correctement identifiée) ; question 2 : 2,5 pts ; question 3 : 2,5 pts ; question 4 : 2 pts.\par
\textbf{Points de vigilance} : le carter doit être décrit comme \emph{fixé de façon non amovible sans outil} ; à la question 3, citer l'arrêt seul ne suffit pas, le verrouillage (consignation) est attendu ; à la question 4, expliciter les deux points de vue (ouvrier \emph{et} employeur).
\end{corrige}

\newpage

\coursec{Fiche bilan}

\begin{popfigure}
\begin{tikzpicture}[mindmap, grow cyclic, every node/.style={concept}, concept color=popblue!60, text=white,
  root concept/.append style={font=\bfseries, minimum size=2.6cm, text width=2.4cm},
  level 1/.append style={level distance=3.6cm, sibling angle=72, font=\small},
  level 1 concept/.append style={minimum size=1.9cm, text width=1.9cm},
  scale=0.9, transform shape]
\node[root concept]{Transmission du mouvement de rotation}
  child[concept color=poporange!80]{ node{Rotation : $\omega$, $N$ (tr/min)} }
  child[concept color=popgreen!70]{ node{Rapport $r=\dfrac{N_s}{N_e}$} }
  child[concept color=poppurple!70]{ node{Engrenages : contact direct} }
  child[concept color=poppink!80]{ node{Courroies : liaison flexible} }
  child[concept color=popgold!85]{ node{Cha\^ines : pignons} };
\end{tikzpicture}
\legende{Carte mentale de la le\c{c}on : les cinq id\'ees ma\^itresses de la transmission du mouvement de rotation.}
\end{popfigure}

\begin{aretenir}[L'essentiel de la le\c{c}on]
\textbf{D\'efinitions cl\'es}
\begin{itemize}
  \item \cle{Mouvement de rotation} : mouvement d'un solide dont tous les points d\'ecrivent des cercles autour d'une droite fixe appel\'ee \cle{axe de rotation}.
  \item \cle{Vitesse angulaire} $\omega$ : angle balay\'e par unit\'e de temps, en rad/s : $\omega = \dfrac{\theta}{t}$.
  \item \cle{Fr\'equence de rotation} $N$ : nombre de tours par unit\'e de temps, en tr/min ou tr/s. Lien : $\omega = \dfrac{2\pi N}{60}$ si $N$ est en tr/min.
  \item \cle{Rapport de transmission} : $r = \dfrac{N_s}{N_e}$ (vitesse de sortie sur vitesse d'entr\'ee).
  \item \cle{R\'educteur} si $r<1$ ; \cle{multiplicateur} si $r>1$.
\end{itemize}

\textbf{Formules \`a conna\^itre par c\oe ur}
\begin{center}
\begin{tabularx}{\linewidth}{|l|X|X|}
\hline
\rowcolor{popblueL} \textbf{Grandeur} & \textbf{Formule} & \textbf{Unit\'es} \\
\hline
Vitesse angulaire & $\omega=\dfrac{2\pi N}{60}$ & $\omega$ en rad/s, $N$ en tr/min \\
\hline
Rapport (engrenages) & $r=\dfrac{N_s}{N_e}=\pm\dfrac{Z_e}{Z_s}$ & $Z$ : nombre de dents \\
\hline
Rapport (poulies-courroie) & $r=\dfrac{N_s}{N_e}=\dfrac{D_e}{D_s}$ & $D$ : diam\`etres des poulies \\
\hline
Rendement & $\eta=\dfrac{P_u}{P_f}$ & sans unit\'e, $\eta<1$ \\
\hline
\end{tabularx}
\end{center}

\textbf{M\'ethodes express}
\begin{itemize}
  \item Engrenage ext\'erieur : les roues tournent en \cle{sens contraires} (signe $-$) ; avec une roue interm\'ediaire, le sens est conserv\'e.
  \item Pour $n$ roues en prise ext\'erieure : le sens de sortie est le m\^eme que le sens d'entr\'ee si $n$ est impair, invers\'e si $n$ est pair.
  \item Courroie ouverte : m\^eme sens ; courroie crois\'ee : sens contraires.
  \item Petite roue menante $\Rightarrow$ r\'eduction de vitesse mais \cle{augmentation du couple} (ex. : petit plateau du v\'elo pour monter une c\^ote).
  \item Calcul type : identifier roue menante/men\'ee $\to$ \'ecrire $r$ $\to$ $N_s = r \times N_e$ $\to$ conclure sur le sens.
\end{itemize}
\end{aretenir}

\begin{attention}[Pi\`eges fr\'equents au BEPC]
\begin{itemize}
  \item Ne pas confondre $N_e$ (entr\'ee, roue \cle{menante}) et $N_s$ (sortie, roue \cle{men\'ee}) : le rapport s'\'ecrit toujours $r=\dfrac{N_s}{N_e}$, jamais l'inverse.
  \item Dans $r=\dfrac{Z_e}{Z_s}$, c'est bien le nombre de dents de la roue \textbf{menante} ($Z_e$) qui est au num\'erateur : c'est l'inverse de la formule avec les vitesses.
  \item Pour convertir en rad/s, $N$ doit \^etre en tr/min dans la formule $\omega=\dfrac{2\pi N}{60}$ ; si $N$ est en tr/s, alors $\omega = 2\pi N$.
  \item Le rendement $\eta$ est toujours inf\'erieur \`a 1 (ou \`a 100 \%) : si ton calcul donne $\eta > 1$, tu as invers\'e $P_u$ et $P_f$.
  \item V\'erifier la coh\'erence du r\'esultat : une petite roue menante entra\^ine une grande roue men\'ee \textbf{plus lente} ($r<1$).
\end{itemize}
\end{attention}

\begin{aretenir}[Checklist avant le BEPC]
\begin{itemize}
  \item[$\square$] Je sais d\'efinir un mouvement de rotation et son axe.
  \item[$\square$] Je sais convertir tr/min en rad/s avec $\omega=\dfrac{2\pi N}{60}$.
  \item[$\square$] Je sais d\'efinir et calculer le rapport de transmission $r=\dfrac{N_s}{N_e}$.
  \item[$\square$] Je sais dire si un syst\`eme est r\'educteur ($r<1$) ou multiplicateur ($r>1$).
  \item[$\square$] Je sais calculer $r$ avec les nombres de dents : $r=\dfrac{Z_e}{Z_s}$.
  \item[$\square$] Je sais calculer $r$ avec les diam\`etres des poulies : $r=\dfrac{D_e}{D_s}$.
  \item[$\square$] Je sais pr\'evoir le sens de rotation d'un engrenage (contact ext\'erieur : sens contraires).
  \item[$\square$] Je distingue courroie ouverte (m\^eme sens) et courroie crois\'ee (sens contraires).
  \item[$\square$] Je sais citer des exemples camerounais : bo\^ite de vitesses de moto-taxi, cha\^ine de v\'elo, courroie de groupe \'electrog\`ene.
  \item[$\square$] Je conna\^is les r\`egles de s\'ecurit\'e : carter de protection, machine arr\^et\'ee avant toute intervention, pas de v\^etements flottants.
  \item[$\square$] Je v\'erifie toujours la coh\'erence de mon r\'esultat (petite roue menante $\Rightarrow$ sortie plus lente).
\end{itemize}
\end{aretenir}

\findecours

\end{document}
